\documentclass[11pt]{article}
\usepackage[letterpaper,portrait,margin=1in]{geometry}
\usepackage[T1]{fontenc}
\usepackage{lmodern}
\usepackage{amsmath,amssymb,amsthm,mathtools}
\usepackage{microtype}
\usepackage{xcolor}
\usepackage{tikz}
\usetikzlibrary{arrows.meta,calc,fit,positioning}
\usepackage{aliascnt}
\usepackage[colorlinks=true,linkcolor=teal,citecolor=teal,urlcolor=teal]{hyperref}
\usepackage[nameinlink,capitalize]{cleveref}

\allowdisplaybreaks

\newtheorem{theorem}{Theorem}[section]
\newaliascnt{lemma}{theorem}
\newtheorem{lemma}[lemma]{Lemma}
\aliascntresetthe{lemma}
\newaliascnt{proposition}{theorem}
\newtheorem{proposition}[proposition]{Proposition}
\aliascntresetthe{proposition}
\newaliascnt{corollary}{theorem}
\newtheorem{corollary}[corollary]{Corollary}
\aliascntresetthe{corollary}
\newtheorem{fact}{Fact}[section]
\renewcommand{\thefact}{\thesection.\Roman{fact}}
\newaliascnt{remark}{theorem}
\theoremstyle{remark}
\newtheorem{remark}[remark]{Remark}
\aliascntresetthe{remark}
\theoremstyle{plain}

\crefname{theorem}{Theorem}{Theorems}
\Crefname{theorem}{Theorem}{Theorems}
\crefname{lemma}{Lemma}{Lemmas}
\Crefname{lemma}{Lemma}{Lemmas}
\crefname{proposition}{Proposition}{Propositions}
\Crefname{proposition}{Proposition}{Propositions}
\crefname{corollary}{Corollary}{Corollaries}
\Crefname{corollary}{Corollary}{Corollaries}
\crefname{fact}{Fact}{Facts}
\Crefname{fact}{Fact}{Facts}
\crefname{remark}{Remark}{Remarks}
\Crefname{remark}{Remark}{Remarks}
\crefname{equation}{equation}{equations}
\Crefname{equation}{Equation}{Equations}

\newcommand{\C}{\mathbb C}
\newcommand{\R}{\mathbb R}
\newcommand{\cA}{\mathcal A}
\newcommand{\cC}{\mathcal C}
\newcommand{\cE}{\mathcal E}
\newcommand{\cF}{\mathcal F}
\newcommand{\cH}{\mathcal H}
\newcommand{\cM}{\mathcal M}
\newcommand{\cS}{\mathcal S}
\newcommand{\cT}{\mathcal T}
\newcommand{\U}{\mathrm U}
\newcommand{\dd}{\mathrm d}
\newcommand{\diag}{\operatorname{diag}}
\newcommand{\Gr}{\operatorname{Gr}}
\newcommand{\id}{\operatorname{id}}
\newcommand{\Tr}{\operatorname{Tr}}
\newcommand{\eps}{\varepsilon}
\newcommand{\ket}[1]{|#1\rangle}
\newcommand{\proj}[1]{|#1\rangle\!\langle #1|}
\newcommand{\ot}{\otimes}
\newcommand{\Forb}{F_{\mathrm{orb}}}
\newcommand{\Fcl}{F_{\mathrm{cl}}}
\newcommand{\Funiv}{F_{\mathrm{univ}}}
\newcommand{\doi}[2]{\href{https://doi.org/#1}{#2}}
\newcommand{\arxiv}[1]{\href{https://arxiv.org/abs/#1}{\texttt{arXiv:#1}}}

\title{Asymptotically Optimal Mixed-State Cloning}
\author{Jiani Fei\(^{*}\)\qquad
  Jinzhao Wang\(^{*,\dagger}\)}
\date{}
\hypersetup{
  pdftitle={Asymptotically Optimal Mixed-State Cloning},
  pdfauthor={Jiani Fei and Jinzhao Wang},
  pdfsubject={Proportional-output cloning on fixed-rank quantum-state strata},
  pdfkeywords={quantum cloning, fixed-rank states, Schur--Weyl duality,
    local asymptotic normality, purify-and-clone}
}
\begin{document}
\maketitle
\begingroup
\renewcommand{\thefootnote}{\fnsymbol{footnote}}
\footnotetext[1]{\textit{Leinweber Institute for Theoretical Physics,
  Stanford University, Stanford, CA 94305, USA}}
\footnotetext[2]{\textit{Zyphra Technology, San Francisco, CA 94105, USA.}
  \enspace Email:
  \href{mailto:wang.jinzhao226@gmail.com}{\texttt{wang.jinzhao226@gmail.com}}}
\endgroup

\begin{abstract}
Unknown quantum states can only be cloned approximately.  Although the
optimal pure-state cloner has long been known, the mixed-state problem remains
challenging.  We determine the asymptotically optimal global
\(n\)-to-\(m_n\) cloning root fidelity on every fixed-rank stratum with a
simple positive spectrum, in fixed dimension and with
\(m_n/n\to\gamma>1\).  If the rank \(r\) and the positive spectrum
\(p_1>\cdots>p_r>0\) are known, the minimax fidelity on the corresponding
unitary orbit is
\[
  \Forb^{(r,d)}(\gamma,p)
  =\gamma^{-r(d-r)/2}\prod_{1\le i<j\le r}
  \frac{\sqrt{\gamma}+\sqrt{q_{ij}(\gamma-1+q_{ij})}}
       {\gamma+q_{ij}},
  \qquad q_{ij}=\frac{p_j}{p_i}.
\]
If the spectrum is unknown but the rank is known, a single rank-aware cloner
guarantees
\[
  \Funiv^{(r,d)}(\gamma,p)
  =\left(\frac{2\sqrt\gamma}{1+\gamma}\right)^{(r-1)/2}
   \Forb^{(r,d)}(\gamma,p),
\]
and this gives the minimax value on every regular compact subset of the
rank-\(r\) spectral chamber.  The case \(r=d\) recovers the full-rank
interior case without a separate argument.  We also solve the positively
degenerate flat-spectrum orbit: for
\(\rho=P/r\), \(P\in\Gr(r,d)\), the optimal root fidelity is
\(\gamma^{-r(d-r)/2}\).  Unlike fidelity-threshold copy-complexity results,
our theorems make no near-unit-fidelity assumption: the fixed gain
\(\gamma>1\) is arbitrary, and the displayed formulas give the explicit,
generally nonunit, asymptotic fidelity.  We further upper-bound the actual
mixed-system output of purify--clone--trace and prove that our cloner has
strictly larger asymptotic fidelity for every simple positive spectrum of
rank \(r>1\) and every flat projector orbit with \(r>1\); for \(r=1\), the
protocols coincide.  The proofs combine Schur--Weyl label processing, an
exact rank-compression identity, Cartan-sector channels, and fixed-rank local
asymptotic normality in its interior and boundary forms.
\end{abstract}

\tableofcontents

\section{Introduction}\label{sec:introduction-results}

The no-cloning theorem rules out perfect copying of an arbitrary unknown
quantum state \cite{WoottersZurek1982,Dieks1982}.  Approximate cloning asks
instead how well a quantum channel can transform \(n\) identical inputs into
an output resembling \(m\) ideal copies.  For pure states, early universal
qubit machines were followed by exact solutions for the standard global and
single-clone criteria
\cite{BuzekHillery1996,GisinMassar1997,
BrussDiVincenzoEkertFuchsMacchiavelloSmolin1998}: Werner solved the global
universal problem in arbitrary dimension \cite{Werner1998}, while
Keyl--Werner established optimality when individual clones are tested
\cite{KeylWerner1999}.  State-dependent global cloning was developed in
\cite{BrussDiVincenzoEkertFuchsMacchiavelloSmolin1998,
CheflesBarnett1999}; see also the reviews
\cite{ScaraniIblisdirGisinAcin2005,FanWangJingYueShiZhangMu2014}.

The choice of figure of merit is essential.  At the level of individual
clones, the many-output limit is closely tied to state estimation
\cite{MassarPopescu1995,BrussEkertMacchiavello1998,BaeAcin2006,
ChiribellaDAriano2006}.  When all outputs are examined jointly, the relation
is subtler: global asymptotic cloning has been analyzed for broad pure-state
families, and an optimal measure-and-prepare output need not consist of
independent copies of the estimated state
\cite{ChiribellaYang2014Optimal}.  The dependence of the optimum on a joint,
marginal, or process-level criterion is already visible for coherent-state
and phase-covariant families
\cite{CerfKrugerNavezWernerWolf2005,KimChitambar2022}.  Here we use root
fidelity with the entire product target \(\rho^{\ot m}\).  This global
criterion is stronger than reproducing one-site marginals and is therefore
distinct from broadcasting, which permits correlations among the outputs
\cite{BarnumCavesFuchsJozsaSchumacher1996}.

Mixed states do not
evade the exact no-cloning obstruction: the general
no-broadcasting theorem rules out a single channel that clones every member
of a noncommuting family of density matrices, with an entropic proof given by
Kalev--Hen
\cite{BarnumCavesFuchsJozsaSchumacher1996,KalevHen2008}.  Approximate
mixed-state cloning poses an additional structural difficulty: both the
eigenvalues and the eigenbasis may be unknown, and the two kinds of
information behave differently under repeated sampling.  In contrast to the
pure-state case, earlier mixed-state results either did
not identify an optimal cloner with a matching fixed-gain value or applied
only to special state families or different criteria.  Work directly on
mixed-state cloning derived global-fidelity, relative-error, and
state-dependent bounds for prescribed finite ensembles
\cite{Rastegin2003Upper,Rastegin2002Relative,
Rastegin2003StateDependent}; constructed cloners for symmetric-subspace inputs
and identical mixed qubits under one-site shrinking-factor criteria
\cite{Fan2003,FanLiuShi2007}; studied a restricted mixed-qubit construction
under an averaged Hilbert--Schmidt criterion \cite{Kazakov2007}; and established
general entropic limitations on simultaneous approximate cloning and
broadcasting of mixed states \cite{LemmWilde2017}.  Closely related work solved
optimal mixed-qubit purification and a local-fidelity cloning/estimation
problem, rather than the joint product-fidelity task considered here
\cite{CiracEkertMacchiavello1999}.  Multi-input mixed-state work has also
focused on broadcasting and superbroadcasting, which optimize marginals
rather than the joint product target
\cite{DArianoMacchiavelloPerinotti2005,
BuscemiDArianoMacchiavelloPerinotti2005Optimal,
BuscemiDArianoMacchiavelloPerinotti2006,DangFan2007}.  More recently, the
random-purification channel of Tang--Wright--Zhandry was combined with
pure-state cloning to give a general all-output guarantee, rather than an
exact fixed-gain optimum, for rank-constrained mixed states
\cite{TangWrightZhandry2025,LiTheilHarrowChuang2026}.  Concurrent work has
since asked how many inputs are needed to guarantee squared fidelity at least
\(1-\epsilon\), uniformly over states of rank at most \(r\), and shown that this
purify--clone--trace construction has optimal copy complexity up to constants
\cite{FanizzaGrinkoScharnhorstSpilecki2026,JeonSohnOh2026}.  The resulting
\(1/\epsilon\) behavior describes the near-unit-fidelity boundary.  Our
primary question imposes no such accuracy requirement: we fix an arbitrary
\(\gamma=m_n/n>1\), determine the explicit asymptotic fidelity curve, which
may remain bounded away from one, and prove that it lies strictly above the
actual PCT fidelity in the cases covered by our comparison theorem.  See
\cref{sec:pct-comparison}.

The mixed-state cloning problem admits two models.  In the known-spectrum problem, the
eigenvalues are supplied and the state ranges over one unitary orbit; cloning
must reproduce only the unknown eigenbasis.  In the unknown-spectrum problem,
a single rank-aware but otherwise spectrum-independent sequence of channels
must handle both the positive eigenvalues and the eigenbasis, uniformly over
compact subsets of a fixed-rank spectral chamber.
The second problem is not merely the first preceded by spectrum estimation:
the measured spectrum retains fluctuations of order \(n^{-1/2}\), which must
themselves be converted from the \(n\)-copy to the \(m_n\)-copy scale.

Here we resolve the mixed-state cloning problem in the proportional-output
asymptotic regime on every fixed-rank stratum with distinct positive
eigenvalues.  For every fixed dimension and \(m_n/n\to\gamma>1\), we
determine the optimal global cloning fidelity on each such orbit and, when
the positive spectrum is unknown but its rank is known, uniformly over
regular compact spectral sets.  We additionally solve the flat positive
spectrum on the Grassmannian projector orbit.
Thus our primary result is a fixed-gain fidelity-curve theorem, not merely a
threshold statement near fidelity one: \(\gamma\) need not be close to one,
and the limiting fidelity may be nonunit.  For every rank \(r>1\) state with
simple positive spectrum, and for every flat projector state with \(r>1\),
our achievable fidelity is also strictly larger than an upper bound on the
actual post-trace PCT output.

Two complementary decompositions drive the solution.  Schur--Weyl duality
separates the spectral information, carried by the Young label, from the
eigenbasis information, carried by the conditional irreducible state
\cite{KeylWerner2001Spectrum}.  Our cloners process the label classically and
apply a multiplicity-one Cartan-sector channel to the irreducible state
\cite{ParthasarathyRangaRaoVaradarajan1967,FultonHarris1991}.
Two-way local asymptotic normality (LAN) gives the corresponding asymptotic
separation into \(r-1\) classical spectral coordinates,
\(\binom r2\) displaced thermal modes within the support, and
\(r(d-r)\) coherent support--kernel modes
\cite{GutaJencova2007LAN,GutaKahn2006LANQubits,
KahnGuta2009LANFiniteDimensional,LahiryNussbaum2024}.
It transfers any putative improvement to a Gaussian dilation or amplification
problem.  Quantum amplifier noise limits
\cite{Caves1982LinearAmplifiers} and the least-noise mixed-Gaussian cloner
\cite{GutaMatsumoto2006MixedGaussianCloning} identify the quantum benchmark;
a unified flat-prior fidelity argument gives the classical and quantum
converses used here.

\subsection{Results}\label{sec:main-results}

Fix positive integers \(m_n\) such that
\[
  \gamma_n:=\frac{m_n}{n}\longrightarrow\gamma>1,
  \qquad r_d:=\binom d2.
\]
In particular, \(m_n>n\) for all sufficiently large \(n\).  For the
finitely many indices with \(m_n\le n\), every cloner below is defined by
the marginal channel, with the empty partial trace interpreted as the
identity.  Thus all nontrivial formulas concern \(m_n>n\).

For \(1\le r\le d\), write
\[
  \mathsf A_r:=\left\{x\in\R^r:\sum_{i=1}^r x_i=1\right\},
  \qquad
  \Delta_r^\circ
  :=\left\{p\in\mathsf A_r:p_1>\cdots>p_r>0\right\},
\]
and, for \(p\in\Delta_r^\circ\), set
\[
  \bar p=(p_1,\ldots,p_r,0,\ldots,0)\in\R^d,
  \qquad
  \rho_{\bar p}=\diag(\bar p),
  \qquad
  \rho_{\bar p,g}=g\rho_{\bar p}g^*.
\]
The active orbital roots are
\begin{equation}\label{eq:active-root-set-main}
  \mathcal I_{r,d}
  :=\{(i,j):1\le i\le r,\ i<j\le d\},
  \qquad
  q_{ij}:=\frac{\bar p_j}{\bar p_i}
  =\begin{cases}
      p_j/p_i,&j\le r,\\
      0,&j>r.
    \end{cases}
\end{equation}
Thus positive--positive roots give thermal modes, support--kernel roots give
vacuum modes, and kernel--kernel roots are absent because they belong to the
stabilizer.
When \(r=d\), we suppress the bar and write \(\rho_p\) and
\(\rho_{p,g}\).
For positive trace-class operators \(A,C\), not necessarily normalized, we
write
\[
  F(A,C)=\|\sqrt A\sqrt C\|_1.
\]
For states this is the unsquared Uhlmann fidelity.  For nonnegative functions
or measures, the same notation denotes their Hellinger affinity; in
particular, \(F(P,Q)=\sum_x\sqrt{P(x)Q(x)}\) on a countable space.  The target
in every statement below is the entire product state.  Define
\[
  f_\gamma(q)
  =\frac{\sqrt\gamma+\sqrt{q(\gamma-1+q)}}{\gamma+q},
  \qquad
  c_\gamma=\frac{2\sqrt\gamma}{1+\gamma},
\]
and, for \(p\in\Delta_r^\circ\),
\begin{align}
  \Forb^{(r,d)}(\gamma,p)
  &=\prod_{(i,j)\in\mathcal I_{r,d}}f_\gamma(q_{ij})
   =\gamma^{-r(d-r)/2}
    \prod_{1\le i<j\le r}f_\gamma\!\left(\frac{p_j}{p_i}\right),
  \label{eq:rank-orb-value}\\
  \Funiv^{(r,d)}(\gamma,p)
  &=c_\gamma^{(r-1)/2}\Forb^{(r,d)}(\gamma,p).
  \label{eq:rank-univ-value}
\end{align}
When \(r=d\), we occasionally use the abbreviations
\begin{equation}\label{eq:main-values}
  \Forb(\gamma,p):=\Forb^{(d,d)}(\gamma,p),\qquad
  \Fcl(\gamma,d):=c_\gamma^{(d-1)/2},\qquad
  \Funiv(\gamma,p):=\Funiv^{(d,d)}(\gamma,p).
\end{equation}
The first product in \eqref{eq:rank-orb-value} is the unified form of the
orbital answer.  Each positive--positive root contributes its thermal factor
\(f_\gamma(p_j/p_i)\), while every support--kernel root has \(q_{ij}=0\)
and contributes \(f_\gamma(0)=\gamma^{-1/2}\).  The additional factor
\(c_\gamma^{(r-1)/2}\) is the cost of dilating the \(r-1\) unknown positive
spectral fluctuations.

\paragraph{Known spectrum}
Fixing \(p\) isolates the genuinely noncommutative part of the problem.  The
target distribution of Young labels is then known and can be sampled exactly
at size \(m_n\); only the conditional representation sectors, which encode
the eigenbasis, need to be cloned.

\begin{theorem}[Known-spectrum fixed-rank optimum]\label{thm:known-optimum}
For every \(1\le r\le d\) and \(p\in\Delta_r^\circ\),
\begin{equation}\label{eq:known-optimum}
  \lim_{n\to\infty}
  \sup_{\cM_n}\inf_{g\in\U(d)}
  F\!\left(\cM_n(\rho_{\bar p,g}^{\ot n}),
             \rho_{\bar p,g}^{\ot m_n}\right)
  =\Forb^{(r,d)}(\gamma,p),
\end{equation}
where the supremum is over all quantum channels from \(n\) to \(m_n\)
\(d\)-level systems.  The explicit Schur--Cartan channel described in
\cref{sec:rank-aware-achievability} asymptotically
attains the displayed value.
\end{theorem}

Thus each positive--positive pair contributes one displaced-thermal factor,
while each support--kernel pair contributes
\(f_\gamma(0)=\gamma^{-1/2}\).  When \(r=d\) there are no
support--kernel factors, while \(r=1\) recovers Werner's pure-state value.

\paragraph{Unknown spectrum}
If \(p\) is not supplied, the target Young-label distribution can no longer
be sampled directly.  The measured input label determines the spectrum only
at the \(n^{-1/2}\) scale, and converting \(n\) copies into
\(m_n\sim\gamma n\) copies requires dilating this classical fluctuation.
This produces the extra, spectrum-independent factor
\(c_\gamma^{(r-1)/2}\).

\begin{theorem}[Unknown-spectrum achievability and compact-set optimum]
\label{thm:unknown-optimum}
For every fixed \(1\le r\le d\), there exists a sequence of rank-aware,
spectrum-independent, \(\U(d)\)-covariant channels \(\cC_n^{(r)}\) with the
following properties.
\begin{enumerate}
\item[(a)] For every compact \(K\Subset\Delta_r^\circ\),
\begin{equation}\label{eq:unknown-uniform-achievability}
  \liminf_{n\to\infty}
  \inf_{\substack{p\in K\\g\in\U(d)}}
  \left[
  F\!\left(
    \cC_n^{(r)}(\rho_{\bar p,g}^{\ot n}),
    \rho_{\bar p,g}^{\ot m_n}
  \right)-\Funiv^{(r,d)}(\gamma,p)
  \right]\ge0.
\end{equation}

\item[(b)] Let \(K\Subset\Delta_r^\circ\) be nonempty and
  satisfy
  \(K=\overline{\operatorname{int}_{\mathsf A_r}K}^{\,\mathsf A_r}\).
Then
\begin{equation}\label{eq:compact-global-value}
  \lim_{n\to\infty}
  \sup_{\cM_n}
  \inf_{\substack{p\in K\\g\in\U(d)}}
  F\!\left(\cM_n(\rho_{\bar p,g}^{\ot n}),
             \rho_{\bar p,g}^{\ot m_n}\right)
  =\inf_{p\in K}\Funiv^{(r,d)}(\gamma,p),
\end{equation}
where the supremum is over all quantum channels from \(n\) to \(m_n\)
\(d\)-level systems.  The same sequence \(\cC_n^{(r)}\), described in
\cref{sec:rank-aware-achievability}, asymptotically
attains the displayed value.
\end{enumerate}
Here spectrum independent means independent of \(p\) within the specified
rank-\(r\) stratum; the channel is allowed to know \(r\).  When \(r=d\), we
also write \(\cC_n^{(d)}=\cC_n^{\mathrm{univ}}\).
\end{theorem}

\begin{theorem}[Proportional cloning on the Grassmannian projector orbit]
\label{thm:grassmann}
For fixed \(1\le r<d\),
\begin{equation}\label{eq:grassmann-value}
  \lim_{n\to\infty}\sup_{\cM_n}\inf_{P\in\Gr(r,d)}
  F\!\left(\cM_n((P/r)^{\ot n}),(P/r)^{\ot m_n}\right)
  =\gamma^{-r(d-r)/2}.
\end{equation}
Equivalently, the squared-fidelity value is \(\gamma^{-r(d-r)}\).
For \(r=d\), the orbit is a singleton and the value is trivially one.
\end{theorem}

The common architecture of the known-spectrum and rank-aware cloners is
summarized in \cref{fig:schur-cartan-cloner}.  After the Schur transform, the
Young label \(\mu\) carries the spectral information, while the conditional
irrep state carries the unknown eigenbasis.  A classical kernel chooses an
output label \(\nu\), and the common covariant sector channel maps the conditional
state into the \(\nu\)-sector before the \(m_n\)-system output is reconstructed.
For rank \(r\), the relevant labels have at most \(r\) rows and are padded by
zeros.  For known \(p\), the kernel samples the exact target law in effective
dimension \(r\); the unknown-spectrum kernel randomly dilates only the
active \(r\) rows.  Exact rank compression supplies the support--kernel
factor.  The details are developed in
\cref{sec:label-kernels,sec:fixed-rank-proof}.

\begin{figure}[!tbp]
\centering
\resizebox{\linewidth}{!}{%
\begin{tikzpicture}[
  flow/.style={-{Latex[length=2mm]},thick},
  dependency/.style={-{Latex[length=1.8mm]},thin,gray},
  guide/.style={thin,gray},
  box/.style={draw,rounded corners=2pt,align=center,inner sep=4pt,
    minimum height=9mm,font=\small},
  labelbox/.style={box,fill=blue!6},
  sectorbox/.style={box,fill=orange!9},
  note/.style={draw,rounded corners=2pt,align=left,inner sep=4pt,
    minimum height=12mm,fill=gray!6,font=\scriptsize}
]
  \node[box,fill=gray!8,text width=17mm] (input) at (-0.5,0)
    {input\\[-1pt]\(\rho_{\bar p,g}^{\ot n}\)};

  \node[labelbox,text width=24mm] (labelin) at (3.5,1.15)
    {Young label\\\(\mu\sim P_{n,p}\)};
  \node[sectorbox,text width=24mm] (sectorin) at (3.5,-1.15)
    {irrep state\\\(\rho_{p,g,\mu}\in\cH_\mu\)};

  \node[labelbox,text width=27mm] (labelout) at (7.1,1.15)
    {label kernel \(K_n\)\\\(\nu\in\mathsf Y_{m_n}\)};
  \node[sectorbox,text width=27mm] (sectorout) at (7.1,-1.15)
    {sector channel\\\(\cE_{\mu,\nu}(\rho_{p,g,\mu})\)};

  \node[box,fill=green!8,text width=28mm] (blockout) at (10.8,0)
    {assemble the \(\nu\)-block\\[-1pt]
     \(\cE_{\mu,\nu}(\rho_{p,g,\mu})
       \ot I_{\cS_\nu}/\dim\cS_\nu\)};
  \node[box,fill=gray!8,text width=22mm,font=\scriptsize]
    (output) at (15.55,0)
    {output state\\[-1pt]
     \(\mathcal Q_n[K_n](\rho_{\bar p,g}^{\ot n})\)};

  \coordinate (schursplit) at (2.0,0);
  \draw[flow] (schursplit) |- (labelin.west);
  \draw[flow] (schursplit) |- (sectorin.west);
  \draw[flow] (input.east) --
    node[below=2pt,pos=.45,font=\scriptsize,align=center]
      {Schur\\transform} (schursplit);
  \draw[flow] (labelin.east) -- (labelout.west);
  \draw[flow] (sectorin.east) -- (sectorout.west);
  \draw[dependency] (labelout.south) --
    node[right,font=\scriptsize,text=gray] {selects \(\nu\)} (sectorout.north);
  \coordinate (assemblemerge) at (8.9,0);
  \draw[thick] (labelout.east) -| (assemblemerge);
  \draw[thick] (sectorout.east) -| (assemblemerge);
  \draw[flow] (assemblemerge) -- (blockout.west);
  \draw[flow] (blockout.east) --
    node[below=2pt,font=\scriptsize,align=center]
      {inverse Schur\\transform} (output.west);

  \node[note,text width=47mm] (known) at (4.45,3.0)
    {\textbf{Known \(p\):}
     \(K_n^p(\nu\mid\mu)=P_{m_n,p}^{(r)}(\nu)\).};
  \node[note,text width=58mm] (unknown) at (10.25,3.0)
    {\textbf{Unknown \(p\):}
     \(K_n^{\mathrm{rnd}}\) randomly rounds \(\gamma_n\mu\).};
  \coordinate (kernelchoice) at ([yshift=4mm]labelout.north);
  \draw[guide] (known.south) -- (kernelchoice);
  \draw[guide] (unknown.south) -- (kernelchoice);
  \draw[dependency] (kernelchoice) -- (labelout.north);

\end{tikzpicture}%
}
\caption{\textbf{Cloner Construction}.
The Schur transform separates the classical Young label from the conditional
irrep state.  A Markov kernel selects the output label, while the generic
sector channel \(\cE_{\mu,\nu}\) maps the conditional state into that sector;
a maximally mixed Specht factor is then restored. Here
\(\mathcal Q_n[K_n]\) denotes the lifted channel associated with the selected
kernel. For known spectrum the
target label is sampled independently of \(\mu\), whereas for unknown
spectrum a randomized dilation is used; these choices produce our cloners
for \(m_n>n\), namely
   the known-spectrum and rank-aware universal channels, respectively.
On the label pairs
carrying asymptotically all probability, the sector fidelity tends to
\(\Forb^{(r,d)}(\gamma,p)\).  Labels have at most \(r\) active rows and
are padded by \(d-r\) zeros; the randomized dilation acts only on those
active rows.  }
\label{fig:schur-cartan-cloner}
\end{figure}


The minimax meaning of \eqref{eq:compact-global-value} is as follows.
The set \(K\) may be known when the channel is designed, but the actual
spectrum \(p\in K\) and eigenbasis \(g\) are not.  The order of optimization
is important: the supremum first allows a channel sequence tailored to
\(K\), and the inner infimum then tests that single sequence against every
\(p\in K\) and \(g\in\U(d)\).  Since the unknown-spectrum benchmark
\(\Funiv^{(r,d)}(\gamma,p)\) varies with \(p\), the resulting worst-case scalar is
governed by its least favorable value, hence the infimum over \(p\) on the
right-hand side.  No further infimum over \(g\) appears there because
unitary covariance makes the benchmark constant along the orbit of each
fixed spectrum.  Thus, even with knowledge of \(K\), no \(K\)-adapted
sequence of channels can asymptotically beat
\(\inf_{p\in K}\Funiv^{(r,d)}(\gamma,p)\).  The single rank-\(r\)
construction \(\cC_n^{(r)}\), which does not depend on \(K\), guarantees
\(\Funiv^{(r,d)}(\gamma,p)\) pointwise, uniformly for \(p\in K\), and hence
asymptotically attains this worst-case value.

The condition
\(K=\overline{\operatorname{int}_{\mathsf A_r}K}^{\,\mathsf A_r}\)
is a clean sufficient regularity assumption for full-dimensional spectral
uncertainty relative to the fixed-rank face.  Its interior contains small
neighborhoods in all \(r-1\)
spectral directions, and the closure condition allows the infimum over
interior points in the converse to be replaced, by continuity of
\(\Funiv^{(r,d)}\),
with the infimum over all of \(K\).  We do not claim that this condition is
necessary.  When \(r\ge2\), however, some thickness assumption cannot be
removed from a geometry-independent statement of this form.  For a
singleton \(K=\{p\}\) with \(r\ge2\),
the channel may be tailored to \(p\), so the problem has the known-spectrum
value \(\Forb^{(r,d)}(\gamma,p)\), not
\(\Funiv^{(r,d)}(\gamma,p)\).  More generally,
lower-dimensional spectral classes encode partial spectral knowledge and
may have different local Gaussian penalties.  Thus
\eqref{eq:compact-global-value} is the broadest geometry-independent
full-dimensional converse proved here; no statement with the same
right-hand side can extend to all compact \(K\) when \(r\ge2\).  For
\(r=1\), the spectral chamber is itself a singleton and
\(\Forb^{(1,d)}=\Funiv^{(1,d)}\), so this distinction disappears.

For example, a natural admissible class keeps the smallest positive
eigenvalue and every adjacent positive spectral gap uniformly bounded away
from zero.  For \(r\ge2\) and
\(0<\kappa<2/[r(r+1)]\),\footnote{At
\(\kappa=2/[r(r+1)]\), the identity
\(
  1=rp_r+\sum_{i=1}^{r-1}i(p_i-p_{i+1})
  \ge\frac{\kappa r(r+1)}2
\)
forces a unique spectrum, so the class is a singleton and its optimum is
\(\Forb^{(r,d)}\), not \(\Funiv^{(r,d)}\).  The compact-set theorem is not asserted uniformly
as the spectrum approaches an eigenvalue collision or the boundary of the
probability simplex.} set
\[
  \Delta_{r,\kappa}
  =\left\{p\in\mathsf A_r:p_r\ge\kappa,\
    p_i-p_{i+1}\ge\kappa\ (i<r)\right\}.
\]
This is a nonempty compact subset of \(\Delta_r^\circ\) and satisfies
\(\Delta_{r,\kappa}
=\overline{\operatorname{int}_{\mathsf A_r}\Delta_{r,\kappa}}
^{\,\mathsf A_r}\),
so \eqref{eq:compact-global-value} applies with
\(K=\Delta_{r,\kappa}\).  Indeed, choose
\(\kappa'\in(\kappa,2/[r(r+1)])\) and set
\(p_i=1/r+((r+1)/2-i)\kappa'\).  Then
\(p_i-p_{i+1}=\kappa'>\kappa\) and
\(p_r>2/[r(r+1)]>\kappa\), so this is a relative-interior point.
The set is closed in the probability simplex, hence compact; being convex
with nonempty relative interior gives the stated closure property.


\subsection{Remarks}
\label{sec:remarks}

\paragraph{Dependence on spectral mixedness.}
The product formula makes precise one sense in which a more mixed positive
spectrum is easier to clone.  For fixed \(\gamma>1\),
\[
  f_\gamma'(q)
  =\frac{(\gamma-1)^2}
  {2\sqrt{q(\gamma-1+q)}
   \bigl(\sqrt{\gamma(\gamma-1+q)}+\sqrt q\bigr)^2}>0,
  \qquad 0<q<1.
\]
Consequently, if \(p,p'\in\Delta_r^\circ\) satisfy
\[
  \frac{p'_j}{p'_i}\ge \frac{p_j}{p_i}
  \qquad\text{for every }1\le i<j\le r,
\]
then
\[
  \Forb^{(r,d)}(\gamma,p')\ge\Forb^{(r,d)}(\gamma,p),
  \qquad
  \Funiv^{(r,d)}(\gamma,p')\ge\Funiv^{(r,d)}(\gamma,p),
\]
with strict inequality if any ratio is strictly larger.  The
support--kernel factor is independent of \(p\).  Ratio-wise flattening is
stronger than majorization when \(r\ge3\); majorization alone need not
increase the cloning fidelity.

\paragraph{Limitations and remaining degeneracies.}
The fixed-rank theorems allow zero eigenvalues but require the positive
eigenvalues to be distinct.  The Grassmannian theorem separately treats the
maximally degenerate positive spectrum \(p_i=1/r\).  We do not resolve
arbitrary intermediate multiplicity patterns, rank-changing local
neighborhoods, or uniformity as \(p_r\downarrow0\), a positive spectral gap
closes, or \(\gamma\downarrow1\).  We also do not determine the exact
finite-copy optimum.  The rank-aware universal channel is allowed to know
\(r\); no theorem here gives a single channel uniformly over changing rank.
The lack of uniformity as \(\gamma\downarrow1\) concerns triangular joint
small-error limits only; it does not restrict any theorem at fixed
\(\gamma>1\).

\paragraph{Conjecture for general positive degeneracies.}
Suppose a rank-\(r\) spectrum has distinct positive values
\(\lambda_1>\cdots>\lambda_s>0\) with multiplicities
\(r_1,\ldots,r_s\), where \(\sum_a r_a=r\).  The exact compression identity
proved below isolates the zero eigenspace, leaving only the unresolved
positive-degeneracy problem in effective dimension \(r\).  We conjecture the
known-spectrum orbital value
\[
  \gamma^{-r(d-r)/2}
  \prod_{a<b}
  f_\gamma\!\left(\frac{\lambda_b}{\lambda_a}\right)^{r_ar_b}.
\]
For an unknown positive spectrum in a rank-preserving neighborhood that may
split the blocks, the conjectural additional classical factor is
\(c_\gamma^{(\sum_a r_a^2-1)/2}\).  Rotations within a degenerate block lie
in the stabilizer and create no orbital mode, whereas block-splitting
directions become classical.  The simple-positive-spectrum formulas above and the
flat known-spectrum projector theorem are the two proved endpoints; the
intermediate positive-degeneracy cases remain open.

\subsection{Comparison to purify-and-clone}
\label{sec:pct-comparison}

Concurrent work of Fanizza--Grinko--Scharnhorst--Spilecki and
Jeon--Sohn--Oh asks how many inputs are needed to guarantee a prescribed
squared fidelity \(1-\epsilon\), uniformly over all states of rank at most
\(r\) \cite{FanizzaGrinkoScharnhorstSpilecki2026,JeonSohnOh2026}.  To put
their fidelity-threshold statements and our fixed-gain statements in a
common convention, let \(N\) be the number of inputs, let \(M\) be the number
of \emph{additional} outputs, and write
\[
  m=N+M,\qquad
  \delta=\frac{M}{N},\qquad
  \gamma=1+\delta,\qquad
  t=\frac{N}{M}=\delta^{-1},\qquad
  \mathcal F:=F^2.
\]
Thus \(\mathcal F\) is the squared Uhlmann fidelity used in those papers.
Uniformly over states of rank at most \(r\), they establish the optimal
scaling
\[
  N=\Theta\!\left(\frac{Mrd}{\epsilon}\right).
\]
Their lower bounds already use flat projector families of rank comparable to
\(r\), while their matching upper bounds use purify--clone--trace (PCT).
This is a uniform worst-case theorem about the resources required to cross a
fidelity threshold; its \(1/\epsilon\) behavior describes the near-unit
boundary.  It does not determine the fidelity as a function of a fixed
nonzero output/input ratio, and copy-complexity optimality of PCT does not
imply its pointwise fidelity optimality.  Our primary theorems instead fix an
arbitrary \(t>0\), or equivalently an arbitrary \(\gamma>1\), and determine the
exact, generally nonunit, limiting fidelity at a fixed spectrum.  No
near-unit-fidelity hypothesis is imposed.  The sample-complexity statements
below are secondary inversions of these fixed-gain curves.

The PCT protocol first applies random purification
\cite{TangWrightZhandry2025}, then Werner's pure-state cloner on the
\(dr\)-dimensional purification space \cite{Werner1998}, and finally
discards the purifying registers
\cite{LiTheilHarrowChuang2026,
FanizzaGrinkoScharnhorstSpilecki2026,JeonSohnOh2026}.  Its published
data-processing guarantee is
\[
  \mathcal F_{\mathrm{PCT}}
  \ge
  \frac{\binom{N+dr-1}{dr-1}}
       {\binom{N+M+dr-1}{dr-1}}
  \longrightarrow
  (1+t^{-1})^{-(dr-1)}.
\]
The binomial ratio is the exact squared fidelity of the pure Werner step;
after the purifying registers are traced out it is only a lower bound on the
mixed-system fidelity.  Thus comparing this guarantee with our exact values
does not by itself order the actual channels.

Appendix~\ref{app:pct-comparison} makes that stronger comparison.  For every
rank-\(r\) state with simple positive spectrum and fixed \(\gamma>1\), it
proves an explicit upper bound
\(B_{\mathrm{PCT}}^{(r,d)}(\gamma,p)\) on the actual PCT root fidelity and,
for \(r>1\),
\[
  \limsup_{n\to\infty}F_{\mathrm{PCT}}
  \le B_{\mathrm{PCT}}^{(r,d)}(\gamma,p)
  <\Funiv^{(r,d)}(\gamma,p)<\Forb^{(r,d)}(\gamma,p).
\]
This is the central comparison: \(B_{\mathrm{PCT}}^{(r,d)}\) upper-bounds
the fidelity of the actual mixed-system PCT output and lies strictly below
the fidelity achieved by our rank-aware cloner.  Because \(\gamma>1\) is
arbitrary, the ordering holds across the full fixed-gain regime, including
when the limiting fidelity is bounded away from one; it is not inferred by
comparing our value with PCT's published lower guarantee.  We next translate
this ordering into the two sample-complexity statements most relevant here.

\paragraph{Projector states.}
Let \(\rho=P/r\), with \(1\le r<d\), and put
\(a=r(d-r)=\dim_{\mathbb C}\Gr(r,d)\).  Our exact Grassmannian theorem gives
\[
  \mathcal F_{\mathrm{ours}}^{\mathrm{proj}}(t)
  =(1+t^{-1})^{-a}.
\]
This is the complete proportional-limit curve for every fixed \(t>0\), not
a near-unit expansion.  The parameter \(\epsilon\) enters only when we
invert that curve into copy-complexity language.  In particular, the curve
reaches squared fidelity \(1-\epsilon\) at the exact ratio
\[
  \frac{N}{M}=t_{\mathrm{ours}}^{\mathrm{proj}}(\epsilon)
  :=\frac{1}{(1-\epsilon)^{-1/a}-1}.
\]
The appendix bounds the actual PCT squared fidelity from above by
\[
  U_{r,d}(t)
  :=(1+t^{-1})^{-a}
    \left(
      \frac{\sqrt{1+2t^{-1}}}{1+t^{-1}}
    \right)^{r-1}.
\]
Let \(t_B(\epsilon):=\inf\{t>0:U_{r,d}(t)\ge1-\epsilon\}\), and let
\[
  t_{\mathrm{guar}}(\epsilon)
  :=\frac{1}{(1-\epsilon)^{-1/(rd-1)}-1}
\]
be the ratio certified by the published PCT lower guarantee.  For
\(1<r<d\),
\[
  t_{\mathrm{ours}}^{\mathrm{proj}}(\epsilon)
  <t_B(\epsilon)\le t_{\mathrm{guar}}(\epsilon).
\]
In words, the upper bound on the \emph{actual} PCT output forces PCT to use
at least \(Mt_B(\epsilon)\) inputs in the proportional limit, whereas the
published guarantee certifies that \(Mt_{\mathrm{guar}}(\epsilon)\) inputs
suffice.  The three small-error expansions make the size of our advantage
explicit:
\begin{align*}
  t_{\mathrm{ours}}^{\mathrm{proj}}(\epsilon)
  &=\frac{a}{\epsilon}-\frac{a+1}{2}+O(\epsilon),\\
  t_B(\epsilon)
  &=\frac{a}{\epsilon}-\frac{a+1}{2}
    +\frac{r-1}{2a}+O(\epsilon),\\
  t_{\mathrm{guar}}(\epsilon)
  &=\frac{rd-1}{\epsilon}-\frac{rd}{2}+O(\epsilon).
\end{align*}
Thus, relative to the published PCT certificate, our exact projector result
improves the leading coefficient from \(rd-1\) to \(r(d-r)\).  Relative to
the actual PCT channel, the upper fidelity bound forces PCT to use at least
\(Mt_B(\epsilon)\) inputs.  Hence our cloner saves at least
\((r-1)M/[2r(d-r)]+O(M\epsilon)\) inputs within the iterated comparison.
The present bounds do not determine whether PCT's actual leading coefficient
equals \(r(d-r)\) or is larger.  For \(r=1\) the two protocols coincide,
while \(r=d\) is the singleton state \(I_d/d\) and needs no input copies.

\paragraph{A fixed simple full-rank state.}
Now let \(r=d\), fix
\(p_1>\cdots>p_d>0\), set \(q_{ij}=p_j/p_i\), and write
\(\sigma_{n,m_n}^{\mathrm{PCT}}(p)\) for the actual PCT output.  For every
fixed \(\delta>0\), with \(m_n/n\to1+\delta\), the actual-channel comparison
already gives
\[
  \limsup_{n\to\infty}
  F\!\left(\sigma_{n,m_n}^{\mathrm{PCT}}(p),
            \rho_p^{\ot m_n}\right)^2
  \le B_{\mathrm{PCT}}(1+\delta,p)^2
  <\Funiv(1+\delta,p)^2.
\]
Thus the fidelity advantage does not depend on a near-unit limit.  To
translate it into small-error copy-complexity notation, we expand our
unknown-spectrum fidelity and the actual-PCT upper bound at
\(\delta=M/N\downarrow0\):
\begin{align*}
  \Funiv(1+\delta,p)^2
  &=1-A_{\mathrm{ours}}(p)\delta^2+O_p(\delta^3),\\
  B_{\mathrm{PCT}}(1+\delta,p)^2
  &=1-A_{\mathrm{PCT}}(p)\delta^2+O_p(\delta^3),
\end{align*}
where
\begin{align*}
  A_{\mathrm{ours}}(p)
  &=\frac{d-1}{8}+\frac14\sum_{i<j}\frac{1}{q_{ij}},\\
  A_{\mathrm{PCT}}(p)
  &=\frac{d-1}{2}
    +\frac14\sum_{i<j}\frac{(1+q_{ij})^2}{q_{ij}},\\
  A_{\mathrm{PCT}}(p)-A_{\mathrm{ours}}(p)
  &=\frac{3(d-1)}{8}+\frac14\sum_{i<j}(2+q_{ij})>0.
\end{align*}
Since the actual PCT limiting fidelity is at most
\(B_{\mathrm{PCT}}\), inversion yields
\[
  \frac{N_{\mathrm{ours}}}{M}
  \sim\sqrt{\frac{A_{\mathrm{ours}}(p)}{\epsilon}},
  \qquad
  \frac{N_{\mathrm{PCT}}}{M}
  \gtrsim\sqrt{\frac{A_{\mathrm{PCT}}(p)}{\epsilon}}.
\]
Thus our cloner has the iterated pointwise scale
\(N/M\sim\sqrt{A_{\mathrm{ours}}(p)/\epsilon}\), whereas any PCT threshold
must satisfy \(N/M\gtrsim\sqrt{A_{\mathrm{PCT}}(p)/\epsilon}\).  Since
\(A_{\mathrm{PCT}}(p)>A_{\mathrm{ours}}(p)\), this proves a strict
coefficient-level advantage for our cloner.  The available bounds do not
determine PCT's exact pointwise \(\epsilon\)-scaling: its necessary
\(M/\sqrt\epsilon\) bound and its published sufficient
\(Md^2/\epsilon\) bound do not match.  If the spectrum is known, the still
smaller orbital coefficient is obtained by deleting \((d-1)/8\) from
\(A_{\mathrm{ours}}(p)\).  The matching worst-case lower bound over
rank-at-most-\(d\) states is driven by lower-rank projector states of rank
comparable to \(d/2\), not by this fixed interior spectrum.

\paragraph{Scope of the sample-complexity translation.}
The strict fixed-gain fidelity comparisons above hold for every fixed
\(\gamma>1\) and require neither \(\epsilon\downarrow0\) nor
\(\gamma\downarrow1\).  Only their small-\(\epsilon\) copy-complexity
translations are iterated-limit consequences: first \(N,M\to\infty\) at
fixed \(t=N/M\), and only then \(t\to\infty\).  Our present proportional
theorem is not uniform for a triangular sequence \(\gamma_N\downarrow1\),
so these inversions are not yet finite-sample theorems in the joint regime
treated by the two concurrent papers.  This qualification does not affect
the pointwise fixed-gain separation from actual PCT.  The inversions isolate
a next-order improvement for projector states and a strict
quadratic-coefficient separation from the PCT fidelity upper bound for every
fixed simple full-rank spectrum.  Full-rank spectra with positive
degeneracies remain outside the proved simple-spectrum comparison.

\subsection{Earlier and related work}\label{sec:related-work}

\paragraph{Pure-state cloning and state estimation.}
Universal qubit cloners were constructed in
\cite{BuzekHillery1996,GisinMassar1997,
BrussDiVincenzoEkertFuchsMacchiavelloSmolin1998}; the arbitrary-dimensional
global and single-clone optima were obtained in
\cite{Werner1998,KeylWerner1999}, while links to finite-copy estimation were
developed in \cite{MassarPopescu1995,BrussEkertMacchiavello1998}.
At the single-clone level, cloning with infinitely many outputs reduces to
estimation \cite{BaeAcin2006,ChiribellaDAriano2006}; global all-output
problems have also been studied for broad pure-state families
\cite{CheflesBarnett1999,ChiribellaYang2014Optimal,KimChitambar2022}.
These settings differ from our proportional-output mixed-state problem; see
\cite{ScaraniIblisdirGisinAcin2005,FanWangJingYueShiZhangMu2014} for reviews.

\paragraph{Finite-dimensional mixed-state cloning.}
The earlier finite-dimensional literature summarized above provides bounds
or constructions for restricted state families and criteria, rather than the
exact proportional-output minimax value under joint product fidelity
\cite{CiracEkertMacchiavello1999,Rastegin2003Upper,
Rastegin2002Relative,Rastegin2003StateDependent,Fan2003,FanLiuShi2007,
Kazakov2007,LemmWilde2017}.
More recently, random purification followed by pure-state cloning gave a
universal all-output guarantee for rank-constrained mixed states, but not an
exact optimum \cite{TangWrightZhandry2025,LiTheilHarrowChuang2026}; companion
work instead amplifies a selected eigenstate
\cite{LiTheilHarrowChuangQPA2026}.  Concurrent work proves that PCT has
optimal copy complexity up to constants for attaining a target squared
fidelity \(1-\epsilon\), uniformly over rank-at-most-\(r\) states
\cite{FanizzaGrinkoScharnhorstSpilecki2026,JeonSohnOh2026}.  This worst-case
fidelity-threshold optimality does not identify an exact fixed-gain fidelity
or imply pointwise optimality of PCT.  Our results impose no
near-unit-fidelity condition: for every fixed \(\gamma>1\), they determine
the explicit limiting fidelity on the proved fixed-rank strata.  By
analyzing the actual post-trace PCT output, rather than only its published
data-processing guarantee, we further prove that its asymptotic fidelity is
strictly smaller for every simple positive spectrum of rank \(r>1\) and
every flat projector orbit with \(r>1\).  Conversely, the concurrent
copy-complexity results are uniform in their joint small-error regime,
whereas our present
\(\gamma\downarrow1\) translation is an iterated limit.  See
\Cref{sec:pct-comparison} for the quantitative comparison.

\paragraph{Broadcasting and superbroadcasting.}
Single-input families are broadcastable precisely when they commute
\cite{BarnumCavesFuchsJozsaSchumacher1996}; an input-independent-fidelity
criterion likewise obstructs universal broadcasting for arbitrary mixed
qubits \cite{ChenChen2007MixedBroadcast}.  Multiple-input superbroadcasting
can nevertheless purify individual mixed-qubit marginals; universal,
phase-covariant, and optimal maps are known
\cite{DArianoMacchiavelloPerinotti2005,
BuscemiDArianoMacchiavelloPerinotti2005Optimal,
BuscemiDArianoMacchiavelloPerinotti2006,DangFan2007}, as is a
continuous-variable version for displaced thermal states
\cite{DArianoPerinottiSacchi2006Continuous}.  These marginal criteria allow
correlated outputs, so they do not determine full tensor-power fidelity.
Virtual broadcasting is a related nonphysical extension linking universal
cloning and optimal physical approximation, without changing this distinction
\cite{ParzygnatFullwoodBuscemiChiribella2024}.

\paragraph{Gaussian cloning and LAN.}
Continuous-variable cloning grew from coherent-state amplifier--beam-splitter
schemes \cite{CerfIpeRottenberg2000,
BraunsteinCerfIblisdirVanLoockMassar2001}; under joint rather than single-clone
fidelity, a non-Gaussian cloner can be optimal
\cite{CerfKrugerNavezWernerWolf2005}.  Caves established the
phase-insensitive amplifier noise limit \cite{Caves1982LinearAmplifiers}, and
Gu{\c t}{\u a}--Matsumoto reduced displaced mixed-Gaussian cloning to
least-noise amplification \cite{GutaMatsumoto2006MixedGaussianCloning}.
Quantum LAN was formulated statistically in \cite{GutaJencova2007LAN} and
developed through likelihood-ratio and explicit-channel formulations in
\cite{YamagataFujiwaraGill2013LAN,GutaKahn2006LANQubits,
KahnGuta2009LANFiniteDimensional}, yielding sharp purification, dilution, and
reversal results \cite{BowlesGutaAdesso2011Purification,
BowlesGutaAdesso2012Reversal}; related Gaussian tasks include amplifying and
purifying noisy coherent states \cite{ZhaoChiribella2017}.  This
LAN--amplification lineage underlies our converse.
Kahn--Gu{\c t}{\u a} supplies the interior antecedent, while the
fixed-rank theorem of Lahiry--Nussbaum gives the unified experiment used
here, including both its \(r=d\) specialization and the support--kernel
boundary modes when \(r<d\)
\cite{KahnGuta2009LANFiniteDimensional,LahiryNussbaum2024}.

\paragraph{Other cloning and replication models.}
Further variants include probabilistic exact cloning \cite{DuanGuo1998};
economical, phase-covariant, asymmetric, and telecloning schemes
\cite{NiuGriffiths1998,BrussCinchettiDArianoMacchiavello2000,
Cerf2000Asymmetric,MuraoJonathanPlenioVedral1999}; and cloning of
transformations and superreplication
\cite{ChiribellaDArianoPerinotti2008,ChiribellaYangYao2013,
GendraCalsamigliaMunozTapiaBaganChiribella2014,
DurSekatskiSkotiniotis2015,ChiribellaYangHuang2015,
ChiribellaYang2016}.  A higher-order framework for general channel cloning
includes state cloning \cite{SekatskiGuryanovaKothakondaSkotiniotis2025};
no-signalling yields general cloning bounds
\cite{SekatskiSkotiniotisDur2015NoSignaling,HuTomamichel2024}, while optimal
pure-state cloning and transposition can be complementary channels
\cite{BrzicGrinkoStudzinskiQuintino2026}.  These models differ from ours in
the success rule, family, target object, or replication rate.

\subsection{Organization}

\Cref{sec:label-kernels} develops the intrinsic positive-support
Schur--Cartan and label-processing ingredients.  These are
dimension-generic active-support lemmas used in the fixed-rank proof.
\Cref{sec:gaussian-LAN} similarly isolates the faithful thermal Gaussian and
LAN-transfer tools.  \Cref{sec:fixed-rank-proof} then proves
\cref{thm:known-optimum,thm:unknown-optimum} directly for every
\(1\le r\le d\): exact compression supplies the ambient
support--kernel factor, active-row label processing gives achievability, and
one thermal--vacuum fixed-rank LAN model gives both converses.  Setting
\(r=d\) simply removes the vacuum modes.  \Cref{sec:grassmann} proves the
exceptional flat projector orbit directly.  The PCT regimes are compared in
\cref{sec:pct-comparison}, which separates the near-unit fidelity-threshold
question from our arbitrary fixed-gain fidelity curves.  The actual-channel
bounds and strict fixed-\(\gamma\) ordering are proved in
Appendix~\ref{app:pct-comparison}.  The remaining appendices supply the PBW
positive-support calculation, the active-label local limit and randomized
rounding estimate, and the unified flat-prior Gaussian converse.

\subsection*{Acknowledgments}

% TODO(authors): add any further collaborators and JF funding support before submission.
We would like to thank Patrick Hayden for discussions and early collaborations.

\subsection*{Generative AI usage}

The original idea for the cloner construction was due to the authors.  We
then used OpenAI's ChatGPT Sol 5.6 and Anthropic's Fable 5 to assist with the
analysis.  Suggestions from these systems informed the evaluation of the
cloning fidelity and the converse argument.  We also used them for language
editing, manuscript organization, and figure preparation.  All AI-assisted
material was reviewed
and revised by the authors, who take full responsibility for the accuracy and
content of the manuscript.

\section*{Glossary of notation}
\addcontentsline{toc}{section}{Glossary of notation}

We collect the notation that recurs across sections.  Symbols local to a
single proof are defined where they first appear.

\begin{description}
\item[Spaces and maps.]
  \(\C\) and \(\R\) denote the complex and real numbers,
  \(\U(d)\) the unitary group, and \(\cT_1(\cH)\) the trace-class
  operators on a Hilbert space \(\cH\).  We write \(A^*\) for the adjoint,
  \(\Tr A\) for the trace, \(\|A\|_1\) for the trace norm, and \(\ot\) for
  tensor products.  A channel is a completely positive trace-preserving map.

\item[Asymptotic scale.]
  The ambient dimension \(d\) and rank \(r\) are fixed, \(n\) is the number of inputs,
  \(m_n\) the number of outputs, and
  \(\gamma_n=m_n/n\to\gamma>1\).  The number of pairs \(i<j\) is
  \(r_d=\binom d2\).  A capital \(N\) denotes a generic sample size.

\item[Spectra and states.]
  \(\mathsf A_r=\{x\in\R^r:\sum_i x_i=1\}\), and
  \(\Delta_r^\circ\) is its strictly ordered, positive spectral chamber.
  For \(p\in\Delta_r^\circ\), \(\bar p=(p,0_{d-r})\),
  \(\rho_{\bar p}=\diag(\bar p)\), and
  \(\rho_{\bar p,g}=g\rho_{\bar p}g^*\).  Bars are suppressed when \(r=d\).
  The notation \(K\Subset\Delta_r^\circ\) means that \(K\) is a compact
  subset of this open chamber; relative interiors and closures are taken in
  \(\mathsf A_r\).

\item[Fidelity and optimal values.]
  The unadorned \(F(A,C)\) denotes root (unsquared) fidelity; for classical
  measures it denotes Hellinger affinity.  With
  \(q_{ij}=\bar p_j/\bar p_i\), including \(q_{ij}=0\) when \(j>r\), the
  one-mode factor is \(f_\gamma(q_{ij})\).  The rank-aware values
  \(\Forb^{(r,d)}\) and \(\Funiv^{(r,d)}\) are defined in
  \eqref{eq:rank-orb-value}--\eqref{eq:rank-univ-value}; \(\Forb\),
  \(\Fcl\), and \(\Funiv\) denote the \(r=d\) abbreviations in
  \eqref{eq:main-values}.

\item[Schur--Weyl notation.]
  \(\mathsf Y_N\) is the set of partitions \(\lambda\vdash N\) with at
  most \(d\) rows.  The corresponding unitary and symmetric-group spaces are
  \(\cH_\lambda\) and \(\cS_\lambda\).  The Young-label law under
  \(\rho_p^{\ot N}\) is \(P_{N,p}\), its conditional irreducible state is
  \(\rho_{p,\lambda}\), and \(\mathsf T_N(p)\) is the typical-label set.
  On the rank-\(r\) face these objects carry a superscript \((r)\); all
  active labels have at most \(r\) rows and are padded by zeros in dimension
  \(d\).

\item[Cloners and label kernels.]
  \(\cM_n\) denotes a generic competing channel.  A Markov kernel
  \(K_n(\nu\mid\mu)\) on Young labels is lifted to the covariant channel
  \(\mathcal Q_n[K_n]\).  For compatible labels, \(\cC_{\mu,\nu}\) is the
  Cartan-sector channel; \(\cE_{\mu,\nu}\) extends it to all label pairs by
  a fallback channel.  The rank-aware unknown-spectrum cloner is
  \(\cC_n^{(r)}\); when \(r=d\) we also write
  \(\cC_n^{\mathrm{univ}}\).  Thus \(K\)
  denotes a spectral set, while \(K_n\) denotes a label kernel.

\item[Local coordinates.]
  \(\mathsf H_r=\{h\in\R^r:\sum_i h_i=0\}\) is the positive-spectral
  tangent hyperplane.  The active root set is
  \(\mathcal I_{r,d}=\{(i,j):1\le i\le r,\ i<j\le d\}\).
  The variables \(h\in\mathsf H_r\) and
  \(z\in\C^{\mathcal I_{r,d}}\) describe local spectral and orbital
  perturbations, respectively;
  \(X_p(z)\) is the source-normalized skew-Hermitian generator in the
  fixed-rank LAN theorem, \(\Omega_L^{(r,d)}\) is a bounded local box, and
  \(v_{n,L}^{(r)}(p)\) is its minimax value.  Alternative physical rotation
  coordinates are related by invertible, spectrum-dependent linear maps;
  whenever such a coordinate is used, the Gaussian displacement is
  reparametrized by the same map.

\item[Gaussian notation.]
  The calligraphic \(\cF\) denotes one-mode Fock space and is distinct from
  the roman \(F\) for fidelity.  The fixed-rank model has
  \(\binom r2\) thermal positive--positive modes and \(r(d-r)\) vacuum
  support--kernel modes.  The state \(\tau_p\) is the thermal product in the
  input and ideal target, while \(\tau_p^{(\gamma)}\) is the noisy amplifier
  output at zero displacement.  The states \(\Phi_z^{(p)}\) and
  \(\Psi_z^{(p,\gamma)}\) are the displaced input and ideal target built from
  \(\tau_p\).  The symbols \(\Theta^{\mathrm{in}}\) and
  \(\Theta^{\mathrm{tar}}\) include the classical Gaussian factor.  The
  number operator is \(\widehat N\), distinct from the sample size \(N\).

\item[Limits and uniformity.]
  Unless stated otherwise, asymptotics are as \(n\to\infty\) with \(d\) and
  \(\gamma\) fixed.  An \(o(1)\) term is uniform only over the parameters
  specified in the accompanying statement.  The fixed-rank results are not
  uniform when the rank changes, a positive spectral gap closes, or
  \(\gamma\) varies.
\end{description}

\section{Schur--Cartan ingredients on the positive support}
\label{sec:label-kernels}

Both fixed-rank cloning protocols use the same finite-dimensional
architecture.  This section establishes the ingredients intrinsic to a
simple positive spectrum.  The working dimension is denoted by \(d\), but
the statements are dimension-uniform: in the ambient rank-\(r\) problem of
\cref{sec:fixed-rank-proof}, they are invoked with the working dimension
replaced by \(r\).  Thus these are dimension-generic active-support lemmas
used in the fixed-rank proof.

We first recall the Schur--Weyl decomposition in
\cref{sec:schur-weyl-decomposition} and lift an arbitrary Markov kernel on
Young labels to a covariant quantum channel in
\cref{sec:cartan-sector-lifting}.  We then choose the exact target-law kernel
for known spectrum and the randomized dilation kernel for unknown spectrum in
\cref{sec:known-label-kernel,sec:unknown-label-kernel}, respectively, and
combine their label estimates with the common sector fidelity in
\cref{sec:uniform-estimates-achievability}.  Exact rank compression and the
ambient fixed-rank assembly are carried out in
\cref{sec:compression,sec:rank-aware-achievability}.  We now develop the
active-support components summarized in
\cref{fig:schur-cartan-cloner}.

\subsection{Schur--Weyl decomposition}
\label{sec:schur-weyl-decomposition}

Let
\[
  \mathsf Y_N=\{\lambda\vdash N:\ell(\lambda)\le d\}.
\]
Schur--Weyl duality gives
\[
  (\C^d)^{\ot N}
  \cong
  \bigoplus_{\lambda\in\mathsf Y_N}\cH_\lambda\ot\cS_\lambda,
\]
where \(\cH_\lambda\) is the irreducible \(\U(d)\)-module of highest
weight \(\lambda\) and \(\cS_\lambda\) is the corresponding Specht module.
For \(\rho_p=\diag(p_1,\ldots,p_d)\),
\begin{equation}\label{eq:schur-state}
  \rho_p^{\ot N}
  =\bigoplus_{\lambda\in\mathsf Y_N}
  P_{N,p}(\lambda)\,
  \rho_{p,\lambda}\ot\frac{I_{\cS_\lambda}}{\dim\cS_\lambda},
  \qquad
  P_{N,p}(\lambda)=\dim\cS_\lambda\,s_\lambda(p),
  \quad
  \rho_{p,\lambda}=\frac{\pi_\lambda(\rho_p)}{s_\lambda(p)}.
\end{equation}
For \(\rho_{p,g}=g\rho_p g^*\), replace \(\rho_{p,\lambda}\) by
\(\rho_{p,g,\lambda}=U_\lambda(g)\rho_{p,\lambda}U_\lambda(g)^*\).  The
Young law \(P_{N,p}\) depends on \(p\), but not on \(g\).

Define
\[
  \mathsf T_N(p)
  =\left\{\lambda\in\mathsf Y_N:
  \left\|\frac\lambda N-p\right\|_\infty\le N^{-1/3}\right\}.
\]

\begin{lemma}[Uniform Young concentration]\label{lem:young-concentration}
There is a constant \(c_d>0\), depending only on \(d\), such that
\[
  \sup_{p\in\Delta_d^\circ}P_{N,p}(\mathsf T_N(p)^c)
  \le e^{-c_d N^{1/3}}
\]
for all sufficiently large \(N\).
\end{lemma}

\begin{proof}
Write \(q=\lambda/N\).  The Weyl dimension formula gives
\(\dim\cH_\lambda\le(N+1)^{d(d-1)/2}\).  Every weight \(\xi\) of
\(\cH_\lambda\) has \(\lambda-\xi\) equal to a nonnegative integral
combination of positive roots.  Since \(p_j/p_i<1\) for \(i<j\), this gives
\(p^\xi\le p^\lambda\).  Hence
\(s_\lambda(p)\le \dim\cH_\lambda\,p^\lambda\).
The standard bounds
\(\dim\cS_\lambda\le\binom{N}{\lambda_1,\ldots,\lambda_d}\) and
\(\binom{N}{\lambda_1,\ldots,\lambda_d}p^\lambda
\le \exp[-ND(q\|p)]\) give
\begin{equation}\label{eq:finite-Young-large-deviation}
  P_{N,p}(\lambda)
  \le (N+1)^{d(d-1)/2}\exp[-ND(q\|p)],
\end{equation}
which is the finite-sample estimate underlying the Keyl--Werner spectrum
large-deviation theorem \cite{KeylWerner2001Spectrum}.

There are at most \((N+1)^{d-1}\) partitions of \(N\) with at most \(d\)
rows.  If \(\lambda\notin\mathsf T_N(p)\), then
\(\|q-p\|_1>N^{-1/3}\), and Pinsker's inequality gives
\(D(q\|p)\ge\tfrac12N^{-2/3}\).  Therefore
\[
  P_{N,p}(\mathsf T_N(p)^c)
  \le (N+1)^{d(d-1)/2+d-1}e^{-N^{1/3}/2}.
\]
For all sufficiently large \(N\), the polynomial factor is at most
\(e^{N^{1/3}/4}\).  Hence the asserted bound holds with \(c_d=1/4\).
\end{proof}

\subsection{Cartan-sector channel and lifting of a label kernel}
\label{sec:cartan-sector-lifting}

Let \(\lambda,\theta\) be dominant weights, with unit highest-weight vectors
\(v_\lambda,v_\theta\).  The Cartan component \(\cH_{\lambda+\theta}\)
occurs with multiplicity one in \(\cH_\lambda\ot\cH_\theta\).  Let
\begin{equation}\label{eq:cartan-inclusion}
  V_{\lambda,\theta}:\cH_{\lambda+\theta}
  \longrightarrow\cH_\lambda\ot\cH_\theta,
  \qquad
  V_{\lambda,\theta}v_{\lambda+\theta}=v_\lambda\ot v_\theta,
\end{equation}
be the normalized intertwining isometry.  For \(\mu\in\mathsf Y_n\) and
\(\nu\in\mathsf Y_m\), write \(\nu-\mu\vdash m-n\) when the coordinatewise
difference \(\nu-\mu\) is a partition of \(m-n\).  If this holds, put
\(\omega=\nu-\mu\) and define
\begin{equation}\label{eq:sector-channel-formula}
  \cC_{\mu,\nu}(X)
  =\frac{\dim\cH_\mu}{\dim\cH_\nu}
  V_{\mu,\omega}^*(X\ot I_{\cH_\omega})V_{\mu,\omega}.
\end{equation}
The partial trace
\(\Tr_{\cH_\omega}(V_{\mu,\omega}V_{\mu,\omega}^*)\) commutes with the
irreducible \(\U(d)\)-action on \(\cH_\mu\); hence Schur's lemma and a trace
comparison give
\(
  \Tr_{\cH_\omega}(V_{\mu,\omega}V_{\mu,\omega}^*)
  =\frac{\dim\cH_\nu}{\dim\cH_\mu}I_{\cH_\mu}
\).
Thus \eqref{eq:sector-channel-formula} is completely positive by
tensoring with \(I_{\cH_\omega}\), applying
\(Y\mapsto V_{\mu,\omega}^*YV_{\mu,\omega}\), and positive rescaling;
the identity above makes it trace
preserving, while equivariance of \(V_{\mu,\omega}\) gives
\(\U(d)\)-covariance.

For every pair \((\mu,\nu)\in\mathsf Y_n\times\mathsf Y_m\), define a
covariant channel \(\cE_{\mu,\nu}:\cT_1(\cH_\mu)\to\cT_1(\cH_\nu)\) by
\[
  \cE_{\mu,\nu}(X)=
  \begin{cases}
    \cC_{\mu,\nu}(X),&\nu-\mu\vdash m-n,\\[1mm]
    \Tr(X)I_{\cH_\nu}/\dim\cH_\nu,&\text{otherwise}.
  \end{cases}
\]
Let \(K_n:\mathsf Y_n\rightsquigarrow\mathsf Y_{m_n}\), with transition
probabilities \(K_n(\nu\mid\mu)\), be any Markov kernel.  Its Schur--Cartan
lifting is the following channel,
denoted by \(\mathcal Q_n[K_n]\): apply the Schur transform, measure
\(\mu\), discard the input Specht factor, sample \(\nu\) from
\(K_n(\cdot\mid\mu)\), apply \(\cE_{\mu,\nu}\), adjoin
\(I_{\cS_\nu}/\dim\cS_\nu\), and apply the inverse Schur transform.  Each
step is completely positive and trace preserving, so their composition is a
channel.  Because the label measurement and sampling are invariant and every
map \(\cE_{\mu,\nu}\) is covariant, \(\mathcal Q_n[K_n]\) is
\(\U(d)\)-covariant.

\begin{proposition}[Fidelity of a lifted label kernel]
\label{prop:lifted-kernel-bound}
Fix \(p\in\Delta_d^\circ\).  Let
\[
  \Pi_{n,p}(\mu,\nu)=P_{n,p}(\mu)K_n(\nu\mid\mu),
  \qquad
  Q_{m_n,p}(\nu)=\sum_\mu\Pi_{n,p}(\mu,\nu).
\]
Suppose \(G\subset\mathsf Y_n\times\mathsf Y_{m_n}\) satisfies
\(\nu-\mu\vdash m_n-n\) for every \((\mu,\nu)\in G\), and
\[
  F\!\left(\cC_{\mu,\nu}(\rho_{p,\mu}),\rho_{p,\nu}\right)\ge b
  \qquad((\mu,\nu)\in G)
\]
for some \(b\in[0,1]\).  Then, for every \(g\in\U(d)\),
\begin{equation}\label{eq:lifted-kernel-bound}
  F\!\left(
    \mathcal Q_n[K_n](\rho_{p,g}^{\ot n}),
    \rho_{p,g}^{\ot m_n}
  \right)
  \ge b\left(F(Q_{m_n,p},P_{m_n,p})-\sqrt{\Pi_{n,p}(G^c)}\right).
\end{equation}
\end{proposition}

\begin{proof}
Covariance makes every sector fidelity independent of \(g\).  In the output
Schur decomposition, the actual state in the \(\nu\)-block is the mixture of
\(\cE_{\mu,\nu}(\rho_{p,g,\mu})\) with weights
\(\Pi_{n,p}(\mu,\nu)\), whereas the ideal block is
\(P_{m_n,p}(\nu)\rho_{p,g,\nu}\).  The Specht factors are identical
and therefore cancel by multiplicativity of the trace norm under tensor
products.

For \(Q_{m_n,p}(\nu)>0\), set
\[
  r(\nu)=\sum_{\mu:(\mu,\nu)\notin G}
  \frac{\Pi_{n,p}(\mu,\nu)}{Q_{m_n,p}(\nu)},
\]
and set \(r(\nu)=0\) otherwise.  The direct-sum identity for fidelity,
its scaling under positive scalar multiples, and joint concavity give
\[
\begin{aligned}
  &F\!\left(
    \mathcal Q_n[K_n](\rho_{p,g}^{\ot n}),
    \rho_{p,g}^{\ot m_n}
  \right)\\
  &\qquad\ge
  b\left(
    F(Q_{m_n,p},P_{m_n,p})
    -\sum_\nu
      \sqrt{Q_{m_n,p}(\nu)P_{m_n,p}(\nu)}\,r(\nu)
  \right).
\end{aligned}
\]
Since \(0\le r\le1\) and
\(\sum_\nu Q_{m_n,p}(\nu)r(\nu)=\Pi_{n,p}(G^c)\),
Cauchy--Schwarz yields
\[
  \sum_\nu\sqrt{Q_{m_n,p}(\nu)P_{m_n,p}(\nu)}\,r(\nu)
  \le
  \sqrt{\Pi_{n,p}(G^c)}
  \left(\sum_\nu P_{m_n,p}(\nu)r(\nu)\right)^{1/2}
  \le\sqrt{\Pi_{n,p}(G^c)}.
\]
This proves \eqref{eq:lifted-kernel-bound}.
\end{proof}

We now specialize the lifted channel by choosing its output-label kernel.
The two protocols differ only in this choice: the known-spectrum construction
samples the exact target law, whereas the unknown-spectrum construction uses
a randomized dilation of the measured input label.

\subsection{Known spectrum: exact target-label sampling}
\label{sec:known-label-kernel}

Fix \(p\in\Delta_d^\circ\).  For every \(n\) such that \(m_n>n\), define
\[
  K_n^p(\nu\mid\mu)=P_{m_n,p}(\nu).
\]
Thus \(\nu\) is independent of \(\mu\) and has the exact target law.  Define
the amplification branch by
\[
  \cC_n^p:=\mathcal Q_n[K_n^p].
\]
Pairs with
\(\nu-\mu\not\vdash m_n-n\) use
\(\cE_{\mu,\nu}(X)=\Tr(X)I_{\cH_\nu}/\dim\cH_\nu\) without changing the
sampled label \(\nu\).  Together with the common convention in
\cref{sec:main-results}, this defines the complete \(\U(d)\)-covariant sequence
\((\cC_n^p)_n\), whose output Young law on \(\rho_{p,g}^{\ot n}\) is
\(P_{m_n,p}\) for every \(n\) and \(g\).

\subsection{Unknown spectrum: randomized rounding}
\label{sec:unknown-label-kernel}

For \(N\in\mathbb N\), let
\[
  \mathsf L_N=\left\{x\in\mathbb Z^d:\sum_i x_i=N\right\}.
\]
For every \(n\) such that \(m_n>n\), set
\(\nu_n^{\mathrm{fb}}=(m_n,0,\ldots,0)\).
Given \(\mu\in\mathsf Y_n\), draw independent variables
\(U_1,\ldots,U_{d-1}\), each uniform on \([-1/2,1/2)\), and set
\begin{equation}\label{eq:randomized-rounding}
  \widetilde\nu_i
  =\left\lfloor\gamma_n(\mu_i+U_i)+\frac12\right\rfloor
  \quad(1\le i<d),
  \qquad
  \widetilde\nu_d=m_n-\sum_{i<d}\widetilde\nu_i.
\end{equation}
This is a spectrum-independent lattice version of the dilation
\(\mu\mapsto\gamma_n\mu\): it sends typical labels near \(np\) to labels near
\(m_np\), with \(\widetilde\nu-\mu\) typically a partition.  The uniform
shifts randomize the rounding error, avoiding lattice bias in the Gaussian
limit.
Then \(\widetilde\nu\in\mathsf L_{m_n}\).  Let
\(D_n(\widetilde\nu\mid\mu)\) denote its distribution.  Define
\(K_n^{\mathrm{rnd}}:\mathsf Y_n\rightsquigarrow\mathsf Y_{m_n}\) by
\begin{equation}\label{eq:rounding-kernel}
  K_n^{\mathrm{rnd}}(\nu\mid\mu)=
  \begin{cases}
    D_n(\nu\mid\mu),
      &\nu\ne\nu_n^{\mathrm{fb}}
       \text{ and }\nu-\mu\vdash m_n-n,\\[1mm]
    1-\displaystyle\sum_{\substack{\lambda\in\mathsf Y_{m_n},\\
            \lambda\ne\nu_n^{\mathrm{fb}}\\
            \lambda-\mu\vdash m_n-n}}
       D_n(\lambda\mid\mu),
      &\nu=\nu_n^{\mathrm{fb}},\\[4mm]
    0,&\text{otherwise}.
  \end{cases}
\end{equation}
Under \eqref{eq:rounding-kernel}, every rounded label outside
\(\mathsf Y_{m_n}\), and every
\(\widetilde\nu\in\mathsf Y_{m_n}\) for which
\(\widetilde\nu-\mu\not\vdash m_n-n\), is replaced by
\(\nu_n^{\mathrm{fb}}\), which also retains any mass rounded onto it.
Define the amplification branch by
\[
  \cC_n^{\mathrm{univ}}:=\mathcal Q_n[K_n^{\mathrm{rnd}}].
\]
The rounding variables, \eqref{eq:rounding-kernel}, and the maps
\(\cE_{\mu,\nu}\) are spectrum independent, while
\(\mathcal Q_n[K_n^{\mathrm{rnd}}]\) is \(\U(d)\)-covariant.  Together
with the common convention in \cref{sec:main-results}, this defines the
complete spectrum-independent, \(\U(d)\)-covariant sequence
\((\cC_n^{\mathrm{univ}})_n\).

\subsection{Sector and label asymptotics}
\label{sec:uniform-estimates-achievability}

The following proposition is the only representation-theoretic asymptotic
estimate used in the cloning protocols.

\begin{proposition}[Uniform Cartan-sector fidelity]
\label{prop:sector-fidelity}
Let \(K\Subset\Delta_d^\circ\) and let \(\delta_n\ge0\) satisfy
\(\delta_n\to0\).  For all sufficiently large \(n\), every
\(p\in K\), \(\mu\in\mathsf Y_n\), and \(\nu\in\mathsf Y_{m_n}\) satisfying
\[
  \left\|\frac\mu n-p\right\|_\infty\le\delta_n,
  \qquad
  \left\|\frac\nu{m_n}-p\right\|_\infty\le\delta_n
\]
also satisfies \(\nu-\mu\vdash m_n-n\).  Moreover,
\[
  F\!\left(\cC_{\mu,\nu}(\rho_{p,\mu}),\rho_{p,\nu}\right)
  =\Forb(\gamma,p)+o(1)
\]
uniformly over all such \(p,\mu,\nu\).
\end{proposition}

The mechanism is already visible on states obtained by a fixed number of
lowering operations.  Writing \(\omega=\nu-\mu\), the Cartan intertwiner
satisfies
\[
  V_{\mu,\omega}F_\alpha^{(\nu)}
  =\bigl(F_\alpha^{(\mu)}\ot I+I\ot F_\alpha^{(\omega)}\bigr)
    V_{\mu,\omega},
\]
where \(F_\alpha^{(\lambda)}\) is the lowering operator for
\(\alpha=e_i-e_j\) on \(\cH_\lambda\).  After normalization, each excitation
splits binomially between the two factors, with fraction
\[
  s_{\alpha,n}
  =\frac{\mu_i-\mu_j}{\nu_i-\nu_j}
  =\frac1\gamma+o(1)
\]
entering the \(\mu\)-factor.  Thus the conditional thermal parameter
\(q_{ij}=p_j/p_i\) is transformed into
\[
  1-s_{\alpha,n}+s_{\alpha,n}q_{ij}
  \longrightarrow
  1-\frac{1-q_{ij}}{\gamma}
  =:q_{ij}^{(\gamma)}.
\]
The low-weight matrices of the sector output and ideal sector therefore
approach those of \(\tau_p^{(\gamma)}\) and \(\tau_p\), whose fidelity is
\(\Forb(\gamma,p)\).  This is a statement about conditional states for
typical label pairs, not convergence of the varying channels
\(\cC_{\mu,\nu}\) to the Gaussian amplifier.  The exact PBW calculation and
uniform thermal-tail argument that complete the proof are given in
Appendix~\ref{app:sector-fidelity}.

The second estimate concerns only \(K_n^{\mathrm{rnd}}\).

\begin{proposition}[Randomized Young-label dilation]
\label{prop:young-rounding}
There is a constant \(c_{\mathrm{rnd}}<\infty\), depending only on \(d\) and
the fixed sequence \((\gamma_n)\), such that, for every compact
\(K\Subset\Delta_d^\circ\) and all sufficiently large \(n\), every
\(p\in K\) and \(\mu\in\mathsf T_n(p)\) satisfy
\[
  \operatorname{supp}K_n^{\mathrm{rnd}}(\,\cdot\mid\mu)
  \subseteq
  \left\{\nu\in\mathsf Y_{m_n}:
    \nu-\mu\vdash m_n-n,\quad
    \|\nu-\gamma_n\mu\|_\infty\le c_{\mathrm{rnd}}
  \right\}.
\]

For \(m_n>n\), let
\[
  Q_{m_n,p}(\nu)
  :=\sum_{\mu\in\mathsf Y_n}
    P_{n,p}(\mu)K_n^{\mathrm{rnd}}(\nu\mid\mu),
  \qquad \nu\in\mathsf Y_{m_n},
\]
be the induced output-label law.  Then, for every compact
\(K\Subset\Delta_d^\circ\),
\begin{equation}\label{eq:rounding-label-fidelity}
  F(Q_{m_n,p},P_{m_n,p})
  =\Fcl(\gamma,d)+o(1)
  \qquad\text{uniformly for }p\in K.
\end{equation}
\end{proposition}

The uniform dither in \eqref{eq:randomized-rounding} makes the label kernel a
discrete realization of continuous dilation.  After centering at \(Np\) and
scaling by \(\sqrt N\), embed a label law as a piecewise-constant density on
the fundamental cells of the label lattice.  Before fallback replacement,
the intersection-volume identity \eqref{eq:cell-kernel} identifies the
embedded rounded law with the target-cell average of the
\(\sqrt{\gamma_n}\)-dilation of the embedded input law.  Write
\[
  \Sigma_p
  :=\left.(\diag(p)-pp^{\mathsf T})\right|_{
    \{x\in\R^d:\,\sum_i x_i=0\}}.
\]
The Young local limit is \(\mathsf N_{0,\Sigma_p}\), so randomized rounding changes the limit
to \(\mathsf N_{0,\gamma\Sigma_p}\), whereas the ideal \(m_n\)-label law
still converges to \(\mathsf N_{0,\Sigma_p}\).  Their Hellinger affinity is
\[
  F\!\left(\mathsf N_{0,\gamma\Sigma_p},
           \mathsf N_{0,\Sigma_p}\right)
  =\left(\frac{2\sqrt\gamma}{1+\gamma}\right)^{(d-1)/2}
  =\Fcl(\gamma,d).
\]
The support check and fallback rule affect only vanishing probability on
typical labels.

The proofs of \cref{prop:sector-fidelity,prop:young-rounding} are given in
Appendices~\ref{app:sector-fidelity} and~\ref{app:young-rounding},
respectively.

\section{Gaussian and LAN ingredients on the positive support}
\label{sec:gaussian-LAN}

We next isolate the Gaussian dilation and LAN tools needed by the fixed-rank
converse.  As in the preceding section, \(d\) here denotes a generic
positive-support dimension.  The Gaussian result treats the faithful thermal
modes within that support, while the unified model in
\cref{sec:fixed-rank-LAN} combines them with the support--kernel vacuum modes
and performs the ambient fixed-rank transfer.  We first introduce the local
coordinates, then solve the Gaussian problems in
\cref{sec:gaussian-amplification}, and finally record the general fidelity
transfer used with fixed-rank LAN in \cref{sec:local-asymptotic-normality}.

The LAN transport common to these converse arguments is summarized in
\cref{fig:LAN-transfer}.  The upper path is the induced Gaussian-model
competitor; the lower path is the ideal target against which its fidelity is
measured.

\begin{figure}[!htbp]
\centering
\begin{tikzpicture}[
  flow/.style={-{Latex[length=2mm]},thick},
  lan/.style={-{Latex[length=2mm]},thick,dashed,teal!70!black},
  compare/.style={{Latex[length=1.8mm]}-{Latex[length=1.8mm]},densely dashed},
  box/.style={draw,rounded corners=2pt,align=center,inner sep=5pt,
    minimum height=10mm,font=\small},
  finite/.style={box,fill=orange!9},
  gaussian/.style={box,fill=blue!6}
]
  \node[gaussian,text width=22mm] (gin) at (0,1.25)
    {Gaussian input\\\(\Phi_\theta\)};
  \node[finite,text width=25mm] (fin) at (3.35,1.25)
    {finite input\\\(\rho_{n,\theta}\)};
  \node[finite,text width=29mm] (factual) at (7.15,1.25)
    {competitor output\\\(\cM_n(\rho_{n,\theta})\)};
  \node[gaussian,text width=31mm] (gactual) at (11.35,1.25)
    {induced output\\\(\Gamma_n(\Phi_\theta)\)};

  \node[finite,text width=29mm] (ftarget) at (7.15,-1.25)
    {finite ideal target\\\(\sigma_{n,\theta}\)};
  \node[gaussian,text width=31mm] (gtarget) at (11.35,-1.25)
    {Gaussian ideal target\\\(\Psi_\theta\)};

  \draw[lan] (gin) -- node[above,font=\scriptsize,text=black] {\(S_n\)} (fin);
  \draw[flow] (fin) -- node[above,font=\scriptsize] {\(\cM_n\)} (factual);
  \draw[flow] (factual) -- node[above,font=\scriptsize] {\(T_n\)} (gactual);
  \draw[lan] (ftarget) -- node[above,font=\scriptsize,text=black] {\(T_n\)}
    (gtarget);
  \draw[compare] (factual) -- node[left,font=\scriptsize] {fidelity}
    (ftarget);
  \draw[compare] (gactual) -- node[right,font=\scriptsize] {fidelity}
    (gtarget);

  \node[draw,rounded corners=2pt,fill=gray!6,align=center,font=\small,
    inner sep=5pt,text width=92mm] at (5.7,-3.0)
    {\(\Gamma_n=T_n\circ\cM_n\circ S_n\).  Hence a Gaussian minimax
     upper bound also bounds the finite-dimensional cloning fidelity.\\[-1pt]
     For cloning, \(\theta=(h,z)\) and the target dilation is
     \((h,z)\mapsto\sqrt\gamma\,(h,z)\); for known spectrum, set \(h=0\).};
\end{tikzpicture}
\caption{Transporting an arbitrary finite-dimensional cloner through
two-way local asymptotic normality.  A dashed arrow identifies its displayed
endpoint up to uniform \(o(1)\) trace-norm error on compact local parameter
sets; the solid \(T_n\) arrow defines the induced output exactly.  The LAN maps
are not exact inverses, and the induced channels \(\Gamma_n\) need not be
Gaussian or converge.  The converse follows from the optimal Gaussian
dilation and amplification bounds.}
\label{fig:LAN-transfer}
\end{figure}

\subsection{Local coordinates}
\label{sec:local-coordinates}

Let
\[
  \mathsf H=\left\{h\in\R^d:\sum_i h_i=0\right\},
  \qquad z\in\C^{r_d},
\]
and, for \(p\in\Delta_d^\circ\), define the source-normalized generator
\[
  X_p(z)=\sum_{j<k}
    \frac{z_{jk}E_{kj}-\bar z_{jk}E_{jk}}{\sqrt{p_j-p_k}},
  \qquad
  g_{n,z}^{(p)}=\exp\!\left(\frac{X_p(z)}{\sqrt n}\right),
\]
where \(E_{jk}\) are the standard matrix units.  In
\(\rho_{p+h/\sqrt n,g_{n,z}^{(p)}}\), the coordinates \(h\) and \(z\)
describe local changes of eigenvalues and eigenbasis, respectively.
Since \(X_p(z)^*=-X_p(z)\), one has
\(g_{n,z}^{(p)}\in\U(d)\).  For \(L>0\), set
\[
  \Omega_L
  =\left\{(h,z)\in\mathsf H\times\C^{r_d}:
  \|h\|_\infty\le L,\
  |\Re z_{ij}|,|\Im z_{ij}|\le L\ \text{for all }i<j\right\}.
\]
For every fixed \(p\in\Delta_d^\circ\) and \(L>0\), the vector
\(p+h/\sqrt n\) belongs to \(\Delta_d^\circ\) for every
\((h,z)\in\Omega_L\) once \(n\) is sufficiently large.  Under LAN, \(h\)
becomes a \((d-1)\)-dimensional classical Gaussian shift, while \(z\) becomes
a collection of displaced oscillator coordinates.  The Gaussian theorem
below optimizes the two corresponding fidelity losses jointly.

\subsection{Gaussian amplification}
\label{sec:gaussian-amplification}

Let \(\cF\) be the one-mode Fock space with number basis
\(\{\ket{k}:k\ge0\}\), creation and annihilation operators \(a^*,a\), and
number operator \(\widehat N=a^*a\).  We use the Weyl convention
\begin{equation}\label{eq:Weyl-convention}
  D(z)=\exp(z a^*-\bar z a).
\end{equation}
For \(0<q<1\), set
\[
  \tau_q=(1-q)\sum_{k\ge0}q^k\proj{k}.
\]
Two thermal states commute, and hence
\begin{equation}\label{eq:thermal-root-fidelity}
  F(\tau_q,\tau_{q'})
  =\frac{\sqrt{(1-q)(1-q')}}{1-\sqrt{qq'}}.
\end{equation}
For \(G\ge1\), put \(q^{(G)}=1-(1-q)/G\).  The quantum-limited
phase-insensitive amplifier \(\cA_G\) satisfies
\begin{equation}\label{eq:amplifier-action}
  \cA_G\!\left(D(z)\tau_qD(z)^*\right)
  =D(\sqrt G\,z)\tau_{q^{(G)}}D(\sqrt G\,z)^*.
\end{equation}
Equivalently,
\(\cA_G(X)=\Tr_b[S_r(X\ot\proj{0}_b)S_r^*]\), where \(S_r\) is a
two-mode squeezing unitary with \(\cosh^2 r=G\).  Squeezing supplies the
gain \(\sqrt G\), while tracing out the vacuum idler produces the minimum
added noise encoded by \(q^{(G)}\).

For \(p\in\Delta_d^\circ\), define on
\(\cF_{\mathrm{orb}}=\cF^{\ot r_d}\)
\[
  \tau_p=\bigotimes_{i<j}\tau_{q_{ij}},
  \qquad
  \tau_p^{(\gamma)}
  =\bigotimes_{i<j}\tau_{q_{ij}^{(\gamma)}}.
\]
With \(D(z)=\bigotimes_{i<j}D(z_{ij})\), the three relevant
states are
\begin{align*}
  \text{input:}\quad
  &\Phi_z^{(p)}=D(z)\tau_pD(z)^*,\\
  \text{ideal gain-\(\sqrt\gamma\) target:}\quad
  &\Psi_z^{(p,\gamma)}=D(\sqrt\gamma z)\tau_pD(\sqrt\gamma z)^*,\\
  \text{amplifier output:}\quad
  &\cA_\gamma^{\ot r_d}(\Phi_z^{(p)})
    =D(\sqrt\gamma z)\tau_p^{(\gamma)}D(\sqrt\gamma z)^*.
\end{align*}
Thus the ideal target retains the thermal state \(\tau_p\); only the actual
amplifier output contains the additional noise \(\tau_p^{(\gamma)}\).  By
\eqref{eq:thermal-root-fidelity},
\begin{equation}\label{eq:one-mode-f}
  F(\tau_{q^{(\gamma)}},\tau_q)=f_\gamma(q).
\end{equation}
Thus, for every \(z\in\C^{r_d}\), the product-amplifier output and the
corresponding ideal state satisfy
\[
  F\!\left(\cA_\gamma^{\ot r_d}(\Phi_z^{(p)}),
           \Psi_z^{(p,\gamma)}\right)
  =\Forb(\gamma,p).
\]

Equip \(\mathsf H\) with its induced Euclidean measure and recall
\[
  \Sigma_p=\left.(\diag(p)-pp^{\mathsf T})\right|_{\mathsf H}>0.
\]
Write \(\mathsf N_{h,\Sigma_p}\) for the Gaussian law on \(\mathsf H\)
with mean \(h\) and covariance \(\Sigma_p\).  A
classical--quantum state has the form
\[
  R(x)=r(x)\rho_x,\qquad
  r(x)\ge0,\quad \int_{\mathsf H}r(x)\,\dd x=1,
\]
where \(x\) and \(r\) are classical, while \(\rho_x\) is a quantum state on
\(\cF_{\mathrm{orb}}\).  A channel \(\Lambda\) is completely positive and
preserves \(\int\Tr R(x)\,\dd x\).  Here complete positivity
means positivity of \(\operatorname{id}_{M_k}\otimes\Lambda\) for every
\(k\ge1\), with the matrix factor acting on the quantum trace-class fibers;
the classical algebra is abelian.  If
\(S(x)=s(x)\sigma_x\), then
\[
  F(R,S)=\int_{\mathsf H}F(R(x),S(x))\,\dd x
  =\int_{\mathsf H}\sqrt{r(x)s(x)}\,F(\rho_x,\sigma_x)\,\dd x.
\]
With \(\|R\|_1:=\int_{\mathsf H}\|R(x)\|_1\,\dd x\), the usual monotonicity
and joint-concavity properties remain valid, and the first equality above is
the continuous direct-sum identity for fidelity.

\par\indent
Set
\[
  \Theta_{h,z}^{\mathrm{in}}
  =\mathsf N_{h,\Sigma_p}\ot\Phi_z^{(p)},
  \qquad
  \Theta_{h,z}^{\mathrm{tar}}
  =\mathsf N_{\sqrt\gamma h,\Sigma_p}\ot\Psi_z^{(p,\gamma)}.
\]
Let
\[
  \mathsf Z_L
  =\{z\in\C^{r_d}:|\Re z_{ij}|,|\Im z_{ij}|\le L\ \text{for all }i<j\}.
\]

\begin{theorem}[Gaussian amplification optimum]
\label{thm:gaussian-amplification}
For every \(p\in\Delta_d^\circ\),
\begin{align}
  \lim_{L\to\infty}\sup_\Lambda
  \frac1{|\mathsf Z_L|}\int_{\mathsf Z_L}
  F\!\left(\Lambda(\Phi_z^{(p)}),\Psi_z^{(p,\gamma)}\right)\,\dd z
  &=\Forb(\gamma,p),
  \label{eq:gaussian-orbital-optimum}\\
  \lim_{L\to\infty}\sup_\Lambda
  \frac1{|\Omega_L|}\int_{\Omega_L}
  F\!\left(\Lambda(\Theta_{h,z}^{\mathrm{in}}),
            \Theta_{h,z}^{\mathrm{tar}}\right)\,\dd h\,\dd z
  &=\Funiv(\gamma,p).
  \label{eq:gaussian-full-optimum}
\end{align}
The first supremum is over all quantum channels on
\(\cF_{\mathrm{orb}}\); the second is over all classical--quantum channels.
No Gaussianity, covariance, or product-structure assumption is imposed on
the competing channels in either supremum.
\end{theorem}

The achievability mechanism is explicit.  The product amplifier
\(\cA_\gamma^{\ot r_d}\) attains \(\Forb(\gamma,p)\) in the orbital model.
In the hybrid classical--quantum model, tensor it with the classical pushforward
\(x\mapsto\sqrt\gamma x\).  The resulting Gaussian affinity is
\(\Fcl(\gamma,d)\), so multiplicativity gives
\(\Fcl(\gamma,d)\Forb(\gamma,p)=\Funiv(\gamma,p)\).

For the converse, average a near-optimal channel over translates of the
expanding flat-prior boxes.  The vanishing boundary-to-volume error yields a
displacement-covariant limiting map.  A weighted-fidelity witness then turns
the orbital fidelity into a noise moment, which the least-noise inequality
bounds by \(\Forb(\gamma,p)\).  In the hybrid model, covariance forces the
classical trace kernel to have convolution form, and the corresponding
Gaussian witness contributes at most \(\Fcl(\gamma,d)\).  The two bounds
therefore multiply to \(\Funiv(\gamma,p)\).  The compactness, possible
lost-mass, and witness arguments are given in
\cref{prop:flat-prior-converse} and Appendix~\ref{app:gaussian-converse}.

\subsection{LAN transfer}
\label{sec:local-asymptotic-normality}

The two-way LAN statement used in the cloning converses is stated once, for
every \(1\le r\le d\), in \cref{thm:fixed-rank-LAN}.  Its \(r=d\)
specialization is ordinary interior LAN.  We use throughout the
source-normalized convention
\[
  \exp\!\left(X_p(z)/\sqrt n\right)
  \quad\longleftrightarrow\quad
  D(z)\tau_pD(z)^*.
\]
Other physical rotation conventions are related to this one by invertible,
spectrum-dependent linear maps.  Such a change must be applied to both the
finite-dimensional family and the Gaussian displacement; we do not identify
coordinates belonging to two different conventions.  The fixed-rank theorem
is the result of Lahiry--Nussbaum, whose \(r=d\) case recovers interior LAN;
Kahn--Gu{\c t}{\u a} provides its faithful-state antecedent
\cite{KahnGuta2009LANFiniteDimensional,LahiryNussbaum2024}.

We will use the following continuity estimate here and in the appendices.  If
\(A,C,A',C'\) are positive trace-class operators with trace at most one, then
\begin{equation}\label{eq:fidelity-continuity}
  |F(A,C)-F(A',C')|
  \le \|A-A'\|_1^{1/2}+\|C-C'\|_1^{1/2}.
\end{equation}
Indeed, the Powers--St\o rmer inequality bounds
\(\|\sqrt A-\sqrt{A'}\|_2\) by \(\|A-A'\|_1^{1/2}\), and similarly for
\(C,C'\).  Adding and subtracting \(\sqrt{A'}\sqrt C\), the reverse triangle
inequality and Schatten H\"older give \eqref{eq:fidelity-continuity}, since
\(\|\sqrt C\|_2,\|\sqrt{A'}\|_2\le1\).

The following lemma formalizes this LAN fidelity transfer.

\begin{lemma}[LAN fidelity transfer]\label{lem:LAN-fidelity-transfer}
Let \(\{\rho_{n,\theta}\}_{\theta\in C}\) and
\(\{\sigma_{n,\theta}\}_{\theta\in C}\) be input and target models on a
compact parameter set \(C\).  Let \(\{\Phi_\theta\}\) and
\(\{\Psi_\theta\}\) be limiting models, and suppose channels \(S_n,T_n\)
satisfy
\[
  \sup_{\theta\in C}\|S_n(\Phi_\theta)-\rho_{n,\theta}\|_1\to0,
  \qquad
  \sup_{\theta\in C}\|T_n(\sigma_{n,\theta})-\Psi_\theta\|_1\to0.
\]
For every channel \(\cM_n\), put \(\Gamma_n=T_n\circ\cM_n\circ S_n\).
Then
\begin{equation}\label{eq:LAN-transfer}
  \inf_{\theta\in C}F(\Gamma_n(\Phi_\theta),\Psi_\theta)
  \ge
  \inf_{\theta\in C}F(\cM_n(\rho_{n,\theta}),\sigma_{n,\theta})-o(1).
\end{equation}
\end{lemma}

\begin{proof}
Let the two uniform approximation errors be \(\delta_n\) and \(\eta_n\).
Since \(T_n\circ\cM_n\) is a channel, uniformly in
\(\theta\),
\(\|\Gamma_n(\Phi_\theta)-T_n\cM_n(\rho_{n,\theta})\|_1
\le \|S_n(\Phi_\theta)-\rho_{n,\theta}\|_1\le\delta_n\).
Also
\(\|T_n(\sigma_{n,\theta})-\Psi_\theta\|_1\le\eta_n\).  Together with
\eqref{eq:fidelity-continuity}, these bounds give
\[
  F(\Gamma_n(\Phi_\theta),\Psi_\theta)
  \ge
  F(T_n\cM_n(\rho_{n,\theta}),T_n(\sigma_{n,\theta}))
  -\sqrt{\delta_n}-\sqrt{\eta_n}.
\]
Finally, monotonicity of fidelity under the channel \(T_n\) gives
\[
  F(T_n\cM_n(\rho_{n,\theta}),T_n(\sigma_{n,\theta}))
  \ge F(\cM_n(\rho_{n,\theta}),\sigma_{n,\theta}).
\]
Taking the infimum over \(\theta\) proves \eqref{eq:LAN-transfer}.
\end{proof}

\paragraph{Conditional sector limits versus LAN.}
Although the same thermal state \(\tau_p\) appears in both arguments, it
represents different limiting operations.  The conditional-sector proof
fixes the input and output Young labels and, on fixed-height cutoffs,
identifies the ideal target with \(\tau_p\) and the Cartan-sector output with
\(\tau_p^{(\gamma)}\); see
\eqref{eq:target-fixed-height-convergence} and
\eqref{eq:output-fixed-height-convergence}.  This comparison directly
explains the achievable orbital factor \(\Forb(\gamma,p)\).
The unified fixed-rank LAN experiment instead
retains the Young-label fluctuation as
  \(\mathsf N_{0,\Sigma_p}\), while orbital LAN marginalizes this classical
factor.  Thus the common occurrence of \(\tau_p\) reflects conditioning on
Young labels in the sector calculation and marginalizing their fluctuation
in orbital LAN.  The fixed-rank proof below combines these ingredients with
exact support compression and the vacuum-mode boundary argument.

\section{Fixed-rank cloning in ambient dimension}
\label{sec:fixed-rank-proof}
Fix \(1\le r\le d\) and \(p\in\Delta_r^\circ\).  We now prove the main
theorems directly on this rank-\(r\) stratum.  Exact compression separates
the positive-support calculation from the support--kernel contribution;
rank-aware label processing gives achievability; and one fixed-rank LAN
model, containing both thermal and vacuum modes, gives the converses.  The
case \(r=d\) is obtained by the same proof with no support--kernel modes,
rather than by a separate argument.
If \(d=r=1\), the model is a singleton and every fidelity below is one; we
therefore assume \(d\ge2\) in the nontrivial proof.
\subsection{Exact rank compression}\label{sec:compression}
We decorate representation spaces and maps by the ambient dimension when
needed.  For a partition \(\lambda\) with at most \(r\) rows, let
\[
  J_\lambda:\cH_\lambda^{(r)}\longrightarrow\cH_\lambda^{(d)}
\]
be the isometric \(\U(r)\)-intertwiner that maps the unit highest-weight
vector to the unit highest-weight vector.  Its range is the
\(\U(r)\)-submodule generated by that vector, and we set
\(P_\lambda=J_\lambda J_\lambda^*\).

For a compatible pair of \(r\)-row partitions \(\mu\vdash n\) and
\(\nu\vdash m\), write \(\omega=\nu-\mu\vdash m-n\).  In ambient
dimension \(s\in\{r,d\}\), the Cartan-sector channel is
\[
  \cC_{\mu,\nu}^{(s)}(X)
  =\frac{\dim\cH_\mu^{(s)}}{\dim\cH_\nu^{(s)}}
   (V_{\mu,\omega}^{(s)})^*
   (X\ot I_{\cH_\omega^{(s)}})
   V_{\mu,\omega}^{(s)}.
\]
For a probability vector \(x\in\R^s\), including one with zero entries, we
use the conditional Schur state
\[
  \rho_{x,\lambda}^{(s)}
  =\frac{\pi_\lambda^{(s)}(\diag x)}{s_\lambda^{(s)}(x)}
\]
whenever the denominator is nonzero.

\begin{lemma}[Exact rank-compression identity]
\label{lem:exact-compression}
For every \(p=(p_1,\ldots,p_r)\) with \(p_i>0\) and
\(\sum_i p_i=1\), and every \(r\)-row partition \(\lambda\),
\begin{equation}
  \rho_{\bar p,\lambda}^{(d)}
  =J_\lambda\rho_{p,\lambda}^{(r)}J_\lambda^*.
  \label{eq:conditional-support}
\end{equation}
For every compatible \(r\)-row pair \(\mu,\nu\) and every trace-class
operator \(X\) on \(\cH_\mu^{(r)}\),
\begin{equation}
  P_\nu\cC_{\mu,\nu}^{(d)}(J_\mu XJ_\mu^*)P_\nu
  =a_{\mu,\nu}^{(r,d)}
   J_\nu\cC_{\mu,\nu}^{(r)}(X)J_\nu^*,
  \label{eq:channel-compression}
\end{equation}
where
\begin{align}
  a_{\mu,\nu}^{(r,d)}
  &=\frac{
     \dim\cH_\mu^{(d)}/\dim\cH_\nu^{(d)}}{
     \dim\cH_\mu^{(r)}/\dim\cH_\nu^{(r)}}
  \label{eq:a-dimension-ratio}\\
  &=\prod_{\substack{1\le i\le r\\r<j\le d}}
    \frac{\mu_i+j-i}{\nu_i+j-i}.
  \label{eq:a-product}
\end{align}
Consequently, there is the exact finite-label fidelity factorization
\begin{align}
 &F\!\left(
   \cC_{\mu,\nu}^{(d)}(\rho_{\bar p,\mu}^{(d)}),
   \rho_{\bar p,\nu}^{(d)}
 \right)\notag\\
 &\qquad=\sqrt{a_{\mu,\nu}^{(r,d)}}\,
 F\!\left(
   \cC_{\mu,\nu}^{(r)}(\rho_{p,\mu}^{(r)}),
   \rho_{p,\nu}^{(r)}
 \right).
 \label{eq:exact-sector-factorization}
\end{align}
\end{lemma}

\begin{proof}
The singular specialization identity for Schur polynomials is
\[
  s_\lambda^{(d)}(p_1,\ldots,p_r,0,\ldots,0)
  =s_\lambda^{(r)}(p_1,\ldots,p_r),
\]
and it vanishes when \(\lambda\) has more than \(r\) rows.  Polynomial
functoriality for the standard inclusion
\(\iota:\C^r\hookrightarrow\C^d\) gives
\[
  \pi_\lambda^{(d)}
    \!\left(\iota\diag(p)\iota^*\right)
  =J_\lambda\pi_\lambda^{(r)}(\diag p)J_\lambda^*.
\]
Thus the ambient operator is supported on the highest-generated
\(\U(r)\)-submodule and restricts there to the \(r\)-dimensional polynomial
representation.  Normalizing by the Schur-polynomial identity proves
\eqref{eq:conditional-support}.

The two maps
\[
  V_{\mu,\omega}^{(d)}J_\nu,
  \qquad
  (J_\mu\ot J_\omega)V_{\mu,\omega}^{(r)}
\]
are \(\U(r)\)-intertwiners from \(\cH_\nu^{(r)}\) to
\(\cH_\mu^{(d)}\ot\cH_\omega^{(d)}\), and both send the unit
highest-weight vector to the same vector
\(v_\mu\ot v_\omega\).  Generation from the highest-weight vector gives
the exact restriction identity
\begin{equation}
  V_{\mu,\omega}^{(d)}J_\nu
  =(J_\mu\ot J_\omega)V_{\mu,\omega}^{(r)}.
  \label{eq:cartan-restriction}
\end{equation}
Inserting \eqref{eq:cartan-restriction} into the two sector-channel
formulas proves \eqref{eq:channel-compression}, with the scalar in
\eqref{eq:a-dimension-ratio}.  The Weyl dimension formula, and the fact
that all rows after \(r\) vanish, give \eqref{eq:a-product}; the
positive--positive Weyl factors cancel and the kernel--kernel factors are
identically one.

If \(A,B\ge0\) and \(PBP=B\), then
\[
  F(A,B)=\Tr\sqrt{\sqrt B A\sqrt B}
        =F(PAP,B).
\]
Apply this with \(P=P_\nu\), then use
\eqref{eq:conditional-support}, \eqref{eq:channel-compression}, invariance
under \(J_\nu\), and \(F(aA,B)=\sqrt a F(A,B)\).  This proves
\eqref{eq:exact-sector-factorization}.
Taking traces in \eqref{eq:channel-compression} also gives
\[
  \Tr\!\left[
    P_\nu\cC_{\mu,\nu}^{(d)}(J_\mu XJ_\mu^*)P_\nu
  \right]
  =a_{\mu,\nu}^{(r,d)}\Tr X.
\]
In particular, compatibility forces \(a_{\mu,\nu}^{(r,d)}\le1\).
\end{proof}

When \(r=d\), every \(J_\lambda\) is the identity and the product in
\eqref{eq:a-product} is empty, so
\(a_{\mu,\nu}^{(d,d)}=1\).  Thus exact compression is part of the common
proof and becomes tautological when \(r=d\).

\begin{corollary}[Uniform ambient fixed-rank sector limit]
\label{cor:rank-sector}
Let \(K\Subset\Delta_r^\circ\), let \(\delta_n\to0\), and let
\(\mu\vdash n\), \(\nu\vdash m_n\) be compatible \(r\)-row labels
satisfying
\[
  \left\|\frac{(\mu_1,\ldots,\mu_r)}n-p\right\|_\infty\le\delta_n,
  \qquad
  \left\|\frac{(\nu_1,\ldots,\nu_r)}{m_n}-p\right\|_\infty\le\delta_n,
  \qquad p\in K.
\]
Then
\begin{align*}
 &F\!\left(
   \cC_{\mu,\nu}^{(d)}(\rho_{\bar p,\mu}^{(d)}),
   \rho_{\bar p,\nu}^{(d)}
 \right)\\
 &\qquad=\prod_{(i,j)\in\mathcal I_{r,d}}f_\gamma(q_{ij})+o(1)
 =\Forb^{(r,d)}(\gamma,p)+o(1),
\end{align*}
uniformly over the displayed family.
\end{corollary}

\begin{proof}
Every factor in \eqref{eq:a-product} tends uniformly to \(1/\gamma\), so
\[
  a_{\mu,\nu}^{(r,d)}=\gamma^{-r(d-r)}+o(1).
\]
The remaining fidelity in \eqref{eq:exact-sector-factorization} is the
intrinsic positive-support sector calculation in dimension \(r\).  If
\(r\ge2\),
\cref{prop:sector-fidelity}, applied with effective dimension \(r\), gives the
product over \(1\le i<j\le r\).  If \(r=1\), all three
\(\U(1)\)-modules are one-dimensional and the lower-dimensional sector
fidelity is exactly one.  In this case the result reduces to
\(\gamma^{-(d-1)/2}\), the proportional limit of Werner's pure-state
cloning root fidelity.
\end{proof}

This proof also makes the stabilizer bookkeeping transparent.  The
positive--positive roots are inherited from \(\U(r)\); the support--kernel
roots occur only in the relative Weyl-dimension ratio; and kernel--kernel
roots cancel exactly.

\subsection{Rank-aware label processing and achievability}
\label{sec:rank-aware-achievability}

For every \(N\), the ambient \(d\)-dimensional Young law of a rank-\(r\)
state reduces exactly to the \(r\)-dimensional law:
\begin{equation}
  P_{N,\bar p}^{(d)}(\lambda)
  =\begin{cases}
    P_{N,p}^{(r)}(\lambda),&\ell(\lambda)\le r,\\
    0,&\ell(\lambda)>r.
  \end{cases}
  \label{eq:young-rank-reduction}
\end{equation}
Indeed, this is the same Schur-polynomial specialization used in
\cref{lem:exact-compression}; the Specht dimension is independent of the
ambient unitary dimension.

\begin{proposition}[Ambient lifted-kernel bound at fixed rank]
\label{prop:fixed-rank-lift}
Fix \(p\in\Delta_r^\circ\), set
\[
  \mathsf Y_N^{(r)}
  =\{\lambda\vdash N:\ell(\lambda)\le r\},
\]
and let \(K_n(\nu\mid\mu)\) be a Markov kernel from
\(\mathsf Y_n^{(r)}\) to \(\mathsf Y_{m_n}^{(r)}\), with labels padded by
zeros in ambient dimension \(d\).  Extend it arbitrarily to the remaining
ambient input labels; for example, when \(\ell(\mu)>r\), send the label
deterministically to \((m_n,0,\ldots,0)\).  Set
\[
  \Pi_{n,p}^{(r)}(\mu,\nu)
  =P_{n,p}^{(r)}(\mu)K_n(\nu\mid\mu),
  \qquad
  Q_{m_n,p}^{(r)}(\nu)=\sum_\mu\Pi_{n,p}^{(r)}(\mu,\nu).
\]
Suppose \(G\subset\mathsf Y_n^{(r)}\times\mathsf Y_{m_n}^{(r)}\),
\(b\in[0,1]\), and every pair in \(G\) is compatible and has ambient
sector fidelity at least \(b\).  On the other pairs, define the lift with
any covariant channel into the prescribed \(\nu\)-block, such as the
maximally mixed replacement channel.  Then the ambient Schur--Cartan lift
\(\mathcal Q_n^{(d)}[K_n]\) satisfies, for every \(g\in\U(d)\),
\[
  F\!\left(
    \mathcal Q_n^{(d)}[K_n](\rho_{\bar p,g}^{\ot n}),
    \rho_{\bar p,g}^{\ot m_n}
  \right)
  \ge b\left[
    F(Q_{m_n,p}^{(r)},P_{m_n,p}^{(r)})
    -\sqrt{\Pi_{n,p}^{(r)}(G^c)}
  \right].
\]
\end{proposition}

\begin{proof}
The proof of \cref{prop:lifted-kernel-bound} uses only
the direct-sum fidelity identity, covariance, joint concavity, and
Cauchy--Schwarz.  By \eqref{eq:young-rank-reduction}, all blocks with more
than \(r\) rows have zero input and target weight.  Omitting those blocks
therefore gives the displayed fixed-rank statement verbatim.
\end{proof}

\subsubsection{Known spectrum}

Sample \(\nu\) independently from the exact law
\(P_{m_n,p}^{(r)}\), pad both measured \(\mu\) and sampled \(\nu\) by zeros,
and use the ambient \(d\)-dimensional Cartan-sector channel whenever
\(\nu-\mu\) is a partition.  Otherwise, use the covariant replacement
channel
\[
  X\longmapsto \Tr(X)\frac{I_{\cH_\nu^{(d)}}}
                              {\dim\cH_\nu^{(d)}}
\]
into the same sampled \(\nu\)-block.  Thus the fallback does not alter the
sampled output label.  On input labels with more than \(r\) rows, define
the kernel by the deterministic extension in \cref{prop:fixed-rank-lift};
these branches have zero probability for the states under consideration.
Because \(p_r\) and all adjacent positive gaps are bounded below on compact
subsets of \(\Delta_r^\circ\), two independent typical labels obey
\[
  \nu-\mu\vdash m_n-n
\]
with probability \(1-o(1)\).  The output-label law is exactly the target
law, so its Hellinger factor is one.  Combining Young concentration,
\cref{cor:rank-sector}, and \cref{prop:fixed-rank-lift} proves the lower
bound in \cref{thm:known-optimum}.

\subsubsection{Unknown spectrum of known rank}

Define the randomized label dilation directly on the \(r\) active rows.  For
\(i<r\), draw independent
\(U_i\sim\operatorname{Unif}[-1/2,1/2)\) and set
\[
  \widetilde\nu_i
  =\left\lfloor\gamma_n(\mu_i+U_i)+\frac12\right\rfloor,
  \qquad
  \widetilde\nu_r=m_n-\sum_{i<r}\widetilde\nu_i.
\]
Pad by \(d-r\) zeros and use the compatibility check and fallback
replacement of \cref{sec:unknown-label-kernel}.  When \(r=1\), simply set
\(\widetilde\nu=(m_n)\).  On measured labels with more than \(r\) rows,
choose the concrete fallback output \((m_n,0,\ldots,0)\); these branches
have zero probability on rank-\(r\) inputs.
Call this padded kernel \(K_n^{\mathrm{rnd},(r)}\), and denote the resulting
ambient lifted channel by
\[
  \cC_n^{(r)}:=\mathcal Q_n^{(d)}[K_n^{\mathrm{rnd},(r)}].
\]
It depends on \(r,d,n,m_n\), but not on the unknown positive spectrum
\(p\).

The dimension-\(r\) Young local limit gives, uniformly for
\(p\in K\Subset\Delta_r^\circ\),
\[
  F(Q_{m_n,p}^{(r)},P_{m_n,p}^{(r)})
  =\left(\frac{2\sqrt\gamma}{1+\gamma}\right)^{(r-1)/2}+o(1).
\]
For \(r=1\), no local-limit theorem is needed: both laws are the point mass
at the one-row label, so the fidelity
is exactly one.
Typical rounded pairs are compatible, and \cref{cor:rank-sector} supplies
the common sector factor.  \Cref{prop:fixed-rank-lift} therefore gives
\eqref{eq:rank-univ-value} and proves the achievability assertion in
\cref{thm:unknown-optimum}.

\subsection{Fixed-rank LAN and the converses}
\label{sec:fixed-rank-LAN}

The same local experiment covers every \(1\le r\le d\).  For \(r=d\) it
reduces to ordinary interior LAN.  For \(r<d\), its support--kernel
coordinates are genuine vacuum boundary modes supplied directly by
fixed-rank LAN, rather than by a singular limit from faithful states
\cite{LahiryNussbaum2024}.

Put
\[
  \mathsf H_r=\left\{h\in\R^r:\sum_{i=1}^r h_i=0\right\},
  \qquad
  \Sigma_p^{(r)}
  =\left.(\diag(p)-pp^{\mathsf T})\right|_{\mathsf H_r}.
\]
Recall the active root set \(\mathcal I_{r,d}\) and the convention
\(p_j=0\) for \(j>r\) from \eqref{eq:active-root-set-main}.  Use the
anti-Hermitian local generator
\begin{equation}
  X_p(z)=\sum_{(i,j)\in\mathcal I_{r,d}}
  \frac{z_{ij}E_{ji}-\bar z_{ij}E_{ij}}{\sqrt{p_i-p_j}}.
  \label{eq:low-rank-generator}
\end{equation}
For fixed \((h,z)\in\mathsf H_r\times\C^{\mathcal I_{r,d}}\), define the
single-copy local state
\begin{equation}
  \rho_{n,h,z}
  =e^{X_p(z)/\sqrt n}
   \diag\!\left(p+\frac h{\sqrt n},0_{d-r}\right)
   e^{-X_p(z)/\sqrt n}.
  \label{eq:low-rank-local-state}
\end{equation}
It is a rank-\(r\) state for every sufficiently large \(n\), uniformly when
\((h,z)\) ranges over a fixed compact set.

Let
\[
  q_{ij}=\begin{cases}p_j/p_i,&j\le r,\\0,&j>r,\end{cases}
  \qquad
  \cF_{r,d}=\bigotimes_{(i,j)\in\mathcal I_{r,d}}\cF_{ij},
\]
and put
\[
  \tau_q=(1-q)\sum_{k\ge0}q^k\proj{k},
  \qquad
  \tau_p^{(r,d)}=\bigotimes_{(i,j)\in\mathcal I_{r,d}}\tau_{q_{ij}},
  \qquad
  D(z)=\bigotimes_{(i,j)\in\mathcal I_{r,d}}D(z_{ij}).
\]
Thus \(\tau_0=\proj0\).  The input and ideal target limit experiments are
\begin{align}
  \Phi_{h,z}^{(r,d)}
  &=\mathsf N_{h,\Sigma_p^{(r)}}\ot
    D(z)\tau_p^{(r,d)}D(z)^*,
  \label{eq:low-rank-LAN-model}\\
  \Psi_{h,z}^{(r,d,\gamma)}
  &=\mathsf N_{\sqrt\gamma h,\Sigma_p^{(r)}}\ot
    D(\sqrt\gamma z)\tau_p^{(r,d)}D(\sqrt\gamma z)^*.
  \label{eq:low-rank-LAN-target}
\end{align}
When \(r<d\), the support--kernel coordinates are shifted pure Gaussian
states, not thermal states with a small positive temperature.  When \(r=d\),
that collection is empty.

\begin{theorem}[Two-way LAN at fixed rank]
\label{thm:fixed-rank-LAN}
Fix \(d,r\), and \(p_1>\cdots>p_r>0\).  There are channels
\begin{align*}
  T_N^{(r)} &: \cT_1((\C^d)^{\ot N})
    \longrightarrow L^1(\mathsf H_r;\cT_1(\cF_{r,d})),\\
  S_N^{(r)} &: L^1(\mathsf H_r;\cT_1(\cF_{r,d}))
    \longrightarrow \cT_1((\C^d)^{\ot N})
\end{align*}
such that, for every compact
\(C\subset\mathsf H_r\times\C^{\mathcal I_{r,d}}\),
\begin{align}
  \sup_{(h,z)\in C}
  \left\|T_n^{(r)}(\rho_{n,h,z}^{\ot n})
             -\Phi_{h,z}^{(r,d)}\right\|_1&\longrightarrow0,
  \label{eq:low-rank-LAN-forward}\\
  \sup_{(h,z)\in C}
  \left\|S_n^{(r)}(\Phi_{h,z}^{(r,d)})
             -\rho_{n,h,z}^{\ot n}\right\|_1&\longrightarrow0,
  \label{eq:low-rank-LAN-reverse}\\
  \sup_{(h,z)\in C}
  \left\|T_{m_n}^{(r)}(\rho_{n,h,z}^{\ot m_n})
             -\Psi_{h,z}^{(r,d,\gamma)}\right\|_1&\longrightarrow0.
  \label{eq:low-rank-LAN-target-sample}
\end{align}
\end{theorem}

\begin{proof}
For every \(1\le r\le d\), the first two assertions are the two-way
fixed-rank LAN theorem of Lahiry and Nussbaum
\cite[Theorem~3.3]{LahiryNussbaum2024}, after the coordinate identification
in \cref{rem:fixed-rank-coordinate}.  They explicitly note that \(r=d\)
recovers ordinary interior LAN.  Their result is pointwise in the fixed base
spectrum \(p\), which is exactly what the converse below needs.  In fact,
they prove uniformity on expanding boxes: the spectral coordinate may grow
as \(n^a\), \(a<1/4\), and the rotation coordinate as \(n^\beta\),
\(\beta<1/9\) (their symbol for \(a\) is \(\gamma\), renamed here to avoid
conflict with the cloning ratio).  We use only the weaker compact-set
consequence.  At \(r=d\), this is ordinary interior LAN in the same
source-normalized coordinate convention.

In both source normalizations, the spectral coordinate
\(u=(u_1,\ldots,u_{r-1})\in\R^{r-1}\) is identified with
\[
  h=\iota(u)
  :=\left(u_1,\ldots,u_{r-1},-\sum_{i<r}u_i\right)\in\mathsf H_r.
\]
If \(V_p=(\delta_{ij}p_i-p_ip_j)_{i,j<r}\) is their multinomial
covariance and \(A\) is the matrix of \(\iota\), then
\[
  AV_pA^{\mathsf T}=\diag(p)-pp^{\mathsf T}.
\]
Consequently, their classical normal experiment pushes forward to
\(\mathsf N_{h,\Sigma_p^{(r)}}\), as used above.

For the third assertion, observe the exact reparametrization
\[
  \rho_{n,h,z}
  =\rho_{m_n,\sqrt{\gamma_n}h,\sqrt{\gamma_n}z},
  \qquad \gamma_n=m_n/n.
\]
Apply the forward LAN channel at sample size \(m_n\).  A fixed compact set
contains all the scaled parameters for large \(n\), and
\(\gamma_n\to\gamma\).  Translation of the Gaussian density and Weyl
displacement are uniformly trace-norm continuous on that compact set, which
gives \eqref{eq:low-rank-LAN-target-sample}.
\end{proof}

\begin{remark}[Coordinate normalization]
\label{rem:fixed-rank-coordinate}
The normalization in \eqref{eq:low-rank-generator} is chosen so that the
finite-dimensional parameter \(z\) is also the oscillator displacement in
\eqref{eq:low-rank-LAN-model}.  This is the source-normalized convention of
the cited fixed-rank LAN theorem and the convention introduced in
\cref{sec:local-coordinates}.  Other physical rotation coordinates are
related by invertible, spectrum-dependent linear maps, but the Gaussian
displacement must then be reparametrized by the same map.  We use
\(X_p(z)\leftrightarrow D(z)\) throughout, including when \(r=d\).
\end{remark}

\subsubsection{Gaussian amplification with thermal and vacuum modes}

For \(q\in[0,1)\), the quantum-limited amplifier obeys
\[
  \cA_\gamma(D(z)\tau_qD(z)^*)
  =D(\sqrt\gamma z)\tau_{q^{(\gamma)}}D(\sqrt\gamma z)^*,
  \qquad
  q^{(\gamma)}=1-\frac{1-q}{\gamma}.
\]
Its fidelity against the ideal target is
\[
  F(\tau_{q^{(\gamma)}},\tau_q)=f_\gamma(q).
\]
At the boundary this reads
\[
  q=0,
  \qquad q^{(\gamma)}=\frac{\gamma-1}{\gamma},
  \qquad F(\tau_{(\gamma-1)/\gamma},\proj0)=\gamma^{-1/2}.
\]

The product amplifier therefore gives achievability uniformly across the
thermal and vacuum modes.  The converse requires one explicit boundary
step because the vacuum target is not faithful.  If \(\widehat N\) is the
number operator, the positive-temperature one-mode witness is
\[
  W_q=
  \sqrt{\frac{1-q}{1-q^{(\gamma)}}}
  \left(\frac{q}{q^{(\gamma)}}\right)^{\widehat N/2},
  \qquad q>0.
\]
Retain this injective witness for every thermal mode and use
\[
  W_0=\sqrt\gamma\,\proj0
\]
for every coherent mode.  Moreover,
\begin{equation}
  \|W_q-W_0\|
  =\sqrt\gamma
   \sqrt{\frac{\gamma q}{\gamma-1+q}}
  \longrightarrow0.
  \label{eq:witness-boundary-continuity}
\end{equation}

Let \(P_{\mathrm{vac}}\) be the identity on the thermal modes tensored with
the joint-vacuum projection on all coherent modes, and take the tensor
product \(W\) of the displayed one-mode witnesses.  Write
\[
  \tau=\tau_p^{(r,d)},
  \qquad
  \tau^{(\gamma)}
  =\bigotimes_{(i,j)\in\mathcal I_{r,d}}\tau_{q_{ij}^{(\gamma)}},
  \qquad
  B=\Forb^{(r,d)}(\gamma,p).
\]
The operator \(W\) is compact: every positive-temperature factor is compact
and every zero-temperature factor is rank one.
The ideal zero-displacement target \(\tau\) is supported on
\(P_{\mathrm{vac}}\), so, for every positive output \(R\),
\[
  F(R,\tau)=F(P_{\mathrm{vac}}RP_{\mathrm{vac}},\tau).
\]
On that support \(W\) is injective.  Modewise,
\[
  \Tr(\tau_{(\gamma-1)/\gamma}W_0)
  =\Tr(\proj0(W_0|_{\C\ket0})^{-1})
  =\gamma^{-1/2}.
\]
Together with the corresponding \(q>0\) identities, this gives
\begin{equation}
  \Tr(\tau^{(\gamma)}W)
  =\Tr\!\left(\tau
    (W|_{P_{\mathrm{vac}}\cF_{r,d}})^{-1}\right)
  =B.
  \label{eq:mixed-witness-identities}
\end{equation}
The weighted-fidelity lemma from
Appendix~\ref{app:gaussian-converse}, applied after support compression,
therefore yields
\begin{equation}
  F(R,\tau)^2
  \le B\,\Tr(RW).
  \label{eq:mixed-weighted-fidelity}
\end{equation}

The least-noise moment inequality also extends to the vacuum modes.  Indeed,
replace every zero temperature by \(q>0\), apply the positive-temperature
inequality, and let \(q\downarrow0\).  Here
\(\tau_q\to\proj0\) and
\(\tau_{q^{(\gamma)}}\to\tau_{(\gamma-1)/\gamma}\) in trace norm, while
\eqref{eq:witness-boundary-continuity} gives operator-norm convergence of
the witnesses.  Hence every completely positive, trace-nonincreasing,
gain-\(\sqrt\gamma\) displacement-covariant map \(\Gamma\) satisfies
\begin{equation}
  \Tr[\Gamma(\tau)W]
  \le \Tr[\Gamma(\tau)]\Tr(\tau^{(\gamma)}W)
  \le B.
  \label{eq:mixed-least-noise}
\end{equation}
Combining \eqref{eq:mixed-weighted-fidelity} and
\eqref{eq:mixed-least-noise} gives
\(F(\Gamma(\tau),\tau)\le B\), including when a compact-limit map loses
trace.

The corresponding fiberwise statement will be used for the hybrid model.
If \(R(y)\) is a positive trace-class-valued density, \(g\) is a probability
density on \(\mathsf H_r\), and \(\chi>0\), then compression followed by
the hybrid weighted-fidelity lemma gives
\begin{equation}
  F(R,g\tau)^2
  \le
  \left(\int_{\mathsf H_r}
    \chi(y)\Tr[R(y)W]\,\dd y\right)
  \left(\int_{\mathsf H_r}\frac{g(y)}{\chi(y)}\,\dd y\right)B.
  \label{eq:mixed-hybrid-weighted-fidelity}
\end{equation}

For \(L>0\), let
\begin{align*}
  \mathsf Z_L^{(r,d)}
  &=\{z:\ |\Re z_{ij}|,|\Im z_{ij}|\le L
          \text{ for every }(i,j)\in\mathcal I_{r,d}\},\\
  \mathsf H_{r,L}
  &=\{h\in\mathsf H_r:\ \|h\|_\infty\le L\},
  \qquad
  \Omega_L^{(r,d)}=\mathsf H_{r,L}\times\mathsf Z_L^{(r,d)}.
\end{align*}
Here \(\mathsf H_r\) carries its induced Lebesgue measure, complex
coordinates carry their standard real Lebesgue measure, and the
zero-dimensional volume is understood to be one when \(r=1\).

\begin{proposition}[Mixed thermal--coherent flat-prior optima]
\label{prop:mixed-gaussian}
With \(\phi_z=D(z)\tau D(z)^*\) and
\(\psi_z=D(\sqrt\gamma z)\tau D(\sqrt\gamma z)^*\),
\begin{align}
 &\lim_{L\to\infty}\sup_\Lambda
  \frac1{|\mathsf Z_L^{(r,d)}|}
  \int_{\mathsf Z_L^{(r,d)}}
  F\!\left(\Lambda(\phi_z),\psi_z\right)\,\dd z
  =B,
  \label{eq:mixed-orbital-average}\\
 &\lim_{L\to\infty}\sup_\Lambda
  \frac1{|\Omega_L^{(r,d)}|}
  \int_{\Omega_L^{(r,d)}}
  F\!\left(\Lambda(\Phi_{h,z}^{(r,d)}),
                 \Psi_{h,z}^{(r,d,\gamma)}\right)\,\dd h\,\dd z
  =c_\gamma^{(r-1)/2}B.
  \label{eq:mixed-full-average}
\end{align}
No Gaussianity, covariance, or product assumption is imposed on the
competing channels.  In \eqref{eq:mixed-orbital-average}, the supremum is
over quantum channels on \(\cT_1(\cF_{r,d})\).  In
\eqref{eq:mixed-full-average}, it is over classical--quantum channels on
\(L^1(\mathsf H_r;\cT_1(\cF_{r,d}))\).
\end{proposition}

\begin{proof}
For achievability, tensor the quantum-limited amplifiers over all active
modes.  On a coherent mode this is the usual amplifier with vacuum input;
its added thermal noise gives \(f_\gamma(0)=\gamma^{-1/2}\).  For the hybrid
model, tensor this channel with the classical pushforward
\(x\mapsto\sqrt\gamma x\).  Multiplicativity gives the two displayed lower
bounds.

For the converse, the F{\o}lner averaging and compact-map argument of
Appendix~\ref{app:gaussian-converse} produces a gain-\(\sqrt\gamma\)
displacement-covariant,
trace-nonincreasing limiting map.  Equations
\eqref{eq:mixed-weighted-fidelity}--\eqref{eq:mixed-least-noise} give the
orbital bound.  In the hybrid argument, compress each output fiber
\(R(y)\mapsto P_{\mathrm{vac}}R(y)P_{\mathrm{vac}}\) before applying the
hybrid weighted-fidelity lemma.  This does not change fidelity with the
target, and it makes the restricted witness injective.  The scalar
convolution part of the proof is otherwise unchanged and contributes
\(c_\gamma^{r-1}\) at the squared-fidelity level, hence
\(c_\gamma^{(r-1)/2}\) at root fidelity.  This proves
\eqref{eq:mixed-full-average}.
\end{proof}

We now spell out the finite-dimensional transfer.  Define
\begin{equation}
  u_n^{(r)}(p)=\sup_{\cM_n}\inf_{g\in\U(d)}
  F\!\left(\cM_n(\rho_{\bar p,g}^{\ot n}),
             \rho_{\bar p,g}^{\ot m_n}\right).
  \label{eq:rank-known-minimax}
\end{equation}
For a fixed \(L\), restrict the orbit to the states
\(\rho_{n,0,z}\), \(z\in\mathsf Z_L^{(r,d)}\).  Discard the
parameter-independent classical factor in \(T_N^{(r)}\) and reinsert it in
  \(S_N^{(r)}\).  The general LAN fidelity-transfer lemma,
\cref{lem:LAN-fidelity-transfer},
turns every nearly optimal finite cloner into a Gaussian channel
\(\Gamma_n\) satisfying
\[
  u_n^{(r)}(p)
  \le\inf_{z\in\mathsf Z_L^{(r,d)}}
       F(\Gamma_n(\phi_z),\psi_z)+o(1).
\]
The infimum is no larger than the flat-prior average.  First take
\(\limsup_n\), then the supremum over Gaussian-model channels, and finally
\(L\to\infty\).  Equation \eqref{eq:mixed-orbital-average} gives
\(\limsup_n u_n^{(r)}(p)\le B\).  Together with the finite-dimensional lower
bound, this proves \cref{thm:known-optimum}.

For all sufficiently large \(n\), define the local fixed-rank minimax value
\begin{equation}
  v_{n,L}^{(r)}(p)
  =\sup_{\cM_n}\inf_{(h,z)\in\Omega_L^{(r,d)}}
  F\!\left(\cM_n(\rho_{n,h,z}^{\ot n}),
             \rho_{n,h,z}^{\ot m_n}\right).
  \label{eq:rank-local-minimax}
\end{equation}

\begin{proposition}[Pointwise fixed-rank local minimax limit]
\label{prop:rank-local-minimax}
For every \(p\in\Delta_r^\circ\),
\[
  \lim_{L\to\infty}\limsup_{n\to\infty}v_{n,L}^{(r)}(p)
  =\lim_{L\to\infty}\liminf_{n\to\infty}v_{n,L}^{(r)}(p)
  =\Funiv^{(r,d)}(\gamma,p).
\]
\end{proposition}

\begin{proof}
For fixed \(L\), all spectra \(p+h/\sqrt n\) lie in one compact subset of
  \(\Delta_r^\circ\) for large \(n\).  The fixed-rank channel constructed
  above and continuity of \eqref{eq:rank-univ-value} therefore give the
lower bound uniformly on the local family.

For the converse, choose a channel within \(1/n\) of
\(v_{n,L}^{(r)}(p)\) and sandwich it between \(S_n^{(r)}\) and
\(T_{m_n}^{(r)}\).  The three approximations in
\cref{thm:fixed-rank-LAN} and fidelity transfer give
\[
  v_{n,L}^{(r)}(p)
  \le\inf_{(h,z)\in\Omega_L^{(r,d)}}
  F\!\left(\Gamma_n(\Phi_{h,z}^{(r,d)}),
             \Psi_{h,z}^{(r,d,\gamma)}\right)+o(1).
\]
Bound the infimum by the average and use
\eqref{eq:mixed-full-average}, first letting \(n\to\infty\) and then
\(L\to\infty\).  This yields the upper bound.  The outer limits exist
because both the limsup and liminf are nonincreasing in \(L\).
\end{proof}

Finally, for a compact set \(K\) as in \cref{thm:unknown-optimum}, write
\[
  w_n^{(r)}(K)=\sup_{\cM_n}
  \inf_{\substack{p\in K\\g\in\U(d)}}
  F\!\left(\cM_n(\rho_{\bar p,g}^{\ot n}),
             \rho_{\bar p,g}^{\ot m_n}\right).
\]
If \(p_0\) lies in the interior of \(K\) relative to the affine hyperplane
\(\sum_i p_i=1\), every fixed local spectral box around \(p_0\) lies in
\(K\) for all sufficiently large \(n\).  Hence
\(w_n^{(r)}(K)\le v_{n,L}^{(r)}(p_0)\).  Let \(n\to\infty\), then
\(L\to\infty\), and finally take the infimum over such \(p_0\).  The
assumption that \(K\) is the closure of this relative interior and the
continuity of \(\Funiv^{(r,d)}\) give
\[
  \limsup_n w_n^{(r)}(K)
  \le\inf_{p\in K}\Funiv^{(r,d)}(\gamma,p).
\]
The uniform rank-aware construction above supplies the reverse inequality
and completes \cref{thm:unknown-optimum}.

\section{The Grassmannian projector orbit}\label{sec:grassmann}
The flat rank-\(r\) state is \(\rho_P=P/r\), with \(P\in\Gr(r,d)\).  For
\(r>1\), its positive eigenvalue has multiplicity \(r\), so it is not covered
by the simple-positive-spectrum theorem.  For \(r=1\), it is the pure-state
orbit already covered by \cref{thm:known-optimum}; the direct calculation
below provides an independent derivation.  We prove \cref{thm:grassmann}
without invoking degenerate LAN.
\subsection{Exact sector calculation}

Let
\begin{equation}
  D_s^\lambda=\dim\cH_\lambda^{(s)},
  \qquad
  R_\lambda=\frac{D_d^\lambda}{D_r^\lambda}
  =\prod_{\substack{1\le i\le r\\r<j\le d}}
    \frac{\lambda_i+j-i}{j-i}.
  \label{eq:R-lambda}
\end{equation}
For the standard projector \(P_0\) onto the first \(r\) coordinates, the
conditional state in the \(\lambda\)-sector is
\[
  \rho_{P_0,\lambda}
  =J_\lambda\frac{I_{\cH_\lambda^{(r)}}}{D_r^\lambda}J_\lambda^*.
\]
The \(r\)-dimensional Cartan channel is \(\U(r)\)-covariant.  It therefore
maps the maximally mixed state on the irreducible input module exactly to
the maximally mixed state on the irreducible output module.  The compression
identity \eqref{eq:exact-sector-factorization} consequently becomes, for
every compatible pair,
\begin{equation}
  F\!\left(
    \cC_{\mu,\nu}^{(d)}(\rho_{P_0,\mu}),
    \rho_{P_0,\nu}
  \right)
  =\sqrt{\frac{R_\mu}{R_\nu}}.
  \label{eq:grassmann-exact-sector}
\end{equation}
For \(\mu_i/n\to1/r\) and \(\nu_i/m_n\to1/r\), this tends to the
right-hand side of \eqref{eq:grassmann-value}.

\subsection{A compatible exact-marginal coupling of flat Schur laws}

Let
\begin{equation}
  \mathsf P_N^{(r)}(\lambda)
  =\frac{\dim\cS_\lambda\,D_r^\lambda}{r^N},
  \qquad \lambda\vdash N,\quad\ell(\lambda)\le r,
  \label{eq:flat-schur-law}
\end{equation}
be the Schur--Weyl law of the maximally mixed \(r\)-level state.  The
complete Schur decomposition is
\begin{equation}
  (P/r)^{\ot N}
  =\bigoplus_{\ell(\lambda)\le r}
   \mathsf P_N^{(r)}(\lambda)
   \frac{\Pi_{\lambda,P}}{D_r^\lambda}
   \ot\frac{I_{\cS_\lambda}}{\dim\cS_\lambda}.
  \label{eq:flat-schur-decomposition}
\end{equation}
Here \(\Pi_{\lambda,P}\) is the rank-\(D_r^\lambda\) support projection
defined intrinsically by the subspace \(P\subset\C^d\).
The
Schur--Weyl/random-word central limit theorem
\cite{Kuperberg2002,Johansson2001} states that
\begin{equation}
  X_N:=\frac{\lambda-(N/r)\mathbf1_r}{\sqrt N}
  \Longrightarrow X,
  \label{eq:traceless-GUE-limit}
\end{equation}
where \(X\) is a fixed scaling of the ordered eigenvalue vector of a
traceless \(r\times r\) GUE matrix.  In particular,
\[
  \Pr(X_1>\cdots>X_r)=1.
\]

Both \(X_n\) and \(X_{m_n}\) converge to the same law.  Their Prokhorov
distance therefore tends to zero, so Strassen's coupling theorem supplies a
coupling \(\Pi_n(\mu,\nu)\) of the \emph{exact} marginals
\(\mathsf P_n^{(r)}\) and \(\mathsf P_{m_n}^{(r)}\) such that
\begin{equation}
  \left\|
    \frac{\mu-(n/r)\mathbf1_r}{\sqrt n}
    -\frac{\nu-(m_n/r)\mathbf1_r}{\sqrt{m_n}}
  \right\|_\infty
  \longrightarrow0
  \quad\text{in probability}.
  \label{eq:shape-coupling}
\end{equation}
Disintegrate this finite coupling to obtain a label kernel
\(K_n(\nu\mid\mu)\).

Put \(\omega=\nu-\mu\).  Tightness in
\eqref{eq:shape-coupling} gives \(\omega_i>0\) for every \(i\) with
probability \(1-o(1)\), because its deterministic leading term is
\((m_n-n)/r\).  For adjacent rows,
\begin{align*}
  \frac{\omega_i-\omega_{i+1}}{\sqrt n}
  &={\sqrt{\gamma_n}}
    \left[
      \frac{\nu_i-\nu_{i+1}}{\sqrt{m_n}}
    \right]
    -\left[
      \frac{\mu_i-\mu_{i+1}}{\sqrt n}
    \right]\\
  &\Longrightarrow
    (\sqrt\gamma-1)(X_i-X_{i+1})>0.
\end{align*}
Thus \(\omega\) is a partition with probability \(1-o(1)\).  Use the
ambient Cartan-sector channel on this good event and, otherwise, the
covariant replacement channel into the same sampled \(\nu\)-block.  The
output-label marginal is exactly the target law, and
\eqref{eq:grassmann-exact-sector} holds uniformly on the typical part.
If \(G_n\) is the intersection of compatibility and the two typical-label
events, the direct-sum fidelity formula and joint concavity give explicitly
\[
  F\!\left(
    \mathcal Q_n^{(d)}[K_n]((P/r)^{\ot n}),
    (P/r)^{\ot m_n}
  \right)
  \ge
  \sum_{(\mu,\nu)\in G_n}
  \Pi_n(\mu,\nu)\sqrt{\frac{R_\mu}{R_\nu}}
  \longrightarrow\gamma^{-r(d-r)/2}.
\]
Indeed, \(\Pi_n(G_n^c)\to0\), while the summand converges uniformly on the
typical part.  On every compatible pair, \(\nu_i\ge\mu_i\) and hence
\[
  0\le\sqrt{\frac{R_\mu}{R_\nu}}\le1.
\]
Thus convergence in probability under the coupling also implies convergence
of the displayed expectations (equivalently, the summands are uniformly
integrable).  This proves achievability in \cref{thm:grassmann}.

\subsection{Direct finite-dimensional converse}

We give the argument because it avoids any appeal to degenerate-spectrum
LAN.  Twirling a cloner over \(\U(d)\), and over input and output
permutations, cannot decrease its worst-case Grassmannian fidelity.  After
Schur pinching, let
\[
  \Phi_{\nu\mu}:
  \cT_1(\cH_\mu^{(d)})\longrightarrow
  \cT_1(\cH_\nu^{(d)})
\]
be its \(\U(d)\)-equivariant completely positive subchannels.  Schur's
lemma gives
\[
  \Phi_{\nu\mu}^*(I_{\cH_\nu^{(d)}})
  =K(\nu\mid\mu)I_{\cH_\mu^{(d)}},
  \qquad
  \sum_\nu K(\nu\mid\mu)=1.
\]
Thus \(K\) is a state-independent label transition.  For every branch with
\(K(\nu\mid\mu)>0\), normalize it to the covariant channel
\[
  \cE_{\mu,\nu}
  =\frac{\Phi_{\nu\mu}}{K(\nu\mid\mu)}.
\]
Zero-probability branches may be defined arbitrarily.

Let \(\Pi_{\lambda,P}\) be the rank-\(D_r^\lambda\) support projection of
the conditional \(\lambda\)-sector state.  Covariance and irreducibility
give
\[
  \cE_{\mu,\nu}(I_{\cH_\mu^{(d)}})
  =\frac{D_d^\mu}{D_d^\nu}I_{\cH_\nu^{(d)}}.
\]
The unnormalized \(\nu\)-block of the actual output is therefore
\[
  Y_{\nu,P}
  =\sum_\mu\mathsf P_n^{(r)}(\mu)K(\nu\mid\mu)
   \cE_{\mu,\nu}\!\left(\frac{\Pi_{\mu,P}}{D_r^\mu}\right),
\]
whereas the target block is
\(\mathsf P_{m_n}^{(r)}(\nu)\Pi_{\nu,P}/D_r^\nu\).
Since
\(\Pi_{\mu,P}/D_r^\mu\le I/D_r^\mu\), positivity gives the uniform
support-overlap bound
\begin{equation}
  \Tr\!\left[
    \Pi_{\nu,P}\,
    \cE_{\mu,\nu}\!\left(\frac{\Pi_{\mu,P}}{D_r^\mu}\right)
  \right]
  \le\frac{R_\mu}{R_\nu}.
  \label{eq:support-overlap}
\end{equation}
For any positive, possibly subnormalized \(A\),
\begin{equation}
  F\!\left(A,\frac{\Pi_{\nu,P}}{D_r^\nu}\right)
  \le\sqrt{\Tr(\Pi_{\nu,P}A)},
  \label{eq:flat-support-fidelity}
\end{equation}
because \(\Tr\sqrt B\le\sqrt{\operatorname{rank}(B)\Tr B}\).
Using the direct-sum fidelity formula, then
\eqref{eq:support-overlap}--\eqref{eq:flat-support-fidelity}, and finally
Cauchy--Schwarz over \(\nu\), gives the finite-\(n\) upper bound
\begin{align}
 &F\!\left(
   \cM_n((P/r)^{\ot n}),
   (P/r)^{\ot m_n}
 \right)\notag\\
 &\quad\le
 \sum_\nu
 \sqrt{\frac{\mathsf P_{m_n}^{(r)}(\nu)}{R_\nu}}
 \sqrt{\sum_\mu
       \mathsf P_n^{(r)}(\mu)K(\nu\mid\mu)R_\mu}
 \notag\\
 &\quad\le
 \left[
   \left(\sum_\mu\mathsf P_n^{(r)}(\mu)R_\mu\right)
   \left(\sum_\nu\frac{\mathsf P_{m_n}^{(r)}(\nu)}{R_\nu}\right)
 \right]^{1/2}.
 \label{eq:grassmann-finite-converse}
\end{align}
Here and in this display, the \(\nu\)-sums may be restricted to
\(\ell(\nu)\le r\), since the target law vanishes otherwise.  Any output
mass in longer diagrams only strengthens the final inequality.

Let \(k=r(d-r)\) and
\[
  C_{r,d}=\prod_{\substack{1\le i\le r\\r<j\le d}}
  \frac{1}{r(j-i)}.
\]
Young concentration at the uniform \(r\)-spectrum and
\eqref{eq:R-lambda} give
\begin{align}
  \frac1{n^k}\sum_\mu\mathsf P_n^{(r)}(\mu)R_\mu
  &\longrightarrow C_{r,d},
  \label{eq:R-positive-moment}\\
  m_n^k\sum_\nu\frac{\mathsf P_{m_n}^{(r)}(\nu)}{R_\nu}
  &\longrightarrow C_{r,d}^{-1}.
  \label{eq:R-negative-moment}
\end{align}
For the first limit, \(n^{-k}R_\mu\) is uniformly bounded.  For the
second, restrict first to \(\nu_r\ge m_n/(2r)\), where the rescaled inverse
is uniformly bounded; the complement has exponentially small probability,
and globally \(m_n^k/R_\nu\le m_n^k\) because \(R_\nu\ge1\).  Thus the
exponential Schur large-deviation bound gives uniform integrability.
Substitution into
\eqref{eq:grassmann-finite-converse} yields
\[
  \limsup_n F\le
  \left(\frac{n}{m_n}\right)^{k/2}
  \longrightarrow\gamma^{-r(d-r)/2},
\]
which completes the converse.

\appendix
\clearpage
\section{Strict comparison with purify--clone--trace}
\label[appendix]{app:pct-comparison}
This appendix upgrades the comparison from PCT's published data-processing
guarantee to the actual mixed-system output after the purifying registers are
discarded.  For every fixed \(\gamma>1\), without assuming a
near-unit cloning fidelity, we obtain an explicit upper bound on that output
and prove that it is strictly below the fidelity achieved by our cloner when
\(r>1\).  For a rank-\(r\) state with simple positive spectrum, the exact
support factorization of
\cref{lem:pct-support-factorization} separates the ambient
support-embedding cost from a PCT channel intrinsic to the positive support.
We analyze that intrinsic channel first and then apply the factorization to
every ambient dimension.  Flat projector states are treated separately by a
commuting-statistics argument.  Thus the comparison is directly between the
fixed-gain asymptotic fidelities of the two cloners, not merely between their
guarantees near fidelity one.

\subsection{Simple spectra on the positive support}

In this subsection \(d\) denotes the intrinsic positive-support dimension,
not a larger ambient dimension.  Fix \(d\ge2\) and
\[
  p_1>\cdots>p_d>0,
  \qquad \sum_{i=1}^d p_i=1,
  \qquad \rho_p=\diag(p_1,\ldots,p_d),
\]
and let \(m_n/n\to\gamma>1\).  We use root fidelity
\(F(\rho,\sigma)=\|\sqrt\rho\sqrt\sigma\|_1\).  Put
\[
  q_{ij}=\frac{p_j}{p_i},
  \qquad
  h_q(x)=\frac{\sqrt{(1-q)(1-x)}}{1-\sqrt{qx}},
\]
and define
\begin{equation}\label{eq:q-pct}
  q_{ij}^{\mathrm{PCT}}
  =\frac{\gamma-1+\gamma q_{ij}}
         {\gamma+(\gamma-1)q_{ij}}.
\end{equation}
The resulting upper bound on the actual post-trace purify--clone--trace output is
\begin{equation}\label{eq:B-pct}
  B_{\mathrm{PCT}}(\gamma,p)
  =\left(\frac{\sqrt{2\gamma-1}}{\gamma}\right)^{(d-1)/2}
   \prod_{i<j}h_{q_{ij}}\!\left(q_{ij}^{\mathrm{PCT}}\right).
\end{equation}

Let \(\sigma_{n,m_n}^{\mathrm{PCT}}(p)\) be the mixed-system output obtained
by random purification in \(\C^d\ot\C^d\), Werner's pure-state
\(n\to m_n\) cloner in dimension \(d^2\), and partial trace over every
purifying register.

\begin{theorem}[Strict comparison on a simple positive support]
\label{thm:general-pct}
For fixed \(d\), \(p\), and \(\gamma>1\) as above,
\begin{equation}\label{eq:main-upper}
  \limsup_{n\to\infty}
  F\!\left(
    \sigma_{n,m_n}^{\mathrm{PCT}}(p),
    \rho_p^{\ot m_n}
  \right)
  \le B_{\mathrm{PCT}}(\gamma,p)
  <\Funiv(\gamma,p)<\Forb(\gamma,p).
\end{equation}
Here
\[
  \Forb(\gamma,p)=\prod_{i<j}f_\gamma(q_{ij}),
  \qquad
  \Funiv(\gamma,p)
  =\left(\frac{2\sqrt\gamma}{1+\gamma}\right)^{(d-1)/2}
   \Forb(\gamma,p),
\]
where
\[
  f_\gamma(q)
  =\frac{\sqrt\gamma+\sqrt{q(\gamma-1+q)}}{\gamma+q}.
\]
\end{theorem}

We first identify the limiting
channel obtained by composing partial trace with interior LAN, and then feed into
that channel the Fock-space limit of Werner's cloner.

\subsubsection{Partial trace followed by interior LAN}

Let \(H=E=\C^d\), set
\[
  S=\diag(\sqrt{p_1},\ldots,\sqrt{p_d}),
  \qquad
  \ket{\psi_p}=\sum_{i=1}^d\sqrt{p_i}\ket{ii},
\]
and identify a vector in \(H\ot E\) with its coefficient matrix:
\[
  \ket Z=\sum_{i,j}Z_{ij}\ket{ij}.
\]
The one-particle purification-tangent space is
\[
  K_p:=\psi_p^\perp
  =K_p^{\mathrm{diag}}
   \oplus\bigoplus_{i<j}\operatorname{span}\{\ket{ij},\ket{ji}\},
\]
where
\[
  K_p^{\mathrm{diag}}
  =\left\{
    \sum_i z_i\ket{ii}:\sum_i\sqrt{p_i}z_i=0
   \right\}.
\]
Let \(\Gamma_s(K_p)\) be the symmetric Fock space over \(K_p\), and let
\(b(u)\) denote the annihilation operator in the one-particle direction
\(u\in K_p\).  We abbreviate
\(b_{ij}=b(\ket{ij})\) for \(i\ne j\).

For a matrix tangent \(Z\in K_p\), the differential of partial trace at the
canonical purification is
\begin{equation}\label{eq:partial-trace-differential}
  D_p(Z)
  :=\left.\frac{\dd}{\dd t}\right|_{t=0}
     \Tr_E\proj{\psi_p+tZ}
  =ZS+SZ^*.
\end{equation}
Its diagonal and positive-CCR orbital coordinates are
\begin{align}
  h_i(Z)&=2\sqrt{p_i}\,\Re Z_{ii},
  \label{eq:h-Z}\\
  \beta_{ij}(Z)
  &=\frac{D_p(Z)_{ji}}{\sqrt{p_i-p_j}}
    =\frac{\sqrt{p_i}Z_{ji}
            +\sqrt{p_j}\,\overline{Z_{ij}}}
           {\sqrt{p_i-p_j}},
  \qquad i<j.
  \label{eq:beta-Z}
\end{align}
Notice that \(\sum_i h_i(Z)=0\).  Set
\[
  \mathsf H_0=\left\{h\in\R^d:\sum_i h_i=0\right\},
  \qquad
  \overline\Sigma_p=\diag(p)-pp^{\mathsf T},
  \qquad
  \Sigma_p=\left.\overline\Sigma_p\right|_{\mathsf H_0}.
\]

\paragraph{The explicit limiting channel.}

We now construct a classical--quantum Gaussian channel
\[
  \Lambda_p:\cT_1(\Gamma_s(K_p))
  \longrightarrow
  L^1\!\left(
    \mathsf H_0;
    \cT_1\bigl(\cF^{\ot\binom d2}\bigr)
  \right).
\]

For the diagonal block, define
\begin{equation}\label{eq:v-i}
  v_i=\sqrt{p_i}\ket{ii}-p_i\ket{\psi_p}
  \in K_p^{\mathrm{diag}},
  \qquad
  X_i=b(v_i)+b(v_i)^*.
\end{equation}
These fields commute.  Indeed,
\begin{equation}\label{eq:v-gram}
  \langle v_i,v_j\rangle
  =p_i\delta_{ij}-p_ip_j=(\overline\Sigma_p)_{ij}
\end{equation}
is real, so \([X_i,X_j]=0\).  Moreover \(\sum_iX_i=0\).  The
diagonal part of \(\Lambda_p\) is the joint spectral measurement of
\((X_1,\ldots,X_d)\), regarded as a classical variable in
\(\mathsf H_0\).

For each \(i<j\), put \(\Delta_{ij}=p_i-p_j\) and define
\begin{align}
  a_{ij}
  &=\frac{\sqrt{p_i}\,b_{ji}
          +\sqrt{p_j}\,b_{ij}^*}
         {\sqrt{\Delta_{ij}}},
  \label{eq:retained-mode}\\
  c_{ij}
  &=\frac{\sqrt{p_i}\,b_{ij}
          +\sqrt{p_j}\,b_{ji}^*}
         {\sqrt{\Delta_{ij}}}.
  \label{eq:complement-mode}
\end{align}
A direct CCR calculation gives
\[
  [a_{ij},a_{ij}^*]=[c_{ij},c_{ij}^*]=1,
  \qquad
  [a_{ij},c_{ij}]=[a_{ij},c_{ij}^*]=0.
\]
Thus \((b_{ij},b_{ji})\mapsto(a_{ij},c_{ij})\) is a two-mode
squeezing transformation.  The \((i,j)\)-part of \(\Lambda_p\) applies
this transformation and discards the complementary mode \(c_{ij}\),
retaining \(a_{ij}\).  Different unordered pairs use disjoint input modes
and commute with the diagonal fields, so these blocks form a product
channel.

Let \(\ket Z\) now denote the Fock coherent vector of amplitude
\(Z\in K_p\).  The preceding construction gives
\begin{equation}\label{eq:Lambda-coherent}
  \Lambda_p(\proj Z)
  =\mathsf N_{h(Z),\Sigma_p}
   \ot\bigotimes_{i<j}
   D(\beta_{ij}(Z))\tau_{q_{ij}}D(\beta_{ij}(Z))^*,
\end{equation}
where
\[
  \tau_q=(1-q)\sum_{r\ge0}q^r\proj r.
\]
To verify the diagonal part, use
\(\langle X_i\rangle_Z=2\Re\langle v_i,Z\rangle=h_i(Z)\) and
\eqref{eq:v-gram}.  For the off-diagonal part,
\(\langle a_{ij}\rangle_Z=\beta_{ij}(Z)\), while on the input vacuum
\[
  \langle a_{ij}^*a_{ij}\rangle
  =\frac{p_j}{p_i-p_j}.
\]
The corresponding geometric parameter is therefore \(p_j/p_i=q_{ij}\).

\paragraph{Strong convergence of the finite channels.}

Let \(R_M\) project onto Fock number vectors of total occupation at most
\(M\), and let \(J_M\) map such a number vector to the normalized symmetric
occupation vector with all remaining particles in \(\psi_p\).  Define the
trace-preserving occupation embedding
\begin{equation}\label{eq:occupation-channel}
  \mathcal J_M(X)
  =J_MR_MXR_MJ_M^*
   +\Tr[(I-R_M)X]\proj{\psi_p}^{\ot M}.
\end{equation}
Let \(T_M^p\) be an interior classical--quantum LAN extraction channel centered
at \(\rho_p\), with its orbital normalization and phase chosen so that its
output annihilator is the LAN limit of the positive-CCR fluctuation
\[
  \frac1{\sqrt{M(p_i-p_j)}}
  \sum_{\ell=1}^M E_{ij}^{(\ell)},
  \qquad i<j.
\]
Define
\begin{equation}\label{eq:finite-Lambda}
  \Lambda_{p,M}
  =T_M^p\circ\Tr_{E^M}\circ\mathcal J_M.
\end{equation}

\begin{lemma}[Interior partial-trace/LAN limit]
\label{lem:partial-trace-LAN}
For every trace-class operator \(X\) on \(\Gamma_s(K_p)\),
\begin{equation}\label{eq:strong-convergence}
  \|\Lambda_{p,M}(X)-\Lambda_p(X)\|_1\longrightarrow0.
\end{equation}
\end{lemma}

\begin{proof}
First take a coherent amplitude \(Z\in K_p\).  Put
\[
  \ket{\psi_{Z,M}}
  =\frac{\ket{\psi_p}+M^{-1/2}\ket Z}
         {\sqrt{1+\|Z\|^2/M}}.
\]
The elementary multinomial-to-Poisson limit gives, uniformly for \(Z\) in
bounded subsets of \(K_p\),
\begin{equation}\label{eq:coherent-product-limit}
  \left\|
    \mathcal J_M(\proj Z)
    -\proj{\psi_{Z,M}}^{\ot M}
  \right\|_1\longrightarrow0.
\end{equation}
After the purifying systems are traced out, the second state in
\eqref{eq:coherent-product-limit} becomes the product state
\(\rho_{Z,M}^{\ot M}\), where
\begin{equation}\label{eq:local-reduced-state}
  \rho_{Z,M}
  =\rho_p+M^{-1/2}(ZS+SZ^*)+O_Z(M^{-1}).
\end{equation}

To make the local chart explicit, define the anti-Hermitian matrix
\begin{equation}\label{eq:G-beta}
  G(\beta)
  =\sum_{i<j}
   \frac{\beta_{ij}E_{ji}-\overline{\beta_{ij}}E_{ij}}
        {\sqrt{p_i-p_j}}.
\end{equation}
It satisfies
\(
  [G(\beta),\rho_p]_{ji}
  =\sqrt{p_i-p_j}\,\beta_{ij}
\), so \eqref{eq:beta-Z} is precisely the off-diagonal coordinate of
\(D_p(Z)\).  Because \(p\) has a simple spectrum, smooth
eigenvalue/eigenvector perturbation theory, with the eigenvector phases fixed
continuously at the standard basis, gives
\[
  \rho_{Z,M}
  =g_M\diag(p+M^{-1/2}h_M)g_M^*,
  \qquad
  h_M\longrightarrow h(Z),
  \qquad
  \sqrt M\log g_M\longrightarrow G(\beta(Z)).
\]
Thus \(\rho_{Z,M}\) lies in the usual interior-LAN local chart with bounded
coordinates converging to \((h(Z),\beta(Z))\).  The
finite-dimensional LAN theorem \cite{KahnGuta2009LANFiniteDimensional}, uniformly on bounded
local parameter sets, therefore gives
\[
  \Lambda_{p,M}(\proj Z)
  \longrightarrow
  \mathsf N_{h(Z),\Sigma_p}
  \ot\bigotimes_{i<j}
  D(\beta_{ij}(Z))\tau_{q_{ij}}D(\beta_{ij}(Z))^*
\]
in trace norm.  This is exactly \eqref{eq:Lambda-coherent}.

It remains to extend the convergence to arbitrary trace-class inputs.  The
complex linear span of coherent projectors is trace-norm dense in the trace
class.  Indeed, a bounded operator annihilating that span has identically
zero Husimi symbol; analyticity then makes every number-basis matrix element
zero, and Hahn--Banach gives density.  Both \(\Lambda_{p,M}\) and
\(\Lambda_p\) are trace-norm contractions.  If \(Y\) is a finite linear
combination of coherent projectors, then
\[
  \limsup_M\|\Lambda_{p,M}(X)-\Lambda_p(X)\|_1
  \le2\|X-Y\|_1.
\]
Approximating \(X\) by such \(Y\) proves \eqref{eq:strong-convergence}.
\end{proof}

\begin{remark}[Coordinate convention]
The proof uses the positive-CCR displacement \(\beta_{ij}\) in
\eqref{eq:beta-Z}.  The generator \(G(\beta)\) in
\eqref{eq:G-beta} is exactly the source-normalized generator
\(X_p(\beta)\) of \cref{sec:local-coordinates}.  Thus the same parameter
\(\beta\) appears in the finite-dimensional local family and in the Gaussian
displacement \(D(\beta)\).
\end{remark}

\paragraph{Action on the thermal tangent state.}

Let \(\tau_\vartheta\) be a one-mode thermal state with mean photon number
\[
  N=\frac{\vartheta}{1-\vartheta}.
\]
Because the product thermal state is invariant under passive changes of
one-particle basis, it factorizes according to the diagonal/off-diagonal
decomposition of \(K_p\).  The joint fields in \eqref{eq:v-i} have covariance
\((2N+1)\Sigma_p\).  For each \(i<j\), the retained mode has mean photon
number
\begin{equation}\label{eq:N-out}
  N_{ij}^{\mathrm{out}}
  =\frac{p_iN+p_j(N+1)}{p_i-p_j}.
\end{equation}
The input is centered Gaussian and gauge invariant, and
\(\langle a_{ij}^2\rangle=0\), so the retained one-mode Gaussian state is
thermal, rather than merely having the displayed mean photon number.
Consequently,
\begin{equation}\label{eq:q-N}
  \frac{N_{ij}^{\mathrm{out}}}{N_{ij}^{\mathrm{out}}+1}
  =\frac{N+(N+1)q_{ij}}{N+1+Nq_{ij}}.
\end{equation}
Taking
\[
  \vartheta=\frac{\gamma-1}{\gamma},
  \qquad N=\gamma-1,
\]
turns \eqref{eq:q-N} into \eqref{eq:q-pct}.  We have proved
\begin{equation}\label{eq:thermal-channel-output}
  \Lambda_p\!\left(
    \tau_\vartheta^{\ot(d^2-1)}
  \right)
  =\mathsf N_{0,(2\gamma-1)\Sigma_p}
   \ot\bigotimes_{i<j}\tau_{q_{ij}^{\mathrm{PCT}}}.
\end{equation}

\subsubsection{The Fock limit of Werner cloning}

We record explicitly the arbitrary-dimensional occupation calculation used
below.  Let \(D\ge2\), fix a unit vector \(\psi\in\C^D\), and choose an
orthonormal basis \(e_1,\ldots,e_{D-1}\) of \(\psi^\perp\).  For
\(\mathbf k=(k_1,\ldots,k_{D-1})\), let \(\ket{\mathbf k}_m\) be the
normalized symmetric occupation vector with \(k_j\) particles in \(e_j\)
and the remaining \(m-|\mathbf k|\) particles in \(\psi\).

\begin{lemma}[Werner occupation law]
\label{lem:Werner-Fock}
Let \(\mathcal W_{n,m}^{(D)}\) be Werner's universal pure-state cloner in
dimension \(D\).  Its output on \(\proj\psi^{\ot n}\) is diagonal in the
occupation basis, with eigenvalue
\begin{equation}\label{eq:occupation-law}
  \pi_{n,m}^{(D)}(\mathbf k)
  =\frac{\binom{m-|\mathbf k|}{n}}
         {\binom{m+D-1}{n+D-1}},
  \qquad |\mathbf k|\le m-n.
\end{equation}
If \(m_n/n\to\gamma>1\), then under the occupation--Fock embedding
\begin{equation}\label{eq:Werner-thermal-limit}
  \mathcal W_{n,m_n}^{(D)}(\proj\psi^{\ot n})
  \longrightarrow
  \tau_{(\gamma-1)/\gamma}^{\ot(D-1)}
\end{equation}
in trace norm.
\end{lemma}

\begin{proof}
Werner's channel \cite{Werner1998} is
\[
  \mathcal W_{n,m}^{(D)}(X)
  =\frac{D[n]}{D[m]}
   P_m(X\ot I^{\ot(m-n)})P_m,
  \qquad
  D[r]=\binom{r+D-1}{D-1},
\]
where \(P_m\) projects onto the symmetric subspace.  On an occupation vector
with total transverse occupation \(t=|\mathbf k|\), the expectation of
\(P_m(\proj\psi^{\ot n}\ot I)P_m\) is
\(\binom{m-t}{n}/\binom mn\).  Multiplying by \(D[n]/D[m]\) gives
\eqref{eq:occupation-law}.  Irreducibility of \(\operatorname{Sym}^t(\C^{D-1})\)
under the stabilizer \(\U(D-1)\) makes the operator scalar on each fixed-\(t\)
subspace, while the stabilizer's central phase eliminates cross-\(t\) blocks.
Thus these are all of its eigenvalues.  Their normalization also follows
from Vandermonde's identity.

For fixed \(\mathbf k\),
\[
  \pi_{n,m_n}^{(D)}(\mathbf k)
  \longrightarrow
  \gamma^{-(D-1)}
  \left(\frac{\gamma-1}{\gamma}\right)^{|\mathbf k|}.
\]
The right side is the product of \(D-1\) normalized geometric laws.  Since
both sides are probability masses on the countable occupation lattice, the
discrete Scheff\'e lemma upgrades pointwise convergence to \(\ell^1\)
convergence, which is exactly the trace-norm convergence
\eqref{eq:Werner-thermal-limit}.
\end{proof}

\subsubsection{Proof of the strict comparison}

\begin{proof}[Proof of Theorem~\ref{thm:general-pct}]
Use the canonical purification \(\psi_p\).  Every other minimal purification
has the form \((I\ot U)\psi_p\), with \(U\in\U(d)\).  Covariance of Werner's
cloner and invariance of partial trace give, for every \(U\),
\begin{align*}
 &\Tr_{E^m}\mathcal W_{n,m}^{(d^2)}\!\left(
   ((I\ot U)\proj{\psi_p}(I\ot U^*))^{\ot n}
 \right)\\
 &\quad=\Tr_{E^m}\!\left[
   (I\ot U)^{\ot m}
   \mathcal W_{n,m}^{(d^2)}(\proj{\psi_p}^{\ot n})
   (I\ot U^*)^{\ot m}
 \right]\\
 &\quad=\Tr_{E^m}
   \mathcal W_{n,m}^{(d^2)}(\proj{\psi_p}^{\ot n}).
\end{align*}
Thus the Haar average in random purification disappears exactly after the
purifying registers are discarded:
\begin{equation}\label{eq:Haar-drops}
  \sigma_{n,m}^{\mathrm{PCT}}(p)
  =\Tr_{E^m}
   \mathcal W_{n,m}^{(d^2)}(\proj{\psi_p}^{\ot n}).
\end{equation}

Pull the Werner output in \eqref{eq:Haar-drops} back to
\(\Gamma_s(K_p)\) using the occupation embedding, and call the resulting
Fock state \(\omega_{n,m}\).  Lemma~\ref{lem:Werner-Fock}, with
\(D=d^2\), gives
\begin{equation}\label{eq:omega-limit}
  \left\|
    \omega_{n,m_n}
    -\tau_{(\gamma-1)/\gamma}^{\ot(d^2-1)}
  \right\|_1\longrightarrow0.
\end{equation}
The definition of \(\Lambda_{p,m}\) and \eqref{eq:Haar-drops} give the exact
identity
\[
  T_{m_n}^p(\sigma_{n,m_n}^{\mathrm{PCT}}(p))
  =\Lambda_{p,m_n}(\omega_{n,m_n}).
\]
By contractivity, Lemma~\ref{lem:partial-trace-LAN}, and
\eqref{eq:omega-limit},
\begin{align*}
 &\left\|
  \Lambda_{p,m_n}(\omega_{n,m_n})
  -\Lambda_p\!\left(
    \tau_{(\gamma-1)/\gamma}^{\ot(d^2-1)}
  \right)
 \right\|_1\\
 &\quad\le
 \left\|
   \omega_{n,m_n}
   -\tau_{(\gamma-1)/\gamma}^{\ot(d^2-1)}
 \right\|_1
 +\left\|
   \Lambda_{p,m_n}\!\left(
     \tau_{(\gamma-1)/\gamma}^{\ot(d^2-1)}
   \right)
   -\Lambda_p\!\left(
     \tau_{(\gamma-1)/\gamma}^{\ot(d^2-1)}
   \right)
 \right\|_1
 \longrightarrow0.
\end{align*}
Using \eqref{eq:thermal-channel-output}, the LAN limit of the actual PCT
output is therefore
\begin{equation}\label{eq:PCT-LAN-output}
  \Theta_p^{\mathrm{PCT}}
  =\mathsf N_{0,(2\gamma-1)\Sigma_p}
   \ot\bigotimes_{i<j}\tau_{q_{ij}^{\mathrm{PCT}}}.
\end{equation}
Ordinary interior LAN applied to the ideal target gives
\begin{equation}\label{eq:target-LAN-output}
  T_{m_n}^p(\rho_p^{\ot m_n})
  \longrightarrow
  \Theta_p
  :=\mathsf N_{0,\Sigma_p}
    \ot\bigotimes_{i<j}\tau_{q_{ij}}
\end{equation}
in trace norm.

Fidelity is nondecreasing under \(T_{m_n}^p\) and continuous under
trace-norm convergence.  Hence
\begin{equation}\label{eq:fidelity-upper-limit}
  \limsup_{n\to\infty}
  F\!\left(
    \sigma_{n,m_n}^{\mathrm{PCT}}(p),
    \rho_p^{\ot m_n}
  \right)
  \le F(\Theta_p^{\mathrm{PCT}},\Theta_p).
\end{equation}
The fidelity on the right factorizes.  In \(d-1\) real dimensions,
\begin{align}
 F\!\left(
   \mathsf N_{0,(2\gamma-1)\Sigma_p},
   \mathsf N_{0,\Sigma_p}
 \right)
 &=\left(
   \frac{2\sqrt{2\gamma-1}}{1+(2\gamma-1)}
  \right)^{(d-1)/2}\notag\\
 &=\left(
   \frac{\sqrt{2\gamma-1}}{\gamma}
  \right)^{(d-1)/2}.
 \label{eq:classical-affinity}
\end{align}
For each thermal pair,
\[
  F(\tau_{q_{ij}^{\mathrm{PCT}}},\tau_{q_{ij}})
  =h_{q_{ij}}(q_{ij}^{\mathrm{PCT}}).
\]
Equations \eqref{eq:fidelity-upper-limit} and
\eqref{eq:classical-affinity} prove the first inequality in
\eqref{eq:main-upper}.

It remains to prove strictness.  The optimal one-mode amplifier has thermal
parameter
\[
  q_{ij}^{\mathrm{opt}}
  =\frac{\gamma-1+q_{ij}}{\gamma},
  \qquad
  f_\gamma(q_{ij})=h_{q_{ij}}(q_{ij}^{\mathrm{opt}}).
\]
Moreover,
\[
  q_{ij}^{\mathrm{opt}}-q_{ij}
  =\frac{(\gamma-1)(1-q_{ij})}{\gamma}>0.
\]
Direct calculation gives
\begin{equation}\label{eq:q-strict}
  q_{ij}^{\mathrm{PCT}}-q_{ij}^{\mathrm{opt}}
  =\frac{q_{ij}(\gamma-1)(1-q_{ij})}
         {\gamma[\gamma+(\gamma-1)q_{ij}]}>0.
\end{equation}
For fixed \(q\in(0,1)\),
\[
  \frac{\dd}{\dd x}\log h_q(x)
  =\frac{\sqrt q-\sqrt x}
         {2\sqrt x(1-x)(1-\sqrt{qx})}<0,
  \qquad x>q.
\]
Thus every quantum factor in \(B_{\mathrm{PCT}}\) is strictly smaller than
the corresponding \(f_\gamma(q_{ij})\).  The classical factors obey
\begin{equation}\label{eq:classical-strict}
  \frac{\sqrt{2\gamma-1}}{\gamma}
  <\frac{2\sqrt\gamma}{1+\gamma},
  \qquad \gamma>1,
\end{equation}
because, after squaring and clearing positive denominators, the difference
is proportional to
\((\gamma-1)^2(2\gamma+1)>0\).  Multiplying the strict modewise inequalities
proves \(B_{\mathrm{PCT}}<\Funiv\).  Finally
\(2\sqrt\gamma/(1+\gamma)<1\), so \(\Funiv<\Forb\).
\end{proof}

\subsection{Exact support factorization}

For every fixed-rank stratum with simple positive spectrum, the ambient PCT fidelity factors into
the support-embedding contribution and the intrinsic positive-support
fidelity analyzed above.  For \(s\ge1\), write
\[
  \mathsf d_s(N)=\dim\operatorname{Sym}^N(\C^s)
  =\binom{N+s-1}{s-1}.
\]
Let \(\sigma_{n,m}^{\mathrm{PCT},(d,r)}(\bar p)\) be the mixed-system PCT
output for the rank-\(r\) state \(\rho_{\bar p}\) on \(\C^d\), using a
minimal purifying register \(\C^r\).  Let
\(\sigma_{n,m}^{\mathrm{PCT},(r,r)}(p)\) denote the corresponding output
when the positive support itself is regarded as the system.

\begin{lemma}[Exact PCT support factorization]
\label{lem:pct-support-factorization}
Fix an isometry \(V:\C^r\to\C^d\), put \(P=VV^*\), and identify
\(\rho_{\bar p}=V\rho_pV^*\).  Then
\begin{align}
 &F\!\left(
   \sigma_{n,m}^{\mathrm{PCT},(d,r)}(\bar p),
   \rho_{\bar p}^{\ot m}
 \right)\notag\\
 &\qquad=\sqrt{s_{n,m}^{(r,d)}}\,
 F\!\left(
   \sigma_{n,m}^{\mathrm{PCT},(r,r)}(p),
   \rho_p^{\ot m}
 \right),
 \label{eq:pct-support-factorization}
\end{align}
where
\begin{equation}
  s_{n,m}^{(r,d)}
  =\frac{\mathsf d_{rd}(n)/\mathsf d_{rd}(m)}
         {\mathsf d_{r^2}(n)/\mathsf d_{r^2}(m)}
  =\prod_{j=r^2}^{rd-1}\frac{n+j}{m+j}.
  \label{eq:pct-support-factor}
\end{equation}
\end{lemma}

\begin{proof}
Choose the canonical minimal purification
\(\ket{\psi_{\bar p}}=(V\ot I_r)\sum_i\sqrt{p_i}\ket{ii}\), and put
\(Q=P\ot I_r\).  The Haar average in random purification disappears after
the purifying registers are traced out, by covariance of Werner's cloner.
The projection \(Q\) has rank \(r^2\), contains
\(\psi_{\bar p}\), and commutes tensorwise with the symmetric projection.
Compressing Werner's channel to \(Q^{\ot m}\) therefore gives
\[
 Q^{\ot m}\mathcal W_{n,m}^{(rd)}
   (\proj{\psi_{\bar p}}^{\ot n})Q^{\ot m}
 =s_{n,m}^{(r,d)}(V\ot I_r)^{\ot m}
  \mathcal W_{n,m}^{(r^2)}(\proj{\psi_p}^{\ot n})
  (V^*\ot I_r)^{\ot m}.
\]
After tracing the purifying systems, the target is supported on
\(P^{\ot m}\).  Fidelity is unchanged by compressing its first argument
to that support, and \(F(aA,B)=\sqrt aF(A,B)\), proving the claim.
\end{proof}

For a rank-\(r\) state with simple positive spectrum, define
\begin{align}
  q_{ij}^{\mathrm{PCT}}
  &=\frac{\gamma-1+\gamma q_{ij}}
          {\gamma+(\gamma-1)q_{ij}},\\
  B_{\mathrm{PCT}}^{(r,d)}(\gamma,p)
  &=\gamma^{-r(d-r)/2}
    \left(\frac{\sqrt{2\gamma-1}}{\gamma}\right)^{(r-1)/2}
    \prod_{1\le i<j\le r}
    h_{q_{ij}}\!\left(q_{ij}^{\mathrm{PCT}}\right).
  \label{eq:rank-pct-bound}
\end{align}

\begin{corollary}[Strict PCT comparison on every simple rank stratum]
\label{cor:rank-pct-comparison}
For \(2\le r\le d\), fixed \(p\in\Delta_r^\circ\), and fixed
\(\gamma>1\),
\[
  \limsup_{n\to\infty}
  F\!\left(
    \sigma_{n,m_n}^{\mathrm{PCT},(d,r)}(\bar p),
    \rho_{\bar p}^{\ot m_n}
  \right)
  \le B_{\mathrm{PCT}}^{(r,d)}(\gamma,p)
  <\Funiv^{(r,d)}(\gamma,p)<\Forb^{(r,d)}(\gamma,p).
\]
For \(r=1\), PCT is Werner's optimal pure-state cloner and equality holds
with \(\Forb^{(1,d)}=\Funiv^{(1,d)}=\gamma^{-(d-1)/2}\).
\end{corollary}

\begin{proof}
The factor in \eqref{eq:pct-support-factor} tends to
\(\gamma^{-r(d-r)}\).  Apply \cref{thm:general-pct} in the positive
support, with its dimension denoted there by \(d=r\), and combine it with
\cref{lem:pct-support-factorization}.  Multiplying its two strict
inequalities by \(\gamma^{-r(d-r)/2}\) gives the result.
\end{proof}
\subsection{Projector states}
\label{sec:grassmann-pct-comparison}

The purify--clone--trace (PCT) cloner combines Werner's pure-state cloner
\cite{Werner1998} with random purification \cite{TangWrightZhandry2025},
then discards the purifying registers
\cite{LiTheilHarrowChuang2026,FanizzaGrinkoScharnhorstSpilecki2026,
JeonSohnOh2026}.  On the projector orbit its actual mixed-system fidelity,
not only its data-processing guarantee, can be compared directly with
\cref{thm:grassmann}.

For \(s\ge1\), write
\[
  \mathsf d_s(N)=\dim\operatorname{Sym}^N(\C^s)
  =\binom{N+s-1}{s-1},
\]
and let \(\Pi_{\mathrm{sym},N}^{(s)}\) be the projection onto this
symmetric subspace.  For \(m\ge n\) and an input \(X\) supported on
\(\operatorname{Sym}^n(\C^s)\), Werner's pure-state \(n\)-to-\(m\) cloner
in dimension \(s\) is
\begin{equation}
  \mathcal W_{n,m}^{(s)}(X)
  =\frac{\mathsf d_s(n)}{\mathsf d_s(m)}
   \Pi_{\mathrm{sym},m}^{(s)}
   \bigl(X\ot I_s^{\ot(m-n)}\bigr)
   \Pi_{\mathrm{sym},m}^{(s)}.
  \label{eq:grassmann-werner-cloner}
\end{equation}
Fix \(P\in\Gr(r,d)\), choose an isometry
\(V_P:\C^r\to\C^d\) with \(V_PV_P^*=P\), and put
\[
  \ket{\Phi_P}
  =\frac1{\sqrt r}\sum_{i=1}^r V_P\ket i\ot\ket i,
  \qquad
  \ket{\Phi_r}=\frac1{\sqrt r}\sum_{i=1}^r\ket{ii}.
\]
Let \(\sigma_{n,m}^{\mathrm{PCT}}(P)\) denote the system output of PCT
on \((P/r)^{\ot n}\), using a purifying register \(\C^r\).

\begin{proposition}[PCT is strictly suboptimal on the projector orbit]
\label{prop:grassmann-pct-comparison}
Define
\begin{equation}
  s_{n,m}^{(r,d)}
  =\frac{\mathsf d_{rd}(n)/\mathsf d_{rd}(m)}
         {\mathsf d_{r^2}(n)/\mathsf d_{r^2}(m)}
  =\prod_{j=r^2}^{rd-1}\frac{n+j}{m+j}.
  \label{eq:grassmann-pct-support-factor}
\end{equation}
There is an exact finite-sample factorization
\begin{equation}
  F\!\left(
    \sigma_{n,m}^{\mathrm{PCT}}(P),(P/r)^{\ot m}
  \right)
  =\sqrt{s_{n,m}^{(r,d)}}\,
   F\!\left(\zeta_{n,m}^{(r)},(I_r/r)^{\ot m}\right),
  \label{eq:grassmann-pct-factorization}
\end{equation}
where
\begin{equation}
  \zeta_{n,m}^{(r)}
  =\Tr_{E^m}\mathcal W_{n,m}^{(r^2)}
    \bigl(\proj{\Phi_r}^{\ot n}\bigr).
  \label{eq:grassmann-pct-internal-state}
\end{equation}
If \(m_n/n\to\gamma>1\), then
\begin{equation}
  \limsup_{n\to\infty}
  F\!\left(
    \sigma_{n,m_n}^{\mathrm{PCT}}(P),(P/r)^{\ot m_n}
  \right)
  \le
  \gamma^{-r(d-r)/2}
  \left(\frac{\sqrt{2\gamma-1}}{\gamma}\right)^{(r-1)/2}.
  \label{eq:grassmann-pct-upper-bound}
\end{equation}
Consequently, for every \(r>1\) and fixed \(\gamma>1\), PCT is strictly
below the optimal Grassmannian value \eqref{eq:grassmann-value}.  For
\(r=1\), the internal factor in
\eqref{eq:grassmann-pct-factorization} is one and PCT reduces to Werner's
optimal pure-state cloner.
\end{proposition}

\begin{proof}
Every purification in the random-purification output has the form
\((I\ot U)\Phi_P\), with \(U\in\U(r)\).  Unitary covariance of
\(\mathcal W_{n,m}^{(rd)}\) and invariance of the partial trace give
\begin{align*}
 &\Tr_{E^m}\mathcal W_{n,m}^{(rd)}\!\left(
   \bigl((I\ot U)\proj{\Phi_P}(I\ot U^*)\bigr)^{\ot n}
 \right)\\
 &\qquad=
 \Tr_{E^m}\!\left[
   (I\ot U)^{\ot m}
   \mathcal W_{n,m}^{(rd)}(\proj{\Phi_P}^{\ot n})
   (I\ot U^*)^{\ot m}
 \right]\\
 &\qquad=
 \Tr_{E^m}\mathcal W_{n,m}^{(rd)}(\proj{\Phi_P}^{\ot n}).
\end{align*}
Thus the Haar average disappears exactly after the purifying registers are
discarded.

Let \(Q=P\ot I_r\) on \(\C^d\ot\C^r\).  Its rank is \(r^2\), its
range contains \(\Phi_P\), and \(Q^{\ot m}\) commutes with the symmetric
projection.  Compressing \eqref{eq:grassmann-werner-cloner} therefore gives
\begin{align}
 &Q^{\ot m}\mathcal W_{n,m}^{(rd)}
   (\proj{\Phi_P}^{\ot n})Q^{\ot m}\notag\\
 &\quad=s_{n,m}^{(r,d)}
 (V_P\ot I_r)^{\ot m}
 \mathcal W_{n,m}^{(r^2)}(\proj{\Phi_r}^{\ot n})
 (V_P^*\ot I_r)^{\ot m}.
 \label{eq:grassmann-pct-purified-compression}
\end{align}
Taking the partial trace over \(E^m\) yields
\begin{equation}
  P^{\ot m}\sigma_{n,m}^{\mathrm{PCT}}(P)P^{\ot m}
  =s_{n,m}^{(r,d)}
   V_P^{\ot m}\zeta_{n,m}^{(r)}(V_P^*)^{\ot m}.
  \label{eq:grassmann-pct-system-compression}
\end{equation}
The target is supported on \(P^{\ot m}\).  Hence root fidelity is
unchanged by compressing the first argument to that support, and its
homogeneity proves \eqref{eq:grassmann-pct-factorization}.

It remains to bound the internal fidelity.  Assume \(r>1\), put
\[
  u_r=\frac1r\mathbf1_r,
  \qquad
  \Sigma_r
  =\left.(\diag(u_r)-u_ru_r^{\mathsf T})\right|_{\mathsf H_r},
\]
and measure every system output in the computational basis.  If
\(C_m=(C_{m,1},\ldots,C_{m,r})\) is the resulting count vector, set
\[
  Z_m=\frac{C_m-mu_r}{\sqrt m}\in\mathsf H_r.
\]
Under the target \((I_r/r)^{\ot m}\), the multinomial central limit theorem
gives
\begin{equation}
  Z_m\Longrightarrow\mathsf N_{0,\Sigma_r}.
  \label{eq:grassmann-target-count-limit}
\end{equation}
We next show that under \(\zeta_{n,m_n}^{(r)}\),
\begin{equation}
  Z_{m_n}\Longrightarrow
  \mathsf N_{0,(2\gamma-1)\Sigma_r}.
  \label{eq:grassmann-pct-count-limit}
\end{equation}

For this purpose put \(R=r^2\) and use the symmetric occupation basis
about the condensate \(\Phi_r\).  Under the corresponding embedding into
the \((R-1)\)-mode Fock space, the purified Werner output is diagonal with
mass
\begin{equation}
  \pi_{n,m}(\mathbf k)
  =\frac{\binom{m-|\mathbf k|}{n}}
         {\binom{m+R-1}{n+R-1}},
  \qquad
  \mathbf k\in\mathbb N^{R-1},\quad |\mathbf k|\le m-n.
  \label{eq:grassmann-pct-occupation-law}
\end{equation}
For every fixed \(\mathbf k\),
\[
  \pi_{n,m_n}(\mathbf k)
  \longrightarrow
  \gamma^{-(R-1)}
  \left(\frac{\gamma-1}{\gamma}\right)^{|\mathbf k|}.
\]
Extend \(\pi_{n,m_n}\) by zero to all of \(\mathbb N^{R-1}\).  Both sides
are normalized probability masses (normalization follows also from
Vandermonde's identity), so the discrete Scheff\'e lemma gives trace-norm
convergence to
\begin{equation}
  \tau_\vartheta^{\ot(R-1)},
  \qquad
  \tau_\vartheta=(1-\vartheta)\sum_{j\ge0}\vartheta^j\proj j,
  \qquad
  \vartheta=\frac{\gamma-1}{\gamma},
  \label{eq:grassmann-pct-fock-limit}
\end{equation}
whose mean occupation in every mode is \(N=\gamma-1\).

To read off the count fluctuations, fix \(h\in\mathsf H_r\) and set
\[
  H_h=\diag(h)\ot I_r,
  \qquad
  v_h=H_h\Phi_r\in\Phi_r^\perp.
\]
Then
\begin{equation}
  \langle v_g,v_h\rangle
  =\frac1r\sum_{i=1}^r g_ih_i
  =g^{\mathsf T}\Sigma_r h.
  \label{eq:grassmann-pct-count-gram}
\end{equation}
Let \(J_m\) be the condensate occupation embedding, \(\mathcal N\) the
Fock number operator, \(b(v)\) the annihilation operator in direction
\(v\), \(\dd\Gamma(A)\) the second quantization of \(A\), and
\(P_\perp=I-\proj{\Phi_r}\).  The standard condensate fluctuation identity
gives, on each finite-particle sector,
\begin{align}
 &J_m^*\left[
   \frac1{\sqrt m}\sum_{\ell=1}^m H_h^{(\ell)}
 \right]J_m\notag\\
 &\quad=
 b^*(v_h)\sqrt{1-\frac{\mathcal N}{m}}
 +\sqrt{1-\frac{\mathcal N}{m}}\,b(v_h)
 +\frac1{\sqrt m}\dd\Gamma(P_\perp H_hP_\perp)
 \longrightarrow b^*(v_h)+b(v_h).
 \label{eq:grassmann-pct-count-field}
\end{align}
The fields on the right commute for \(h\in\mathsf H_r\), because the
inner products in \eqref{eq:grassmann-pct-count-gram} are real.  The
pulled-back operators may be extended to the entire Fock space by replacing
\(\sqrt{1-\mathcal N/m}\) with
\(\sqrt{(1-\mathcal N/m)_+}\) and compressing the last term to
\(\{\mathcal N\le m\}\).  The finite-particle vectors form a common
analytic core, so \eqref{eq:grassmann-pct-count-field} implies
strong-resolvent convergence and hence strong convergence of the bounded
fluctuation exponentials.
Together with the trace-norm state convergence in
\eqref{eq:grassmann-pct-fock-limit}, this gives convergence of the joint
characteristic functions.  In the gauge-invariant thermal state of mean
\(N\),
\[
  \Tr\!\left[
    \tau_\vartheta^{\ot(R-1)}
    e^{it(b^*(v_h)+b(v_h))}
  \right]
  =\exp\!\left(-\frac{t^2}{2}(2N+1)\|v_h\|^2\right).
\]
Since measuring the system counts after tracing out \(E^m\) is the same as
measuring \(H_h\) before that trace, the Cram\'er--Wold theorem,
\(2N+1=2\gamma-1\), and
\eqref{eq:grassmann-pct-count-gram} prove
\eqref{eq:grassmann-pct-count-limit}.

Weak convergence suffices after an arbitrarily small common classical
smoothing.  For \(t>0\), add to \(Z_m\) an independent
\(\mathsf N_{0,t\Sigma_r}\) random vector on \(\mathsf H_r\).  This is a
classical channel, so fidelity monotonicity bounds the internal quantum
fidelity by the root fidelity of the two smoothed count laws.  Convolution
with the fixed nondegenerate Gaussian has densities that converge pointwise
under weak convergence, and hence in \(L^1\) by Scheff\'e's lemma.  The
limiting laws have covariances \((1+t)\Sigma_r\) and
\((2\gamma-1+t)\Sigma_r\).  The \(L^1\) convergence of both densities
implies convergence of their root fidelity, and the limiting Gaussian
root fidelity is
\[
  \left(
    \frac{\sqrt{(1+t)(2\gamma-1+t)}}{\gamma+t}
  \right)^{(r-1)/2}.
\]
It follows, after \(t\downarrow0\), that
\begin{equation}
  \limsup_{n\to\infty}
  F\!\left(\zeta_{n,m_n}^{(r)},(I_r/r)^{\ot m_n}\right)
  \le
  \left(\frac{\sqrt{2\gamma-1}}{\gamma}\right)^{(r-1)/2}.
  \label{eq:grassmann-pct-internal-bound}
\end{equation}
Finally, \eqref{eq:grassmann-pct-support-factor} gives
\[
  s_{n,m_n}^{(r,d)}\longrightarrow\gamma^{-r(d-r)}.
\]
Combining this with \eqref{eq:grassmann-pct-factorization} and
\eqref{eq:grassmann-pct-internal-bound} proves
\eqref{eq:grassmann-pct-upper-bound}.  Its right-hand side is strictly
smaller than \(\gamma^{-r(d-r)/2}\) when \(r>1\), because
\(2\gamma-1<\gamma^2\).  The \(r=1\) assertion follows directly from
\eqref{eq:grassmann-pct-internal-state}.
\end{proof}

The preceding proposition is a direct comparison at every fixed
\(\gamma>1\), and it remains informative when the limiting fidelity is
bounded away from one.  By contrast, the two recent sample-complexity
analyses use the PCT data-processing guarantee for their fidelity-threshold
objective.  In their squared-fidelity convention \(F_{\mathrm{sq}}:=F^2\),
\[
  F_{\mathrm{sq}}\!\left(
    \sigma_{n,m}^{\mathrm{PCT}}(P),(P/r)^{\ot m}
  \right)
  \ge\frac{\mathsf d_{rd}(n)}{\mathsf d_{rd}(m)}
  \longrightarrow\gamma^{-(rd-1)}
  \quad\cite{FanizzaGrinkoScharnhorstSpilecki2026,JeonSohnOh2026}.
\]
Our Grassmannian theorem instead gives the exact squared-fidelity value
\(\gamma^{-r(d-r)}\).  At the limiting squared-fidelity level, this is
larger than the published PCT guarantee by the factor
\(\gamma^{r^2-1}\) for \(r>1\); that comparison of guarantees alone would
not order the actual protocols.  The actual strict ordering is supplied by
\eqref{eq:grassmann-pct-upper-bound}.

\begin{remark}[The near-unit-gain boundary]
The strict comparison itself holds for every fixed \(\gamma>1\) and requires
no high-fidelity assumption.  As \(\gamma\downarrow1\), however, the
particular strictness certificate above becomes second order.  Indeed,
writing \(\gamma=1+\delta\), the additional factor detected by the count
statistic satisfies
\[
  \left(\frac{\sqrt{2\gamma-1}}{\gamma}\right)^{(r-1)/2}
  =1-\frac{r-1}{4}\delta^2+O(\delta^3).
\]
It therefore does not by itself settle the leading sample-complexity
comparison in a joint near-unit-gain limit.  The proof uses only the
\(r-1\) commuting diagonal count directions; a sharper analysis of all
internal directions would require a degenerate central-limit or LAN
statement and is not claimed here.
\end{remark}

\section{Exact-weight PBW calculation of the sector fidelity}
\label[appendix]{app:sector-fidelity}

This appendix proves \cref{prop:sector-fidelity} by comparing the
target sector \(\rho_{p,\nu}\) with the Cartan-sector output
\(\cC_{\mu,\nu}(\rho_{p,\mu})\) on finitely many weight spaces near their
highest weights, then controlling the remaining target mass.
Throughout this appendix, \(d\) is a generic intrinsic support dimension; the
ambient rank-\(r\) construction applies the result with \(d=r\) and then
uses \cref{lem:exact-compression}.
\Cref{lem:PBW-fixed-height} shows that normalized lowering-operator vectors
become asymptotically orthonormal on these spaces, while
\cref{lem:uniform-character-normalization} gives the uniform limit
\(p^\lambda/s_\lambda(p)\to\prod_{i<j}(1-p_j/p_i)\).  The Cartan inclusion
\(V_{\mu,\nu-\mu}\) then splits
each lowering operation between \(\cH_\mu\) and \(\cH_{\nu-\mu}\) with
asymptotically binomial weights, producing the thermal limits \(\tau_p\) and
\(\tau_p^{(\gamma)}\) and hence the fidelity \(\Forb(\gamma,p)\).

Because the calculations involve only weights of height bounded
in terms of \(L\), \(O_L(\cdot)\) may depend on \(L,d,c,C\) in
\(cN\le\Delta_\alpha^\lambda\le CN\) for \(\alpha\in\mathsf R_+\), but is
uniform in \(N\) and all dominant \(\lambda\) satisfying these inequalities.

\subsection{PBW vectors on complete weight spaces}

Identify the positive roots of \(\U(d)\) with
\[
  \mathsf R_+=\{e_i-e_j:1\le i<j\le d\}.
\]
For \(\alpha=e_i-e_j\), set
\[
  F_\alpha=E_{ji},
  \qquad E_\alpha=E_{ij},
  \qquad H_\alpha=[E_\alpha,F_\alpha]=E_{ii}-E_{jj},
  \qquad \Delta_\alpha^\lambda=\lambda_i-\lambda_j.
\]
Fix a total order on \(\mathsf R_+\).  For
\(\mathbf r\in\mathbb Z_{\ge0}^{\mathsf R_+}\), write
\[
  F^{\mathbf r}=\prod_{\alpha\in\mathsf R_+}^{\prec}F_\alpha^{r_\alpha},
  \qquad |\mathbf r|=\sum_\alpha r_\alpha.
\]
For a nonnegative simple-root combination \(\eta\), let
\(\operatorname{ht}(\eta)\) be its simple-root height and put
\[
  \mathsf P(\eta)
  =\left\{\mathbf r:\sum_{\alpha\in\mathsf R_+}r_\alpha\alpha=\eta\right\}.
\]
The set \(\mathsf P(\eta)\) is finite, and the PBW theorem implies that
\(\{F^{\mathbf r}v_\lambda:\mathbf r\in\mathsf P(\eta)\}\) spans
\(\cH_\lambda[\lambda-\eta]\) \cite{FultonHarris1991}.  Define
\begin{equation}\label{eq:PBW-height-subspace}
  \mathcal K_{\lambda,L}
  =\bigoplus_{\operatorname{ht}(\eta)\le L}
    \cH_\lambda[\lambda-\eta].
\end{equation}
Let \(P_{\lambda,L}\) denote the orthogonal projection onto
\(\mathcal K_{\lambda,L}\).

\begin{lemma}[PBW Gram matrices at fixed height]
\label{lem:PBW-fixed-height}
Let \(\lambda=\lambda^{(N)}\) range over dominant weights satisfying
\[
  cN\le\Delta_\alpha^\lambda\le CN
  \qquad(\alpha\in\mathsf R_+)
\]
for fixed \(0<c<C<\infty\).  Define
\begin{equation}\label{eq:normalized-PBW-vectors}
  \widetilde e_{\mathbf r}^\lambda
  =\prod_{\alpha\in\mathsf R_+}^{\prec}
   \frac{F_\alpha^{r_\alpha}}
        {\sqrt{r_\alpha!}(\Delta_\alpha^\lambda)^{r_\alpha/2}}
   v_\lambda.
\end{equation}
For every fixed \(L\), uniformly over
\(\operatorname{ht}(\eta)\le L\) and
\(\mathbf r,\mathbf s\in\mathsf P(\eta)\),
\begin{equation}\label{eq:PBW-Gram-fixed-height}
  \langle\widetilde e_{\mathbf r}^\lambda,
         \widetilde e_{\mathbf s}^\lambda\rangle
  =\delta_{\mathbf r,\mathbf s}+O_L(N^{-1/2}).
\end{equation}
Consequently, for all sufficiently large \(N\), the complete Gram matrix
\(G_{\lambda,\eta}\) on every such exact-weight block is invertible.  If
\begin{equation}\label{eq:exact-block-orthonormalization}
  (e_{\mathbf r}^\lambda)_{\mathbf r\in\mathsf P(\eta)}
  =
  (\widetilde e_{\mathbf r}^\lambda)_{\mathbf r\in\mathsf P(\eta)}
  G_{\lambda,\eta}^{-1/2},
\end{equation}
then these vectors form an orthonormal basis of
\(\cH_\lambda[\lambda-\eta]\), and
\begin{equation}\label{eq:PBW-orthonormal-close}
  e_{\mathbf r}^\lambda
  =\widetilde e_{\mathbf r}^\lambda+O_L(N^{-1/2})
\end{equation}
uniformly for \(\operatorname{ht}(\eta)\le L\) and
\(\mathbf r\in\mathsf P(\eta)\).
\end{lemma}

\begin{proof}
For \(\mathbf r\in\mathsf P(\eta)\), every positive root has positive
simple-root height, so
\(|\mathbf r|\le\operatorname{ht}(\eta)\).  Put
\[
  \widehat F_\alpha=\frac{F_\alpha}{\sqrt{\Delta_\alpha^\lambda}},
  \qquad
  \widehat E_\alpha=\frac{E_\alpha}{\sqrt{\Delta_\alpha^\lambda}}.
\]
Each unnormalized Gram entry
\(\langle v_\lambda,(\widehat F^{\mathbf r})^*
\widehat F^{\mathbf s}v_\lambda\rangle\) is therefore a word of length at
most \(2L\).  Commuting root operators cannot increase the word length, and
each root has height at most \(d-1\).  Hence every intermediate vector lies
in
\[
  \mathcal K_{\lambda,R_L},
  \qquad R_L:=2L(d-1).
\]

For \(0\le R\le R_L\), the restrictions
\[
  \widehat E_\alpha:\mathcal K_{\lambda,R}\to\mathcal K_{\lambda,R},
  \qquad
  \widehat F_\alpha:\mathcal K_{\lambda,R}
  \to\mathcal K_{\lambda,R+d-1}
\]
satisfy
\(\|\widehat E_\alpha\|+\|\widehat F_\alpha\|=O_L(1)\).  To see this,
decompose \(\cH_\lambda\) orthogonally into irreducible
  \(\mathfrak{sl}_2(\alpha)\)-summands.  If a summand meets the weight space
  of \(\lambda-\eta\), where \(\operatorname{ht}(\eta)\le R\), let \(t\) be
  its level there and \(M\) its highest \(H_\alpha\)-eigenvalue.
  The top of this string, of weight
  \(\lambda-\eta+t\alpha\), is again a weight of \(\cH_\lambda\); hence
  \(\eta-t\alpha\) is a nonnegative simple-root combination and
  \(t\operatorname{ht}(\alpha)\le\operatorname{ht}(\eta)\le R\).  Moreover,
\[
  M=\Delta_\alpha^\lambda-\eta(H_\alpha)+2t
   =\Delta_\alpha^\lambda+O_L(1).
\]
At level \(t\), the \(F_\alpha\)- and \(E_\alpha\)-coefficients are
\(\sqrt{(t+1)(M-t)}\) and \(\sqrt{t(M-t+1)}\).  After division by
\(\sqrt{\Delta_\alpha^\lambda}\), both are \(O_L(1)\), proving the norm
bound.  All operator norms below are taken after restriction to
\(\mathcal K_{\lambda,R_L}\), with the codomain norm inherited from
\(\cH_\lambda\).

The matrix-unit commutator formula now gives
\begin{equation}\label{eq:normalized-root-commutator}
  [\widehat E_\alpha,\widehat F_\beta]
  =\delta_{\alpha\beta}I+O_L(N^{-1/2}).
\end{equation}
For \(\alpha=\beta\), the left side is
\(H_\alpha/\Delta_\alpha^\lambda=I+O_L(N^{-1})\) on the cutoff.
For \(\alpha\ne\beta\), the unnormalized commutator is zero or,
up to sign, a root operator \(E_\gamma\) or \(F_\gamma\).  In the nonzero
case, the preceding bound gives
\(\|[\widehat E_\alpha,\widehat F_\beta]\|
=O_L\bigl(\sqrt{\Delta_\gamma^\lambda}/
\sqrt{\Delta_\alpha^\lambda\Delta_\beta^\lambda}\bigr)
=O_L(\sqrt N/N)=O_L(N^{-1/2})\).

Commute every \(\widehat E_\alpha\) to the right in
\(\langle v_\lambda,(\widehat F^{\mathbf r})^*
\widehat F^{\mathbf s}v_\lambda\rangle\).  The scalar commutators in
\eqref{eq:normalized-root-commutator} contribute only when
\(\mathbf r=\mathbf s\), in which case they give
\(\prod_\alpha r_\alpha!\).  Every other term vanishes or contains an
\(O_L(N^{-1/2})\) remainder.  Only \(O_L(1)\) terms occur, and all remaining
operator factors have norm \(O_L(1)\); consequently,
\[
  \langle v_\lambda,(\widehat F^{\mathbf r})^*
  \widehat F^{\mathbf s}v_\lambda\rangle
  =\delta_{\mathbf r,\mathbf s}\prod_\alpha r_\alpha!
   +O_L(N^{-1/2}).
\]
Dividing by the factorials in \eqref{eq:normalized-PBW-vectors} proves
\eqref{eq:PBW-Gram-fixed-height}.

The number and sizes of the exact-weight blocks are bounded in terms of
\(L\), so the entrywise estimate gives
\(G_{\lambda,\eta}=I+O_L(N^{-1/2})\) in operator norm.  Thus
\(G_{\lambda,\eta}\) is invertible for large \(N\), and the PBW spanning
vectors form a basis.  Finally, functional calculus gives
\(G_{\lambda,\eta}^{-1/2}=I+O_L(N^{-1/2})\); hence
\eqref{eq:exact-block-orthonormalization} is an orthonormal basis and
satisfies \eqref{eq:PBW-orthonormal-close}.
\end{proof}

\subsection{Character normalization and fixed-height matrices}

For \(a=(a_1,\ldots,a_d)\), write \(p^a=\prod_i p_i^{a_i}\).

\begin{lemma}[Uniform character normalization]
\label{lem:uniform-character-normalization}
Let \(K\Subset\Delta_d^\circ\), let \(\delta_N\to0\), and suppose
\(p\in K\), \(\lambda\vdash N\), and
\(\|\lambda/N-p\|_\infty\le\delta_N\).  Then, uniformly over this family,
\begin{equation}\label{eq:uniform-character-normalization}
  \frac{p^\lambda}{s_\lambda(p)}
  =\prod_{i<j}\left(1-\frac{p_j}{p_i}\right)+o(1).
\end{equation}
\end{lemma}

\begin{proof}
Write the numerator in the Weyl character formula as
\[
  \sum_{\sigma\in\mathfrak S_d}\operatorname{sgn}(\sigma)A_\sigma,
  \qquad
  A_\sigma=\prod_i p_i^{\lambda_{\sigma(i)}+d-\sigma(i)}.
\]
For \(\sigma\ne\mathrm{id}\),
\[
  \frac1N\log\frac{A_{\mathrm{id}}}{A_\sigma}
  =\sum_i\frac{\lambda_i}{N}\log p_i
   -\sum_i\frac{\lambda_{\sigma(i)}}{N}\log p_i+O_K(N^{-1}).
\]
For \(x\in\mathbb R^d\), put
\[
  D_\sigma(x,p)
  =\sum_i x_i\log p_i-\sum_i x_{\sigma(i)}\log p_i.
\]
Strict rearrangement gives \(D_\sigma(p,p)>0\).  Since \(K\) is compact,
the permutations are finite, and the logarithms are uniformly bounded on
\(K\), the assumptions \(\|\lambda/N-p\|_\infty\le\delta_N\) and
\(\delta_N\to0\) give \(c_K>0\) such that, for every
\(\sigma\ne\mathrm{id}\) and all sufficiently large \(N\),
\[
  \frac1N\log\frac{A_{\mathrm{id}}}{A_\sigma}\ge c_K,
  \qquad
  A_\sigma\le e^{-c_KN}A_{\mathrm{id}}.
\]
Consequently,
\[
  \sum_{\sigma\in\mathfrak S_d}\operatorname{sgn}(\sigma)A_\sigma
  =A_{\mathrm{id}}\bigl(1+O_K(e^{-c_KN})\bigr),
  \qquad
  A_{\mathrm{id}}=p^\lambda\prod_i p_i^{d-i}.
\]
Since
\[
  \det(p_i^{d-j})_{i,j=1}^d
  =\prod_i p_i^{d-i}
   \prod_{i<j}\left(1-\frac{p_j}{p_i}\right),
\]
the Weyl character formula gives, uniformly,
\[
  s_\lambda(p)
  =p^\lambda
   \prod_{i<j}\left(1-\frac{p_j}{p_i}\right)^{-1}
   \bigl(1+O_K(e^{-c_KN})\bigr).
\]
This proves \eqref{eq:uniform-character-normalization}.
\end{proof}

For the remainder of the appendix, index the number basis of
\(\cF_{\mathrm{orb}}\) by
\(\mathbf r\in\mathbb Z_{\ge0}^{\mathsf R_+}\), and write it as
\(\{\ket{\mathbf r}\}\).  Define the finite-rank projection
\begin{equation}\label{eq:Fock-height-projection}
  P_L^{\cF}
  =\sum_{\operatorname{ht}(\sum_\alpha r_\alpha\alpha)\le L}
    \proj{\mathbf r}.
\end{equation}
For every fixed \(L\) and all sufficiently large \(N\), each highest
weight \(\lambda\) satisfying the gap bounds in
\cref{lem:PBW-fixed-height} has the orthonormal basis
\(e_{\mathbf r}^\lambda\) of \(\mathcal K_{\lambda,L}\) constructed there.
Define
\begin{equation}\label{eq:finite-coordinate-unitary}
  J_{\lambda,L}:\mathcal K_{\lambda,L}\longrightarrow
  P_L^{\cF}\cF_{\mathrm{orb}},
  \qquad
  J_{\lambda,L}e_{\mathbf r}^\lambda=\ket{\mathbf r}.
\end{equation}
By \cref{lem:PBW-fixed-height}, the two spaces have bases indexed
by the same \(\mathbf r\), so \(J_{\lambda,L}\) is unitary onto
\(P_L^{\cF}\cF_{\mathrm{orb}}\).  It is defined only on
\(\mathcal K_{\lambda,L}\), and it depends on \(\lambda\) and \(L\).

\begin{proof}[Proof of \cref{prop:sector-fidelity}]
Fix \(K\) and \(\delta_n\) as in the proposition, write \(m=m_n\), and put
\(\omega=\nu-\mu\).  Since \(m/n=\gamma_n\to\gamma>1\), we have
\(m/n-1\ge(\gamma-1)/2\) for all sufficiently large \(n\).
Uniformly over the allowed parameters,
\[
\begin{aligned}
  \left\|\frac\omega{m-n}-p\right\|_\infty
  &=\left\|
    \frac{m/n}{m/n-1}\left(\frac\nu m-p\right)
    -\frac1{m/n-1}\left(\frac\mu n-p\right)
  \right\|_\infty\\
  &\le\frac{m/n+1}{m/n-1}\,\delta_n=o(1).
\end{aligned}
\]
Every coordinate and every adjacent gap is uniformly positive on \(K\).
Hence \(\omega\vdash m-n\) for all sufficiently large \(n\).  For
\(\alpha=e_i-e_j\), set
\begin{equation}\label{eq:splitting-fraction}
\begin{aligned}
  s_{\alpha,n}
  &=\frac{\Delta_\alpha^\mu}{\Delta_\alpha^\nu}
    =\frac{n}{m}
      \frac{\mu_i/n-\mu_j/n}{\nu_i/m-\nu_j/m}
    =\frac1\gamma+o(1),\\
  1-s_{\alpha,n}
  &=\frac{\Delta_\alpha^\omega}{\Delta_\alpha^\nu}.
\end{aligned}
\end{equation}
Consequently, there are constants
\(0<c_K<C_K<\infty\), depending only on \(K\) and
\(\gamma\), such that, for all large \(n\),
\begin{equation}\label{eq:uniform-sector-gaps}
  c_Kn
  \le\Delta_\alpha^\mu,\Delta_\alpha^\nu,\Delta_\alpha^\omega
  \le C_Kn
  \qquad(\alpha\in\mathsf R_+).
\end{equation}

Let
\(V=V_{\mu,\omega}:\cH_\nu\to\cH_\mu\ot\cH_\omega\), and write
\(F_\alpha^{(\lambda)}=\dd\pi_\lambda(F_\alpha)\) for the action of the root
operator on \(\cH_\lambda\).  Intertwining and the highest-weight
normalization give
\[
  VF_\alpha^{(\nu)}
  =\bigl(F_\alpha^{(\mu)}\ot I+I\ot F_\alpha^{(\omega)}\bigr)V.
\]
Expanding each
\(\bigl(F_\alpha^{(\mu)}\ot I+I\ot F_\alpha^{(\omega)}\bigr)^{r_\alpha}\)
binomially gives the exact identity
\begin{equation}\label{eq:cartan-PBW-exact}
  V\widetilde e_{\mathbf r}^{\nu}
  =\sum_{\mathbf a\le\mathbf r}
    \beta_{\mathbf r,\mathbf a}^{(n)}
    \widetilde e_{\mathbf a}^{\mu}
    \ot\widetilde e_{\mathbf r-\mathbf a}^{\omega},
\end{equation}
where
\begin{equation}\label{eq:cartan-coefficients}
  \beta_{\mathbf r,\mathbf a}^{(n)}
  =\prod_{\alpha\in\mathsf R_+}
   \sqrt{\binom{r_\alpha}{a_\alpha}
   s_{\alpha,n}^{a_\alpha}
   (1-s_{\alpha,n})^{r_\alpha-a_\alpha}}.
\end{equation}
If \(\operatorname{ht}(\sum_\alpha r_\alpha\alpha)\le L\),
then \eqref{eq:cartan-PBW-exact} has \(O_L(1)\) terms, all of height at
most \(L\), with
\(0\le\beta_{\mathbf r,\mathbf a}^{(n)}\le1\).
Since \(m/n\to\gamma>1\), the sizes \(n,m,m-n\) are uniformly comparable, so
\eqref{eq:uniform-sector-gaps} and \cref{lem:PBW-fixed-height} give
\(e_{\mathbf b}^{\lambda}=\widetilde e_{\mathbf b}^{\lambda}
+O_L(n^{-1/2})\) for \(\lambda\in\{\mu,\nu,\omega\}\) and all relevant
\(\mathbf b\).  Substituting these relations into
\eqref{eq:cartan-PBW-exact} and using \(\|Vx\|=\|x\|\) gives
\begin{equation}\label{eq:cartan-orthonormal-fixed-height}
  Ve_{\mathbf r}^{\nu}
  =\sum_{\mathbf a\le\mathbf r}
    \beta_{\mathbf r,\mathbf a}^{(n)}
    e_{\mathbf a}^{\mu}\ot e_{\mathbf r-\mathbf a}^{\omega}
    +O_L(n^{-1/2})
\end{equation}
uniformly in norm for
\(\operatorname{ht}(\sum_\alpha r_\alpha\alpha)\le L\).

Put
\[
  \sigma_{\mu,\nu}=\cC_{\mu,\nu}(\rho_{p,\mu}),
  \qquad
  c_\lambda(p)=\frac{p^\lambda}{s_\lambda(p)},
  \qquad
  q_\alpha=\frac{p_j}{p_i}\quad(\alpha=e_i-e_j).
\]
For \(\lambda\in\{\mu,\nu\}\), the weight of
\(e_{\mathbf b}^\lambda\) is
\(\lambda-\sum_\alpha b_\alpha\alpha\), and hence
\[
  \rho_{p,\lambda}e_{\mathbf b}^\lambda
  =c_\lambda(p)\prod_\alpha q_\alpha^{b_\alpha}e_{\mathbf b}^\lambda.
\]
The sector channel gives the local formula
\[
  \sigma_{\mu,\nu}
  =\frac{\dim\cH_\mu}{\dim\cH_\nu}
   V^*(\rho_{p,\mu}\ot I_{\cH_\omega})V.
\]
All asymptotic estimates below are uniform over the allowed family.  The
character estimate, \cref{lem:uniform-character-normalization},
gives
\[
  c_\lambda(p)=\prod_\alpha(1-q_\alpha)+o(1)
  \quad(\lambda=\mu,\nu),
\]
while the Weyl dimension formula, \eqref{eq:uniform-sector-gaps}, and
\eqref{eq:splitting-fraction} give
\[
  \frac{\dim\cH_\mu}{\dim\cH_\nu}
  =\prod_{\alpha=e_i-e_j}
    \frac{\Delta_\alpha^\mu+j-i}{\Delta_\alpha^\nu+j-i}
  =\prod_{\alpha\in\mathsf R_+}
    \bigl(s_{\alpha,n}+O(n^{-1})\bigr)
  =\gamma^{-r_d}+o(1).
\]
Thus, substituting
\eqref{eq:cartan-orthonormal-fixed-height} and the weight action above into
this formula, and using
\(\|\rho_{p,\mu}\|\le1\), gives, for fixed \(L\),
\begin{align}
  \langle e_{\mathbf r}^\nu,
    \sigma_{\mu,\nu}e_{\mathbf s}^\nu\rangle
  &={}\frac{\dim\cH_\mu}{\dim\cH_\nu}
  \sum_{\substack{\mathbf a\le\mathbf r\\
                   \mathbf b\le\mathbf s}}
  \beta_{\mathbf r,\mathbf a}^{(n)}
  \beta_{\mathbf s,\mathbf b}^{(n)}
  \langle e_{\mathbf a}^\mu,
    \rho_{p,\mu}e_{\mathbf b}^\mu\rangle
  \langle e_{\mathbf r-\mathbf a}^\omega,
    e_{\mathbf s-\mathbf b}^\omega\rangle
  +O_L(n^{-1/2})\notag\\
  &={}
  \delta_{\mathbf r,\mathbf s}
  \frac{\dim\cH_\mu}{\dim\cH_\nu}c_\mu(p)
  \prod_{\alpha\in\mathsf R_+}
  \bigl(1-s_{\alpha,n}+s_{\alpha,n}q_\alpha\bigr)^{r_\alpha}
  +O_L(n^{-1/2}),
  \label{eq:sector-output-fixed-matrix}\\
  \langle e_{\mathbf r}^\nu,
    \rho_{p,\nu}e_{\mathbf s}^\nu\rangle
  &=\delta_{\mathbf r,\mathbf s}
    c_\nu(p)\prod_{\alpha\in\mathsf R_+}q_\alpha^{r_\alpha}.
  \label{eq:sector-target-fixed-matrix}
\end{align}

Recall the following notation from \cref{sec:gaussian-amplification}:
\[
  q_\alpha^{(\gamma)}=1-\frac{1-q_\alpha}{\gamma},
  \qquad
  \tau_q=(1-q)\sum_{k\ge0}q^k\proj{k},
  \qquad
  \tau_p=\bigotimes_{\alpha\in\mathsf R_+}\tau_{q_\alpha},
  \quad
  \tau_p^{(\gamma)}
  =\bigotimes_{\alpha\in\mathsf R_+}\tau_{q_\alpha^{(\gamma)}}.
\]
Since \(s_{\alpha,n}=\gamma^{-1}+o(1)\),
\eqref{eq:sector-output-fixed-matrix} and
\eqref{eq:sector-target-fixed-matrix} now give
\[
\begin{aligned}
  \langle e_{\mathbf r}^\nu,
    \sigma_{\mu,\nu}e_{\mathbf s}^\nu\rangle
  &=\delta_{\mathbf r,\mathbf s}
    \prod_\alpha(1-q_\alpha^{(\gamma)})
    \bigl(q_\alpha^{(\gamma)}\bigr)^{r_\alpha}+o(1),\\
  \langle e_{\mathbf r}^\nu,
    \rho_{p,\nu}e_{\mathbf s}^\nu\rangle
  &=\delta_{\mathbf r,\mathbf s}
    \prod_\alpha(1-q_\alpha)q_\alpha^{r_\alpha}+o(1).
\end{aligned}
\]
These are precisely the number-basis matrix elements of
\(\tau_p^{(\gamma)}\) and \(\tau_p\).
For fixed \(L\), the cutoff space has fixed finite dimension, so the uniform
entrywise limits imply convergence in trace norm:
\begin{align}
  \left\|
  J_{\nu,L}P_{\nu,L}\sigma_{\mu,\nu}P_{\nu,L}J_{\nu,L}^*
  -P_L^{\cF}\tau_p^{(\gamma)}P_L^{\cF}
  \right\|_1&\longrightarrow0,
  \label{eq:output-fixed-height-convergence}\\
  \left\|
  J_{\nu,L}P_{\nu,L}\rho_{p,\nu}P_{\nu,L}J_{\nu,L}^*
  -P_L^{\cF}\tau_pP_L^{\cF}
  \right\|_1&\longrightarrow0.
  \label{eq:target-fixed-height-convergence}
\end{align}

Both \(\sigma_{\mu,\nu}\) and \(\rho_{p,\nu}\) commute with
the diagonal torus.  The second assertion is immediate; the first follows
from torus invariance of \(\rho_{p,\mu}\) and covariance of
\(\cC_{\mu,\nu}\).
Consequently, \(P_{\nu,L}\) reduces both states, and the direct-sum identity
gives
\[
\begin{aligned}
  F(\sigma_{\mu,\nu},\rho_{p,\nu})
  ={}&F(P_{\nu,L}\sigma_{\mu,\nu}P_{\nu,L},
        P_{\nu,L}\rho_{p,\nu}P_{\nu,L})\\
  &+F((I-P_{\nu,L})\sigma_{\mu,\nu}(I-P_{\nu,L}),
       (I-P_{\nu,L})\rho_{p,\nu}(I-P_{\nu,L})).
\end{aligned}
\]
Schatten H\"older gives \(F(A,C)\le\sqrt{\Tr A\,\Tr C}\), so the second
term is at most
\begin{equation}\label{eq:sector-tail-fidelity-bound}
  \sqrt{\Tr[(I-P_{\nu,L})\rho_{p,\nu}]},
\end{equation}
because \(\Tr[(I-P_{\nu,L})\sigma_{\mu,\nu}]\le1\).

Taking traces in \eqref{eq:target-fixed-height-convergence} gives, for fixed
\(L\),
\[
  \Tr[(I-P_{\nu,L})\rho_{p,\nu}]
  =\Tr[(I-P_L^{\cF})\tau_p]+o(1)
\]
uniformly as \(n\to\infty\).  Since \(K\) is
compact in \(\Delta_d^\circ\), there is \(b_K<1\) such that
\(q_\alpha\le b_K\) for every \(p\in K\) and every positive root
\(\alpha\).  The occupation numbers under \(\tau_p\) are independent
geometric variables with parameters at most \(b_K\).  If their total
occupation is \(s\), then their total root height is at most \((d-1)s\).
Moreover, the probability of any fixed occupation vector of total occupation
\(s\) is at most \(b_K^s\), and there are
\(\binom{s+r_d-1}{r_d-1}\) such vectors.  Therefore
\begin{equation}\label{eq:uniform-thermal-height-tail}
  \sup_{p\in K}\Tr[(I-P_L^{\cF})\tau_p]
  \le
  \sum_{s>L/(d-1)}\binom{s+r_d-1}{r_d-1}b_K^s
  \longrightarrow0\qquad(L\to\infty).
\end{equation}
The same direct-sum decomposition for
\(F(\tau_p^{(\gamma)},\tau_p)\) bounds its complementary part by the square
root of the same target thermal tail.  Consequently,
\[
\begin{aligned}
 &\left|F(\sigma_{\mu,\nu},\rho_{p,\nu})
       -F(\tau_p^{(\gamma)},\tau_p)\right|\\
 &\quad\le
 \left|F(P_{\nu,L}\sigma_{\mu,\nu}P_{\nu,L},
         P_{\nu,L}\rho_{p,\nu}P_{\nu,L})
       -F(P_L^{\cF}\tau_p^{(\gamma)}P_L^{\cF},
          P_L^{\cF}\tau_pP_L^{\cF})\right|\\
 &\qquad+
 \sqrt{\Tr[(I-P_{\nu,L})\rho_{p,\nu}]}
 +\sqrt{\Tr[(I-P_L^{\cF})\tau_p]}.
\end{aligned}
\]

By \eqref{eq:fidelity-continuity} and
\eqref{eq:output-fixed-height-convergence}--\eqref{eq:target-fixed-height-convergence},
for every fixed \(L\),
\[
\begin{aligned}
  &F(P_{\nu,L}\sigma_{\mu,\nu}P_{\nu,L},
     P_{\nu,L}\rho_{p,\nu}P_{\nu,L})\\
  &\qquad\longrightarrow
  F(P_L^{\cF}\tau_p^{(\gamma)}P_L^{\cF},
    P_L^{\cF}\tau_pP_L^{\cF})
\end{aligned}
\]
uniformly.  First let \(n\to\infty\) at fixed \(L\), and then let
\(L\to\infty\) using
\eqref{eq:sector-tail-fidelity-bound}--\eqref{eq:uniform-thermal-height-tail}.
This proves
\[
  \left|F(\sigma_{\mu,\nu},\rho_{p,\nu})
  -F(\tau_p^{(\gamma)},\tau_p)\right|\longrightarrow0
\]
uniformly.  Finally,
\[
  F(\tau_p^{(\gamma)},\tau_p)
  =\prod_{i<j}F(\tau_{q_{ij}^{(\gamma)}},\tau_{q_{ij}})
  =\Forb(\gamma,p).
\]
Here multiplicativity of fidelity follows directly from multiplicativity of
the trace norm under tensor products, and the final equality uses
\eqref{eq:thermal-root-fidelity} and \eqref{eq:one-mode-f}.
\end{proof}

\section{Young local limit and randomized rounding}
\label[appendix]{app:young-rounding}

This appendix proves \cref{prop:young-rounding} by expressing the
randomized rounding law \(D_n(\widetilde\nu\mid\mu)\) as
the normalized intersection volume of the dilated source cell
around \(\gamma_n\mu\) with the target cell labeled by \(\widetilde\nu\).
Here \(d\) is again the active label dimension; in the fixed-rank
construction it is instantiated as \(r\), independently of the ambient
Hilbert-space dimension.
\Cref{lem:young-local-limit} shows that the centered and rescaled target label
law \(P_{N,p}\) converges uniformly to \(\mathsf N_{0,\Sigma_p}\), while
\cref{lem:randomized-dilation} shows that sending \(P_{n,p}\) through \(D_n\)
changes this limit to \(\mathsf N_{0,\gamma\Sigma_p}\) and yields the label
affinity \(\Fcl(\gamma,d)\).  The final proof checks that typical rounded
labels satisfy \(\nu-\mu\vdash m_n-n\) and that replacement by
\(\nu_n^{\mathrm{fb}}\) moves only vanishing mass.

Define the half-open
fundamental cell
\[
  \mathsf C
  =\left\{\sum_{i=1}^{d-1}t_i(e_i-e_d):-\frac12\le t_i<\frac12\right\}
  \subset\mathsf H.
\]
Volumes are taken with the induced \((d-1)\)-dimensional Lebesgue measure.
The rounding rule \eqref{eq:randomized-rounding} is equivalent, outside a
null set of cell boundaries, to choosing the unique
\(\widetilde\nu\in\mathsf L_{m_n}\) for which
\[
  \gamma_n(\mu+U)\in\widetilde\nu+\mathsf C,
  \qquad
  U=\sum_{i=1}^{d-1}U_i(e_i-e_d).
\]
For \(i<d\), this is precisely the nearest-integer rule because
\(\lfloor x+\tfrac12\rfloor\) is the unique integer in
\((x-\tfrac12,x+\tfrac12]\); the \(d\)-th coordinate is then automatic
because both sides sum to \(m_n\).
Since \(U\) is uniform on \(\mathsf C\), the change of variables
\(x=\gamma_n(\mu+u)\) gives \(\dd x=\gamma_n^{d-1}\dd u\) and sends the
event domain to
\(\gamma_n(\mu+\mathsf C)\cap(\widetilde\nu+\mathsf C)\).
Consequently,
\begin{equation}\label{eq:cell-kernel}
  D_n(\widetilde\nu\mid\mu)
  =\frac{|\gamma_n(\mu+\mathsf C)\cap(\widetilde\nu+\mathsf C)|}
         {\gamma_n^{d-1}|\mathsf C|}.
\end{equation}
All label laws below are regarded as functions on
\(\mathsf L_N\): if a law \(R\) is originally defined on \(\mathsf Y_N\),
set \(R(\lambda)=0\) for every
\(\lambda\in\mathsf L_N\setminus\mathsf Y_N\).

For \(\lambda\in\mathsf L_N\), define the recentered and
rescaled lattice cell
\[
  \mathsf C_{N,p}(\lambda)=\frac{\lambda-Np+\mathsf C}{\sqrt N}.
\]
For any \(R:\mathsf L_N\to\R\) satisfying
\(\sum_{\lambda\in\mathsf L_N}|R(\lambda)|<\infty\), define
\[
  (\mathcal E_{N,p}R)(x)
  =\frac{N^{(d-1)/2}}{|\mathsf C|}R(\lambda)
  \quad(x\in\mathsf C_{N,p}(\lambda)).
\]
The cells \(\mathsf C_{N,p}(\lambda)\), \(\lambda\in\mathsf L_N\),
tile \(\mathsf H\), and their boundaries have zero induced Lebesgue measure.
The map \(\mathcal E_{N,p}\) preserves \(\ell^1\)-distance,
\(\|\mathcal E_{N,p}R-\mathcal E_{N,p}S\|_1=\|R-S\|_{\ell^1}\), and, for
nonnegative \(R,S\), Hellinger affinity,
\(\int_{\mathsf H}\sqrt{(\mathcal E_{N,p}R)(\mathcal E_{N,p}S)}\,\dd x
=\sum_{\lambda\in\mathsf L_N}\sqrt{R(\lambda)S(\lambda)}\).
Indeed, each \(\mathsf C_{N,p}(\lambda)\) has volume
\(|\mathsf C|N^{-(d-1)/2}\), which cancels the prefactor in both cellwise
integrals.

\begin{lemma}[Uniform Young local limit]\label{lem:young-local-limit}
For every compact \(K\Subset\Delta_d^\circ\),
\begin{equation}\label{eq:young-local-limit}
  \sup_{p\in K}
  \left\|\mathcal E_{N,p}P_{N,p}-\varphi_{0,\Sigma_p}\right\|_1
  \longrightarrow0.
\end{equation}
\end{lemma}

\begin{proof}
Let \(M_{N,p}\) be the multinomial law on nonnegative integer
compositions of \(N\), extended by zero to \(\mathsf L_N\).  We first show
\begin{equation}\label{eq:young-multinomial-l1}
  \sup_{p\in K}\|P_{N,p}-M_{N,p}\|_{\ell^1}\longrightarrow0.
\end{equation}
Set \(\mathsf W_{N,p}=\{\lambda\in\mathsf L_N:
\|\lambda-Np\|_\infty\le N^{2/3}\}\).  For all large \(N\),
\(\mathsf W_{N,p}\cap\operatorname{supp}M_{N,p}\subset\mathsf Y_N\), since
\(\lambda_i-\lambda_{i+1}\ge N(p_i-p_{i+1})-2N^{2/3}>0\) for \(i<d\),
uniformly over \(p\in K\).

For \(\lambda\in\mathsf Y_N\), the Weyl dimension and character formulas
give
\[
\begin{gathered}
  P_{N,p}(\lambda)=\dim\cS_\lambda\,s_\lambda(p),\qquad
  \dim\cS_\lambda
  =N!\frac{\prod_{i<j}(\lambda_i-\lambda_j+j-i)}
  {\prod_i(\lambda_i+d-i)!},\\
  s_\lambda(p)
  =\frac{\det[p_i^{\lambda_j+d-j}]_{i,j}}
  {\prod_{i<j}(p_i-p_j)}.
\end{gathered}
\]
Uniformly on \(\mathsf W_{N,p}\cap\operatorname{supp}M_{N,p}\),
\[
  s_\lambda(p)
  =\frac{\prod_i p_i^{\lambda_i+d-i}}
  {\prod_{i<j}(p_i-p_j)}\bigl(1+O_K(e^{-c_KN})\bigr)
\]
for some \(c_K>0\).  Indeed, both \((\lambda_i+d-i)/N\) and
\(\log p_i\) are uniformly strictly decreasing there, so strict
rearrangement makes every nonidentity determinant term exponentially
smaller than the identity term.

Comparison with
\(M_{N,p}(\lambda)=N!\prod_i p_i^{\lambda_i}/\prod_i\lambda_i!\) gives
\[
  \frac{P_{N,p}(\lambda)}{M_{N,p}(\lambda)}
  =\bigl(1+O_K(e^{-c_KN})\bigr)
  \prod_{i<j}\frac{\lambda_i-\lambda_j+j-i}{N(p_i-p_j)}
  \prod_i\frac{\lambda_i!(Np_i)^{d-i}}{(\lambda_i+d-i)!}.
\]
Writing the last product as
\(\prod_i\prod_{a=1}^{d-i}\frac{Np_i}{\lambda_i+a}\), each factor
\(\frac{\lambda_i-\lambda_j+j-i}{N(p_i-p_j)}\), \(i<j\), and
\(\frac{Np_i}{\lambda_i+a}\), \(1\le a\le d-i\), equals
\(1+O_K(N^{-1/3})\).  Therefore
\[
\begin{aligned}
  \|P_{N,p}-M_{N,p}\|_{\ell^1}
  \le{}&\sup_{\lambda\in\mathsf W_{N,p}\cap\operatorname{supp}M_{N,p}}
       \left|\frac{P_{N,p}(\lambda)}{M_{N,p}(\lambda)}-1\right|\\
  &+P_{N,p}(\mathsf W_{N,p}^c)+M_{N,p}(\mathsf W_{N,p}^c)
  \longrightarrow0
\end{aligned}
\]
uniformly for \(p\in K\), since the last two terms vanish by
\cref{lem:young-concentration} and Hoeffding's inequality.

By \eqref{eq:young-multinomial-l1} and the \(\ell^1\)-isometry
of \(\mathcal E_{N,p}\), it remains to show
\begin{equation}\label{eq:embedded-multinomial-limit}
  \sup_{p\in K}
  \|\mathcal E_{N,p}M_{N,p}-\varphi_{0,\Sigma_p}\|_1\longrightarrow0.
\end{equation}
Fix
\(\beta_0\in(1/2,2/3)\) and set
\[
  V_{N,p}=\{\lambda\in\mathsf L_N:
  \|\lambda-Np\|_\infty\le N^{\beta_0}\},
  \qquad x_\lambda=\frac{\lambda-Np}{\sqrt N}.
\]
Since \(\inf_{p\in K}p_d>0\),
\(V_{N,p}\subset\operatorname{supp}M_{N,p}\) for all sufficiently large
\(N\), uniformly in \(p\in K\).
Uniformly for \(p\in K\) and \(\lambda\in V_{N,p}\),
\begin{equation}\label{eq:multinomial-local-estimate}
  \frac{N^{(d-1)/2}}{|\mathsf C|}M_{N,p}(\lambda)
  =\varphi_{0,\Sigma_p}(x_\lambda)(1+o(1)).
\end{equation}
To see this, uniform Stirling and third-order Taylor expansions give
\[
  \log M_{N,p}(\lambda)
  =-\frac{d-1}{2}\log(2\pi N)-\frac12\sum_i\log p_i
   -\frac12\sum_i\frac{x_{\lambda,i}^2}{p_i}
   +O_K\!\left(N^{3\beta_0-2}+N^{\beta_0-1}+N^{-1}\right).
\]
This is the Gaussian exponent since
\(\langle x,\Sigma_p^{-1}x\rangle=\sum_i x_i^2/p_i\).
The normalizing constants agree since
\(|\mathsf C|^2=\det[\langle e_i-e_d,e_j-e_d\rangle]_{i,j<d}=d\), while
\(\ker\Sigma_p=\operatorname{span}\{(1,\ldots,1)\}\) and
\(\operatorname{cof}_{ii}(\Sigma_p)=\prod_jp_j\) give
\(\det_{\mathsf H}\Sigma_p=d\prod_jp_j\).
The remainder is
\(o(1)\) since \(1/2<\beta_0<2/3\).

Set
\[
  U_{N,p}=\bigcup_{\lambda\in V_{N,p}}\mathsf C_{N,p}(\lambda).
\]
The multinomial mass outside \(V_{N,p}\) and the Gaussian mass
outside \(U_{N,p}\) are exponentially small: for some
\(c_K,C_K>0\),
\begin{equation}\label{eq:multinomial-gaussian-tail}
  \sup_{p\in K}\left[
  M_{N,p}(V_{N,p}^c)
  +\int_{U_{N,p}^c}\varphi_{0,\Sigma_p}(x)\,\dd x
  \right]
  \le C_K e^{-c_K N^{2\beta_0-1}}.
\end{equation}
Indeed, if \(\mathsf C_{N,p}(\lambda)\) contains \(x\), then
\(\|x_\lambda\|_\infty\le\|x\|_\infty+(d-1)/(2\sqrt N)\).  Thus
\[
  \{x:\|x\|_\infty\le
  N^{\beta_0-1/2}-(d-1)/(2\sqrt N)\}\subset U_{N,p}.
\]
On the set at left \(\|x_\lambda\|_\infty\le N^{\beta_0-1/2}\),
so \(\lambda\in V_{N,p}\).  Hoeffding controls \(M_{N,p}(V_{N,p}^c)\);
the complementary inclusion and the uniform eigenvalue bounds for
\(\Sigma_p\) control the Gaussian tail on \(U_{N,p}^c\).

The Gaussian density has vanishing relative oscillation on
\(\mathsf C_{N,p}(\lambda)\), \(\lambda\in V_{N,p}\).
Indeed, uniformly for \(p\in K\), \(\lambda\in V_{N,p}\), and
\(x,y,z\in\mathsf C_{N,p}(\lambda)\),
\(\|x-y\|=O(N^{-1/2})\) and
\(\|z\|=O(N^{\beta_0-1/2})\).
Since \(\nabla\log\varphi_{0,\Sigma_p}(z)=-\Sigma_p^{-1}z\),
\begin{equation}\label{eq:gaussian-cell-oscillation}
\begin{aligned}
  \left|\log\frac{\varphi_{0,\Sigma_p}(x)}
                       {\varphi_{0,\Sigma_p}(y)}\right|
  &\le \sup_{z\in\mathsf C_{N,p}(\lambda)}
      \|\nabla\log\varphi_{0,\Sigma_p}(z)\|\,\|x-y\|\\
  &\le O_K(N^{\beta_0-1/2})\,O(N^{-1/2})
   =O_K(N^{\beta_0-1}),\\
  \left|\frac{\varphi_{0,\Sigma_p}(x)}
                   {\varphi_{0,\Sigma_p}(y)}-1\right|
  &=O_K(N^{\beta_0-1})=o(1).
\end{aligned}
\end{equation}
The second estimate follows by exponentiation.

The embedded multinomial density is consequently uniformly accurate on
\(U_{N,p}\).  Indeed, \(0\in\mathsf C\), so the point
\(x_\lambda=(\lambda-Np)/\sqrt N\) lies in
\(\mathsf C_{N,p}(\lambda)\).  Uniformly for \(p\in K\),
\(\lambda\in V_{N,p}\), and \(x\) in the same cell,
\begin{equation}\label{eq:embedded-multinomial-cellwise}
\begin{aligned}
  \frac{N^{(d-1)/2}M_{N,p}(\lambda)}
       {|\mathsf C|\varphi_{0,\Sigma_p}(x)}
  &=\frac{N^{(d-1)/2}M_{N,p}(\lambda)}
          {|\mathsf C|\varphi_{0,\Sigma_p}(x_\lambda)}
    \frac{\varphi_{0,\Sigma_p}(x_\lambda)}
         {\varphi_{0,\Sigma_p}(x)}\\
  &=(1+o(1))(1+o(1))=1+o(1).
\end{aligned}
\end{equation}
The two factors are \(1+o(1)\) by the multinomial local
estimate \eqref{eq:multinomial-local-estimate} and the Gaussian
cell-oscillation bound \eqref{eq:gaussian-cell-oscillation}, respectively.

For each \(p\in K\),
\[
\begin{aligned}
  \|\mathcal E_{N,p}M_{N,p}-\varphi_{0,\Sigma_p}\|_1
  ={}&\int_{U_{N,p}}
      |(\mathcal E_{N,p}M_{N,p})(x)-\varphi_{0,\Sigma_p}(x)|\,\dd x\\
     &+\int_{U_{N,p}^c}
      |(\mathcal E_{N,p}M_{N,p})(x)-\varphi_{0,\Sigma_p}(x)|\,\dd x
  \\[-2pt]
  \le{}&M_{N,p}(V_{N,p}^c)
  +\int_{U_{N,p}^c}\varphi_{0,\Sigma_p}(x)\,\dd x\\
  &+\sup_{\substack{\lambda\in V_{N,p}\\
                     x\in\mathsf C_{N,p}(\lambda)}}
  \left|
  \frac{N^{(d-1)/2}M_{N,p}(\lambda)}
       {|\mathsf C|\varphi_{0,\Sigma_p}(x)}-1
  \right|,
\end{aligned}
\]
Here the inequality uses
\(\int_{U_{N,p}}\varphi_{0,\Sigma_p}(x)\,\dd x\le1\) and the cell
normalization
\(\int_{U_{N,p}^c}(\mathcal E_{N,p}M_{N,p})(x)\,\dd x
=M_{N,p}(V_{N,p}^c)\).
Estimates \eqref{eq:multinomial-gaussian-tail} and
\eqref{eq:embedded-multinomial-cellwise} make all three terms on the right
vanish uniformly, proving
\eqref{eq:embedded-multinomial-limit}.

Finally, the \(\ell^1\)-isometry of \(\mathcal E_{N,p}\) gives
\[
\begin{aligned}
  \sup_{p\in K}
  \|\mathcal E_{N,p}P_{N,p}-\varphi_{0,\Sigma_p}\|_1
  \le{}&\sup_{p\in K}\|P_{N,p}-M_{N,p}\|_{\ell^1}\\
  &+\sup_{p\in K}
    \|\mathcal E_{N,p}M_{N,p}-\varphi_{0,\Sigma_p}\|_1
  \longrightarrow0
\end{aligned}
\]
by \eqref{eq:young-multinomial-l1} and
\eqref{eq:embedded-multinomial-limit}.
\end{proof}

\begin{lemma}[Continuous limit of randomized dilation]
\label{lem:randomized-dilation}
Let \(\widetilde Q_{m_n,p}=P_{n,p}D_n\) denote the
\(\widetilde\nu\)-marginal of the joint law obtained by composing the kernel
\(D_n\) with \(P_{n,p}\), before the replacement in
\eqref{eq:rounding-kernel}.  For every compact
\(K\Subset\Delta_d^\circ\),
\begin{align}
  \sup_{p\in K}
  \left\|\mathcal E_{m_n,p}\widetilde Q_{m_n,p}
  -\varphi_{0,\gamma\Sigma_p}\right\|_1&\longrightarrow0,
  \label{eq:rounded-density-limit}\\
  \sup_{p\in K}
  \left|F(\widetilde Q_{m_n,p},P_{m_n,p})-\Fcl(\gamma,d)\right|
  &\longrightarrow0.
  \label{eq:raw-rounded-affinity}
\end{align}
\end{lemma}

\begin{proof}
For \(a>0\) and \(f\in L^1(\mathsf H)\), set
\[
  (\operatorname{dil}_a f)(y)=a^{-(d-1)/2}f(y/\sqrt a),
\]
and define the target-cell averaging operator explicitly by
\[
  (\operatorname{Av}_{m_n,p}f)(x)
  :=\frac{m_n^{(d-1)/2}}{|\mathsf C|}
   \int_{\mathsf C_{m_n,p}(\nu)}f(y)\,\dd y,
  \qquad x\in\mathsf C_{m_n,p}(\nu).
\]
Since \(m_n=\gamma_n n\), the cell definitions and
\eqref{eq:cell-kernel} give
\[
\begin{aligned}
  \gamma_n^{-1/2}\mathsf C_{m_n,p}(\nu)
  &=\frac{\gamma_n^{-1}(\nu+\mathsf C)-np}{\sqrt n},\\
  \left|\gamma_n^{-1/2}\mathsf C_{m_n,p}(\nu)
       \cap\mathsf C_{n,p}(\mu)\right|
  &=\frac{|\mathsf C|}{n^{(d-1)/2}}D_n(\nu\mid\mu).
\end{aligned}
\]
Thus, for \(x\in\mathsf C_{m_n,p}(\nu)\), changing variables
by \(z=y/\sqrt{\gamma_n}\) and summing over the source cells gives
\begin{equation}\label{eq:rounding-intertwining}
\begin{aligned}
 &\bigl[\operatorname{Av}_{m_n,p}\operatorname{dil}_{\gamma_n}
   (\mathcal E_{n,p}P_{n,p})\bigr](x)\\
 &\quad=\frac{m_n^{(d-1)/2}}{|\mathsf C|}
   \sum_{\mu\in\mathsf L_n}\frac{n^{(d-1)/2}}{|\mathsf C|}
   P_{n,p}(\mu)
   \left|\gamma_n^{-1/2}\mathsf C_{m_n,p}(\nu)
        \cap\mathsf C_{n,p}(\mu)\right|\\
 &\quad=\frac{m_n^{(d-1)/2}}{|\mathsf C|}
   \sum_{\mu\in\mathsf L_n}P_{n,p}(\mu)D_n(\nu\mid\mu)
  =(\mathcal E_{m_n,p}\widetilde Q_{m_n,p})(x).
\end{aligned}
\end{equation}
For \(f,g\in L^1(\mathsf H)\),
\[
  \|\operatorname{Av}_{m_n,p}f-\operatorname{Av}_{m_n,p}g\|_1
  \le\|f-g\|_1,
  \qquad
  \|\operatorname{dil}_{\gamma_n}f
      -\operatorname{dil}_{\gamma_n}g\|_1
  =\|f-g\|_1.
\]
By \cref{lem:young-local-limit},
\(\|\mathcal E_{n,p}P_{n,p}-\varphi_{0,\Sigma_p}\|_1=o(1)\) uniformly on
\(K\).  Since
\(\operatorname{dil}_{\gamma_n}\varphi_{0,\Sigma_p}
=\varphi_{0,\gamma_n\Sigma_p}\), dilation isometry gives
\(\|\operatorname{dil}_{\gamma_n}(\mathcal E_{n,p}P_{n,p})
-\varphi_{0,\gamma_n\Sigma_p}\|_1=o(1)\).  Moreover,
\(\|\varphi_{0,\gamma_n\Sigma_p}-\varphi_{0,\gamma\Sigma_p}\|_1=o(1)\)
uniformly because \(\gamma_n\to\gamma\).  Cell averaging satisfies
\[
  \|\operatorname{Av}_{m_n,p}f-f\|_1
  \le 2^{d-1}\!\sup_{\substack{h\in\mathsf H\\
       \|h\|\le\operatorname{diam}(\mathsf C)/\sqrt{m_n}}}
       \|f(\,\cdot+h)-f\|_1.
\]
Since \(\gamma_n\to\gamma\) and \(K\) is compact,
\(\sup_{p\in K,\,n\ge n_0}
\|\nabla\varphi_{0,\gamma_n\Sigma_p}\|_1<\infty\) for some \(n_0\).
Together with
\(\|f(\,\cdot+h)-f\|_1\le\|h\|\,\|\nabla f\|_1\), this gives
\(\|\operatorname{Av}_{m_n,p}\varphi_{0,\gamma_n\Sigma_p}
-\varphi_{0,\gamma_n\Sigma_p}\|_1=o(1)\).  Using the contraction of
\(\operatorname{Av}_{m_n,p}\) and \eqref{eq:rounding-intertwining},
\[
\begin{aligned}
  \|\mathcal E_{m_n,p}\widetilde Q_{m_n,p}
       -\varphi_{0,\gamma\Sigma_p}\|_1
  ={}&\|\operatorname{Av}_{m_n,p}\operatorname{dil}_{\gamma_n}
       (\mathcal E_{n,p}P_{n,p})
       -\varphi_{0,\gamma\Sigma_p}\|_1\\
  \le{}&\|\operatorname{Av}_{m_n,p}
       [\operatorname{dil}_{\gamma_n}(\mathcal E_{n,p}P_{n,p})
       -\varphi_{0,\gamma_n\Sigma_p}]\|_1\\
  &+\|\operatorname{Av}_{m_n,p}\varphi_{0,\gamma_n\Sigma_p}
       -\varphi_{0,\gamma_n\Sigma_p}\|_1
   +\|\varphi_{0,\gamma_n\Sigma_p}
       -\varphi_{0,\gamma\Sigma_p}\|_1\\
  \le{}&\|\operatorname{dil}_{\gamma_n}(\mathcal E_{n,p}P_{n,p})
       -\varphi_{0,\gamma_n\Sigma_p}\|_1\\
  &+\|\operatorname{Av}_{m_n,p}\varphi_{0,\gamma_n\Sigma_p}
       -\varphi_{0,\gamma_n\Sigma_p}\|_1
   +\|\varphi_{0,\gamma_n\Sigma_p}
       -\varphi_{0,\gamma\Sigma_p}\|_1=o(1).
\end{aligned}
\]
The equality is \eqref{eq:rounding-intertwining}.
This proves \eqref{eq:rounded-density-limit}.

Applying \cref{lem:young-local-limit} at sample size \(m_n\) gives
\[
  \mathcal E_{m_n,p}P_{m_n,p}\longrightarrow\varphi_{0,\Sigma_p}
\]
uniformly in \(L^1\).  The embedding preserves Hellinger affinity, and
\eqref{eq:fidelity-continuity} also holds for nonnegative \(L^1\) densities.
Therefore
\[
  F(\widetilde Q_{m_n,p},P_{m_n,p})
  \longrightarrow
  F(\mathsf N_{0,\gamma\Sigma_p},\mathsf N_{0,\Sigma_p})
  =\Fcl(\gamma,d)
\]
uniformly on \(K\).
\end{proof}

\begin{proof}[Proof of \cref{prop:young-rounding}]
For \(i<d\), the nearest-integer rule gives
\[
  |\widetilde\nu_i-\gamma_n\mu_i|
  \le\frac{\gamma_n+1}{2}.
\]
Since the coordinate sums of \(\widetilde\nu\) and \(\gamma_n\mu\) are both
\(m_n\),
\[
  |\widetilde\nu_d-\gamma_n\mu_d|
  \le\frac{(d-1)(\gamma_n+1)}2.
\]
Thus one may take
\[
  c_{\mathrm{rnd}}
  =\frac{d-1}{2}\left(1+\sup_n\gamma_n\right)<\infty.
\]

Compactness of \(K\) gives \(c_0>0\) such that, for all large \(n\), every
\(p\in K\) and \(\mu\in\mathsf T_n(p)\) satisfy
\[
  \mu_d\ge c_0n,
  \qquad
  \mu_i-\mu_{i+1}\ge c_0n\quad(i<d).
\]
For every
\(\widetilde\nu\in\operatorname{supp}D_n(\cdot\mid\mu)\),
\[
  (\widetilde\nu-\mu)_d
  \ge(\gamma_n-1)\mu_d-c_{\mathrm{rnd}},
\]
and
\[
  (\widetilde\nu_i-\mu_i)
  -(\widetilde\nu_{i+1}-\mu_{i+1})
  \ge(\gamma_n-1)(\mu_i-\mu_{i+1})-2c_{\mathrm{rnd}}.
\]
These quantities are positive for all large \(n\), uniformly on \(K\).
Hence \(\widetilde\nu-\mu\) is a partition; since \(\mu\) is a partition,
so is \(\widetilde\nu\).
Moreover, \(\widetilde\nu_d>\mu_d>0\), so
\(\widetilde\nu\ne\nu_n^{\mathrm{fb}}\).  Hence, for
\(\mu\in\mathsf T_n(p)\), the law \(D_n(\cdot\mid\mu)\) is supported on
\(\mathsf Y_{m_n}\), assigns no mass to \(\nu_n^{\mathrm{fb}}\), and
undergoes no replacement in \eqref{eq:rounding-kernel}.  Consequently,
\[
  K_n^{\mathrm{rnd}}(\nu\mid\mu)=D_n(\nu\mid\mu),
  \qquad \nu\in\mathsf Y_{m_n},
\]
and every \(\nu\) in this support satisfies
\(\nu-\mu\vdash m_n-n\) and
\(\|\nu-\gamma_n\mu\|_\infty\le c_{\mathrm{rnd}}\).

Let \(\eps_{n,p}\) be the probability, under \(\mu\sim P_{n,p}\) and
\(\widetilde\nu\sim D_n(\cdot\mid\mu)\), that \(\widetilde\nu\) is
replaced in \eqref{eq:rounding-kernel}.  The preceding argument gives
\[
  \eps_{n,p}\le P_{n,p}(\mathsf T_n(p)^c),
\]
which tends uniformly to zero by \cref{lem:young-concentration}.

By \eqref{eq:raw-rounded-affinity}, the unreplaced output law has affinity
\(\Fcl(\gamma,d)+o(1)\) with \(P_{m_n,p}\), uniformly on \(K\).  Moving the
rejected mass to \(\nu_n^{\mathrm{fb}}\) changes the law by at most
\(2\eps_{n,p}\) in \(\ell^1\).  Applying \eqref{eq:fidelity-continuity} to
classical distributions therefore gives
\eqref{eq:rounding-label-fidelity}.
\end{proof}

\section{Unified Gaussian flat-prior converse}
\label[appendix]{app:gaussian-converse}

This appendix proves the faithful thermal building block used in the two
converse inequalities of \cref{thm:gaussian-amplification} by averaging any
candidate channel \(\Lambda\) over
larger parameter boxes and testing the unavoidable noise in its output.
The fixed-rank theorem uses this argument for all positive-temperature modes
and the explicit vacuum-support compression of
\cref{sec:fixed-rank-LAN} when some \(q_{ij}=0\).
\Cref{lem:weighted-fidelity} turns fidelity into such a test expectation.
\Cref{lem:compact-cp-limit} shows that any bounded sequence of
nearly displacement-covariant maps has a subsequence converging to an exactly
covariant map with the same trace bound.
\Cref{prop:least-noise-moment} shows that every covariant map with
the required gain is at least as noisy as the amplifier, as measured by the
number test, even allowing for lost trace.  \Cref{lem:scalar-convolution}
governs the classical-distribution part of a classical--quantum channel:
translation covariance forces it to scale \(x\) to \(\sqrt\gamma x\), add
independent noise, and possibly lose mass.
\Cref{prop:flat-prior-converse} combines these bounds, giving
\(\Forb(\gamma,p)\) for
\(\Phi_z^{(p)}\mapsto\Psi_z^{(p,\gamma)}\) and
\(\Fcl(\gamma,d)\Forb(\gamma,p)\) for
\(\Theta_{h,z}^{\mathrm{in}}\mapsto\Theta_{h,z}^{\mathrm{tar}}\).

The classical parameter space has
dimension \(k=0\) for the orbital problem and \(k=d-1\) for the hybrid
classical--quantum problem.
All trace-class-valued integrals below are Bochner integrals in trace norm.

\subsection{Weighted fidelity and compact limits}

\begin{lemma}[Weighted root-fidelity bound]\label{lem:weighted-fidelity}
Let \(R,T\) be positive trace-class operators, and let \(W\) be bounded,
positive, and injective.  Define
\[
  \Tr(TW^{-1})
  :=\lim_{\varepsilon\downarrow0}
  \Tr\bigl(T(W+\varepsilon I)^{-1}\bigr)\in[0,\infty].
\]
If this quantity is finite, then
\begin{equation}\label{eq:weighted-quantum-fidelity}
  F(R,T)^2\le\Tr(RW)\Tr(TW^{-1}).
\end{equation}

Let \(\mathsf E\) be a finite-dimensional Euclidean space,
let \(R\in L^1(\mathsf E;\cT_1(\cF_{\mathrm{orb}}))\) be positive, let
\(g\) be a probability density on \(\mathsf E\), and let \(\chi>0\) be
measurable.  For positive \(R,S\) in this \(L^1\)-space, set
\(F(R,S)=\int_{\mathsf E}F(R(y),S(y))\,\dd y\), extending the hybrid-state
convention in \cref{sec:gaussian-LAN}.
The notation \(gT\) means the density \(y\mapsto g(y)T\).
Whenever the right-hand side below is finite,
\begin{equation}\label{eq:weighted-hybrid-fidelity}
  F(R,gT)^2
  \le
  \left(\int_{\mathsf E}\chi(y)\Tr[R(y)W] \,\dd y\right)
  \left(\int_{\mathsf E}\frac{g(y)}{\chi(y)}\,\dd y\right)
  \Tr(TW^{-1}).
\end{equation}
\end{lemma}

\begin{proof}
For \(W_\varepsilon=W+\varepsilon I\), the Schatten H\"older inequality gives
\[
  F(R,T)
  =\|(R^{1/2}W_\varepsilon^{1/2})(W_\varepsilon^{-1/2}T^{1/2})\|_1
  \le\sqrt{\Tr(RW_\varepsilon)}\sqrt{\Tr(TW_\varepsilon^{-1})}.
\]
Squaring and letting \(\varepsilon\downarrow0\), using
\(\Tr(RW_\varepsilon)=\Tr(RW)+\varepsilon\Tr R\) and monotone convergence
for the second trace, proves \eqref{eq:weighted-quantum-fidelity}.  For the
hybrid state,
\[
  F(R,gT)
  =\int_{\mathsf E}\sqrt{g(y)}\,F(R(y),T)\,\dd y.
\]
Apply the first inequality pointwise and then Cauchy--Schwarz to the factors
\(\sqrt{\chi(y)\Tr[R(y)W]}\) and
\(\sqrt{g(y)/\chi(y)}\).  This gives
\eqref{eq:weighted-hybrid-fidelity}.
\end{proof}

\begin{lemma}[Compact-observable limits of completely positive maps]
\label{lem:compact-cp-limit}
Let \(\cH_0,\cH_1\) be separable Hilbert spaces and let
\(\Gamma_j:\cT_1(\cH_0)\to\cT_1(\cH_1)\) be completely positive maps with
\(\sup_j\|\Gamma_j\|_{1\to1}<\infty\).  There are a subsequence and a
completely positive map
\(\Gamma:\cT_1(\cH_0)\to\cT_1(\cH_1)\) such that
\begin{equation}\label{eq:compact-map-convergence}
  \Tr[\Gamma_j(X)K]\longrightarrow\Tr[\Gamma(X)K]
\end{equation}
for every \(X\in\cT_1(\cH_0)\) and every compact
\(K\in\mathcal K(\cH_1)\).  If
\(\Tr\Gamma_j(X)\le c\Tr X\) for every \(X\ge0\), then
\(\Tr\Gamma(X)\le c\Tr X\).

Suppose in addition that
\(U_g\in\U(\cH_0)\) and \(V_g\in\U(\cH_1)\) are families of
unitaries indexed by a set \(G\), and that, for every fixed \(g\in G\) and
\(X\in\cT_1(\cH_0)\),
\[
  \|\Gamma_j(U_g X U_g^*)-V_g\Gamma_j(X)V_g^*\|_1
  \longrightarrow0.
\]
Then
\[
  \Gamma(U_g X U_g^*)=V_g\Gamma(X)V_g^*
\]
for every \(g\in G\) and \(X\in\cT_1(\cH_0)\).
\end{lemma}

\begin{proof}
Let \(M=\sup_j\|\Gamma_j\|_{1\to1}\), and let
\(\mathcal D\subset\mathcal K(\cH_1)\) be a countable norm-dense
\(\mathbb Q(\mathrm i)\)-linear \(*\)-algebra of finite-rank operators.
The predual \(\cT_1(\cH_0)\) is separable, so a diagonal weak-* compactness
argument gives a subsequence for which
\(\Gamma_j^*(K)\) converges weak-* in \(\mathcal B(\cH_0)\) for every
\(K\in\mathcal D\).  The bound
\(\|\Gamma_j^*(K)\|\le M\|K\|\) extends both the limit and the convergence
to every compact \(K\), producing a
bounded complex-linear map
\(\Gamma^*:\mathcal K(\cH_1)\to\mathcal B(\cH_0)\).

To check complete positivity, fix \(r\) and \(A\ge0\) in
\(M_r(\mathcal K(\cH_1))\).  Since \(A^{1/2}\) is compact, choose
\(B_\ell\in M_r(\mathcal D)\) with \(B_\ell\to A^{1/2}\).  Since
\(\Gamma_j^{*(r)}(B_\ell^*B_\ell)\ge0\), first letting \(j\to\infty\) and
then \(\ell\to\infty\) gives \(\Gamma^{*(r)}(A)\ge0\).

Since \(\mathcal K(\cH_1)^*=\cT_1(\cH_1)\), define \(\Gamma(X)\) by
\[
  \Tr[\Gamma(X)K]=\Tr[X\Gamma^*(K)]
  \qquad(K\in\mathcal K(\cH_1)).
\]
Then \(\Gamma\) is completely positive and
\eqref{eq:compact-map-convergence} holds.  If \(P_N\uparrow I\) are
finite-rank projections and \(X\ge0\), then
\[
  \Tr\Gamma(X)
  =\sup_N\lim_j\Tr[\Gamma_j(X)P_N]
  \le c\Tr X.
\]
Finally, for fixed \(g,X\) and compact \(K\), the operator
\(V_g^*KV_g\) is compact, and hence
\[
\begin{aligned}
  &\Tr\!\left[
    \bigl(\Gamma(U_gXU_g^*)-V_g\Gamma(X)V_g^*\bigr)K
  \right]\\
  &\qquad=\lim_j\Tr\!\left[
    \bigl(\Gamma_j(U_gXU_g^*)-V_g\Gamma_j(X)V_g^*\bigr)K
  \right]=0.
\end{aligned}
\]
Compact operators separate trace-class operators, so the limiting map is
covariant.
\end{proof}

\subsection{The least-noise oscillator moment}

Fix \(\gamma>1\).  For \(s\ge1\) and
\(q_1,\ldots,q_s\in(0,1)\), put
\[
  q_\ell^{(\gamma)}=1-\frac{1-q_\ell}{\gamma},
  \qquad
  \tau=\bigotimes_{\ell=1}^s\tau_{q_\ell},
  \qquad
  \tau^{(\gamma)}
  =\bigotimes_{\ell=1}^s\tau_{q_\ell^{(\gamma)}}.
\]

The proof below uses the following two external inputs, isolated here to make
the dependence on \cite{GutaMatsumoto2006MixedGaussianCloning} explicit.

\begin{fact}[Covariant-amplifier representation]
\label{fact:covariant-amplifier-representation}
Every displacement-covariant one-mode channel of amplitude gain
\(\sqrt\gamma\) admits an idler realization
\[
  a_{\mathrm{out}}
  =\sqrt\gamma\,a_{\mathrm{in}}+\sqrt{\gamma-1}\,b^*.
\]
The idler state need not be Gaussian; if the channel is also phase covariant,
it may be chosen number diagonal
\cite[Sec.~II.C and Appendix; Sec.~III for general gain]
{GutaMatsumoto2006MixedGaussianCloning}.
\end{fact}

\begin{fact}[Vacuum-idler stochastic order]
\label{fact:vacuum-idler-stochastic-order}
For a thermal input \(\tau_q\) and a number-diagonal idler, let \(p_k\) be
the output number distribution of the amplifier in
\cref{fact:covariant-amplifier-representation}.  Then
\[
  \sum_{k=0}^m(1-q^{(\gamma)})(q^{(\gamma)})^k
  \ge \sum_{k=0}^m p_k
  \qquad(m\ge0).
\]
Thus the vacuum-idler output \(\tau_{q^{(\gamma)}}\) is stochastically
smallest among these outputs
\cite[Lem.~2.3 and Sec.~III]{GutaMatsumoto2006MixedGaussianCloning}.
\end{fact}

\begin{proposition}[Least-noise moment]
\label{prop:least-noise-moment}
Let \(\Gamma\) be a completely positive trace-nonincreasing map on the
\(s\)-mode trace class satisfying
\[
  \Gamma(D(z)XD(z)^*)
  =D(\sqrt\gamma\,z)\Gamma(X)D(\sqrt\gamma\,z)^*
  \qquad(z\in\mathbb C^s).
\]
If
\[
  W(\widehat{\mathbf N})
  =\prod_{\ell=1}^s w_\ell(\widehat N_\ell),
\]
where every \(w_\ell\) is bounded, nonnegative, and nonincreasing on
\(\mathbb Z_{\ge0}\), then
\begin{equation}\label{eq:least-noise-moment}
  \Tr[\Gamma(\tau)W(\widehat{\mathbf N})]
  \le\Tr[\Gamma(\tau)]
       \Tr[\tau^{(\gamma)}W(\widehat{\mathbf N})].
\end{equation}
\end{proposition}

\begin{proof}
Averaging \(\Gamma\) over independent phase rotations preserves its
trace bound and displacement covariance, and it does not change either side
of \eqref{eq:least-noise-moment}.  We may therefore assume that \(\Gamma\)
is phase covariant in every mode.

Consider one mode.  Write \(\Tr\Gamma(X)=\Tr(AX)\) with \(0\le A\le I\).
Displacement covariance implies \(D(z)^*AD(z)=A\); hence irreducibility of
the Weyl representation gives \(A=cI\), where
  \(c=\Tr\Gamma(\tau_q)\).  If \(c>0\), then
  \(\Lambda=c^{-1}\Gamma\) is a phase-covariant amplifier of amplitude gain
  \(\sqrt\gamma\).  By
  \cref{fact:covariant-amplifier-representation}, it has a number-diagonal
  idler realization.  \Cref{fact:vacuum-idler-stochastic-order} therefore
  gives
\begin{equation}\label{eq:imported-stochastic-order}
  \sum_{n=0}^m(1-q^{(\gamma)})(q^{(\gamma)})^n
  \ge
  \sum_{n=0}^m\langle n|\Lambda(\tau_q)|n\rangle
  \qquad(m\ge0).
\end{equation}
The distribution on the left is the vacuum-idler output, namely that of
\(\tau_{q^{(\gamma)}}\), and both output states are number diagonal.  For a
bounded nonincreasing \(w\), summation by parts writes its expectation as a
nonnegative linear combination of these cumulative probabilities.  Both
distributions in \eqref{eq:imported-stochastic-order} have total mass one,
so the common limiting constant cancels.  Hence
\[
  \Tr[\Lambda(\tau_q)w(\widehat N)]
  \le\Tr[\tau_{q^{(\gamma)}}w(\widehat N)].
\]
Multiplication by \(c\) proves the one-mode case; the case \(c=0\) is
immediate.

For the multimode case, the claim is immediate if some \(w_\ell\)
vanishes identically; otherwise, divide each \(w_\ell\) by
\(w_\ell(0)\) and assume \(0\le w_\ell\le1\).  Put
\[
  W_{\ge\ell}
  =I_1\ot\cdots\ot I_{\ell-1}\ot
   \bigotimes_{j=\ell}^s w_j(\widehat N_j),
  \qquad W_{\ge s+1}=I,
\]
and, for an operator \(X\) on mode \(\ell\), define
\[
  \Gamma_\ell(X)
  =\Tr_{\ne\ell}\!\left[
    W_{\ge\ell+1}^{1/2}
    \Gamma(\tau_{q_1}\ot\cdots\ot\tau_{q_{\ell-1}}\ot X\ot
           \tau_{q_{\ell+1}}\ot\cdots\ot\tau_{q_s})
    W_{\ge\ell+1}^{1/2}
  \right].
\]
Because \(0\le W_{\ge\ell+1}\le I\) and this operator is the identity on
mode \(\ell\), each \(\Gamma_\ell\) is completely positive, trace
nonincreasing, and displacement covariant in that mode.  The phase-covariant
reduction following \eqref{eq:least-noise-moment} also makes \(\Gamma_\ell\)
phase covariant because the fixed thermal inputs and \(W_{\ge\ell+1}\) are
phase invariant.  Moreover,
\[
  \Tr[\Gamma_\ell(\tau_{q_\ell})w_\ell(\widehat N_\ell)]
  =\Tr[\Gamma(\tau)W_{\ge\ell}],
  \qquad
  \Tr\Gamma_\ell(\tau_{q_\ell})
  =\Tr[\Gamma(\tau)W_{\ge\ell+1}].
\]
Applying the one-mode inequality to \(\Gamma_\ell\) and iterating over
\(\ell=1,\ldots,s\) gives \eqref{eq:least-noise-moment}.
\end{proof}

Specialize to \(s=r_d\), \(q_\ell=q_{ij}\),
\(\tau=\tau_p\), and \(\tau^{(\gamma)}=\tau_p^{(\gamma)}\).
For
\(\mathbf k=(k_{ij})_{i<j}\in\mathbb Z_{\ge0}^{r_d}\), let
\(t_{\mathbf k}\) and \(t_{\mathbf k}^{(\gamma)}\) be the corresponding
number-basis eigenvalues, and define
\begin{equation}\label{eq:quantum-witness}
  W_\gamma
  =\sum_{\mathbf k}
  \sqrt{\frac{t_{\mathbf k}}{t_{\mathbf k}^{(\gamma)}}}
  \proj{\mathbf k}.
\end{equation}
Equivalently,
\[
  W_\gamma=\prod_{i<j}w_{ij}(\widehat N_{ij}),
  \qquad
  w_{ij}(n)=
  \sqrt{\frac{1-q_{ij}}{1-q_{ij}^{(\gamma)}}}
  \left(\frac{q_{ij}}{q_{ij}^{(\gamma)}}\right)^{n/2}.
\]
Here \(\widehat N_{ij}\) is the number operator of the mode indexed by
\(i<j\), and \(\lvert\mathbf k\rvert:=\sum_{i<j}k_{ij}\) is the total
occupation number.
Since \(q_{ij}<q_{ij}^{(\gamma)}\), each \(w_{ij}(n)\) is positive, bounded,
and strictly decreasing in \(n\), and the eigenvalues of \(W_\gamma\) tend
to zero as \(\lvert\mathbf k\rvert\to\infty\).
Thus \(W_\gamma\) is bounded,
injective, and compact.  Its definition gives
\begin{equation}\label{eq:quantum-witness-identities}
  \Tr(\tau_p^{(\gamma)}W_\gamma)
  =\Tr(\tau_p W_\gamma^{-1})
  =\sum_{\mathbf k}\sqrt{t_{\mathbf k}t_{\mathbf k}^{(\gamma)}}
  =\Forb(\gamma,p),
\end{equation}
where the inverse is understood as in \cref{lem:weighted-fidelity}.

\subsection{Classical trace kernels}

\begin{lemma}[Translation-covariant scalar kernels]
\label{lem:scalar-convolution}
Let \(\mathsf E\) be a finite-dimensional Euclidean space and let
\(\mathcal R:L^1(\mathsf E)\to C_0(\mathsf E)^*\) be a positive
contraction.  For a finite Radon measure \(\lambda\), define
\(T_a\lambda(A)=\lambda(A-a)\) for Borel \(A\).  Suppose
\begin{equation}\label{eq:scalar-measure-covariance}
  \mathcal R(u(\,\cdot-h))=T_{\sqrt\gamma h}\mathcal R(u)
  \qquad(u\in L^1(\mathsf E),\ h\in\mathsf E)
\end{equation}
as an equality of finite Radon measures.  Then there is a finite positive
measure \(\mu\), with \(\mu(\mathsf E)\le1\), such that
\begin{equation}\label{eq:scalar-convolution}
  \mathcal R(u)(A)
  =\int_{\mathsf E}u(x)\,\mu(A-\sqrt\gamma x)\,\dd x
\end{equation}
for every \(u\in L^1(\mathsf E)\) and every Borel set \(A\).
\end{lemma}

\begin{proof}
If \(\mathsf E=\{0\}\), take \(\mu=\mathcal R(1)\).  Otherwise choose a
nonnegative approximate identity \(u_j\in C_c(\mathsf E)\) with \(\int u_j=1\)
and \(\operatorname{supp}u_j\subset\{x:\|x\|\le j^{-1}\}\).  The positive
measures \(\mathcal R(u_j)\) lie in the unit ball of \(C_0(\mathsf E)^*\), so
weak-* compactness and separability yield a positive subsequential weak-* limit
\(\mu\) with \(\mu(\mathsf E)\le1\).

For \(u\in L^1(\mathsf E)\), the identity
\(u*u_j=\int_{\mathsf E}u(x)u_j(\,\cdot-x)\,\dd x\) holds as a Bochner
integral in \(L^1(\mathsf E)\), so bounded linearity and covariance give
\[
  \mathcal R(u*u_j)=\int_{\mathsf E}u(x)T_{\sqrt\gamma x}\mathcal R(u_j)\,\dd x.
\]
Pair with \(\phi\in C_0(\mathsf E)\).  The approximate-identity property and
contractivity give \(\mathcal R(u*u_j)\to\mathcal R(u)\) in total variation.
On the right, weak-* convergence applies pointwise because
\(y\mapsto\phi(y+\sqrt\gamma x)\) lies in \(C_0(\mathsf E)\), and the
integrand is bounded by \(|u(x)|\|\phi\|_\infty\).  Dominated convergence gives
\[
  \langle\phi,\mathcal R(u)\rangle=\int_{\mathsf E}u(x)\langle\phi,T_{\sqrt\gamma x}\mu\rangle\,\dd x.
\]
By definition of \(T_a\), the right-hand side pairs \(\phi\) with the finite
Radon measure in \eqref{eq:scalar-convolution}, whose total variation is at
most \(\|u\|_1\mu(\mathsf E)\).  Riesz uniqueness proves the claim.
\end{proof}

\subsection{One flat-prior converse for both Gaussian models}

\begin{proposition}[Flat-prior converse inequality]
\label{prop:flat-prior-converse}
Let \(k\in\{0,d-1\}\).  For \(k=0\), let
\(\mathsf E=\{0\}\), omit the classical factor, and interpret \(h\)-integration
as evaluation at \(0\).  For \(k=d-1\), let \(\mathsf E=\mathsf H\) with its
induced Euclidean measure.  Let
\[
  \mathsf E_L=\{h\in\mathsf E:\|h\|_\infty\le L\},
  \qquad
  \Xi_L=\mathsf E_L\times\mathsf Z_L,
\]
with \(|\mathsf E_L|=1\) when \(k=0\).  Define
\[
  R_{0,z}^{(0)}=\Phi_z^{(p)},
  \qquad S_{0,z}^{(0)}=\Psi_z^{(p,\gamma)},
\]
\[
  R_{h,z}^{(d-1)}=\Theta_{h,z}^{\mathrm{in}},
  \qquad S_{h,z}^{(d-1)}=\Theta_{h,z}^{\mathrm{tar}},
  \qquad
  c_k=\left(\frac{2\sqrt\gamma}{1+\gamma}\right)^{k/2}.
\]
For \(k=0\), the supremum below is over quantum channels on
\(\cF_{\mathrm{orb}}\); for \(k=d-1\), it is over classical--quantum
channels.  Then
\begin{equation}\label{eq:unified-flat-prior-converse}
  \limsup_{L\to\infty}\sup_\Lambda
  \frac1{|\Xi_L|}\int_{\Xi_L}
  F\!\left(\Lambda(R_{h,z}^{(k)}),S_{h,z}^{(k)}\right)
  \,\dd h\,\dd z
  \le c_k\Forb(\gamma,p).
\end{equation}
\end{proposition}

\begin{proof}
For \(k=0\), set \(f_1=f_\gamma=\chi=1\) on the one-point space.  For
\(k=d-1\), let \(f_1\) and \(f_\gamma\) be the centered Gaussian densities
with covariances \(\Sigma_p\) and \(\gamma\Sigma_p\), and set
\[
  \chi(y)=\sqrt{\frac{f_1(y)}{f_\gamma(y)}}
  =\gamma^{k/4}
   \exp\!\left[-\frac{\gamma-1}{4\gamma}
     \langle y,\Sigma_p^{-1}y\rangle\right].
\]
This function is bounded and strictly positive; for \(k=d-1\), it also
belongs to \(C_0(\mathsf E)\).  In both cases,
\begin{equation}\label{eq:classical-witness-identity}
  \int_{\mathsf E}\frac{f_1(y)}{\chi(y)}\,\dd y
  =\int_{\mathsf E}\sqrt{f_1(y)f_\gamma(y)}\,\dd y
  =c_k.
\end{equation}

For \(\xi=(h,z)\), define
\[
  (\mathcal U_\xi^{\mathrm{in}}R)(x)
  =D(z)R(x-h)D(z)^*,
  \qquad
  (\mathcal U_\xi^{\mathrm{out}}R)(y)
  =D(\sqrt\gamma z)R(y-\sqrt\gamma h)D(\sqrt\gamma z)^*.
\]
Then
\(R_\xi^{(k)}=\mathcal U_\xi^{\mathrm{in}}(f_1\tau_p)\) and
\(S_\xi^{(k)}=\mathcal U_\xi^{\mathrm{out}}(f_1\tau_p)\).  For a channel
\(\Lambda\), let
\[
  \overline\Lambda_L
  =\frac1{|\Xi_L|}\int_{\Xi_L}
   \mathcal U_{-\eta}^{\mathrm{out}}\circ\Lambda\circ
   \mathcal U_\eta^{\mathrm{in}}\,\dd\eta.
\]
Invariance of \(F\) under the isometries
\(\mathcal U_\xi^{\mathrm{out}}\) and joint concavity, applied to the Bochner
average, give
\begin{equation}\label{eq:unified-folner-concavity}
  F_L(\Lambda)
  =\frac1{|\Xi_L|}\int_{\Xi_L}
  F\!\left(\Lambda(R_\xi^{(k)}),S_\xi^{(k)}\right)\,\dd\xi
  \le F(\overline\Lambda_L(f_1\tau_p),f_1\tau_p).
\end{equation}

Let \(F_*=\limsup_{L\to\infty}\sup_\Lambda F_L(\Lambda)\).  Choose
\(L_j\to\infty\) with \(\sup_\Lambda F_{L_j}(\Lambda)\to F_*\), and channels
\(\widetilde\Lambda_j\) satisfying
\(F_{L_j}(\widetilde\Lambda_j)\ge\sup_\Lambda F_{L_j}(\Lambda)-j^{-1}\).
Set
\[
  \Lambda_j=\overline{\widetilde\Lambda_j}_{L_j},
  \qquad R_j=\Lambda_j(f_1\tau_p).
\]
Then
\begin{equation}\label{eq:unified-maximizing-sequence}
  F_{L_j}(\widetilde\Lambda_j)\longrightarrow F_*,
  \qquad
  F_{L_j}(\widetilde\Lambda_j)\le F(R_j,f_1\tau_p).
\end{equation}
For \(\eta\in\mathsf E\times\mathbb C^{r_d}\), write
\[
  A_{j,\eta}
  =\mathcal U_{-\eta}^{\mathrm{out}}\circ
   \widetilde\Lambda_j\circ\mathcal U_\eta^{\mathrm{in}},
  \qquad
  \Lambda_j=\frac1{|\Xi_{L_j}|}\int_{\Xi_{L_j}}A_{j,\eta}\,\dd\eta.
\]
The group law and the change of variables \(u=\eta+\xi\) give
\[
\begin{aligned}
  \Lambda_j(\mathcal U_\xi^{\mathrm{in}}R)
  &=\mathcal U_\xi^{\mathrm{out}}
    \frac1{|\Xi_{L_j}|}\int_{\Xi_{L_j}+\xi}A_{j,u}(R)\,\dd u,\\
  \mathcal U_\xi^{\mathrm{out}}\Lambda_j(R)
  &=\mathcal U_\xi^{\mathrm{out}}
    \frac1{|\Xi_{L_j}|}\int_{\Xi_{L_j}}A_{j,u}(R)\,\dd u.
\end{aligned}
\]
The two integrals cancel on the intersection of the two averaging sets,
and \(\|A_{j,u}\|_{1\to1}\le1\).  Therefore, for fixed \(\xi\),
\begin{equation}\label{eq:unified-approx-covariance}
  \|\Lambda_j(\mathcal U_\xi^{\mathrm{in}}R)
    -\mathcal U_\xi^{\mathrm{out}}\Lambda_j(R)\|_1
  \le
  \frac{|(\Xi_{L_j}+\xi)\mathbin\triangle\Xi_{L_j}|}{|\Xi_{L_j}|}
  \|R\|_1
  \longrightarrow0.
\end{equation}
Indeed, \(\Xi_1\) is a bounded polytope and \(\Xi_L=L\Xi_1\), so the
ratio is \(O_\xi(L^{-1})\).

Define
\[
  \Gamma_j^\chi(X)
  =\int_{\mathsf E}\chi(y)[\Lambda_j(f_1X)](y)\,\dd y.
\]
These maps are completely positive and satisfy
\(\|\Gamma_j^\chi\|_{1\to1}\le\|\chi\|_\infty\).  Taking
\(\xi=(0,z)\) in \eqref{eq:unified-approx-covariance}, multiplying by
\(\chi\), and integrating gives
\[
\begin{aligned}
 &\|\Gamma_j^\chi(D(z)XD(z)^*)
   -D(\sqrt\gamma z)\Gamma_j^\chi(X)D(\sqrt\gamma z)^*\|_1\\
 &\qquad\le\|\chi\|_\infty
   \frac{|(\Xi_{L_j}+(0,z))\mathbin\triangle\Xi_{L_j}|}{|\Xi_{L_j}|}
   \|X\|_1\longrightarrow0.
\end{aligned}
\]
Hence \cref{lem:compact-cp-limit} gives, after passing to a subsequence, a
completely positive displacement-covariant limit \(\Gamma^\chi\) such that
\begin{equation}\label{eq:weighted-map-limit}
  \Tr[\Gamma_j^\chi(X)K]
  \longrightarrow\Tr[\Gamma^\chi(X)K]
\end{equation}
for trace-class \(X\) and compact \(K\).  With
\[
  M_j
  =\int_{\mathsf E}\chi(y)\Tr[R_j(y)W_\gamma] \,\dd y
  =\Tr[\Gamma_j^\chi(\tau_p)W_\gamma],
\]
compactness of \(W_\gamma\) gives
\begin{equation}\label{eq:moment-limit}
  M_j\longrightarrow\Tr[\Gamma^\chi(\tau_p)W_\gamma].
\end{equation}
Moreover, \cref{lem:weighted-fidelity},
\eqref{eq:classical-witness-identity}, and
\eqref{eq:quantum-witness-identities} give
\begin{equation}\label{eq:unified-fidelity-to-moment}
  F(R_j,f_1\tau_p)^2
  \le M_j\,c_k\Forb(\gamma,p).
\end{equation}

If \(k=0\), then \(\Gamma_j^\chi=\Lambda_j\) are channels.  The trace-bound
clause of \cref{lem:compact-cp-limit} with \(c=1\) makes the limit trace
nonincreasing.  By
\cref{prop:least-noise-moment},
\[
\begin{aligned}
  \lim_j M_j
  &=\Tr[\Gamma^\chi(\tau_p)W_\gamma]\\
  &\le\Tr\Gamma^\chi(\tau_p)\,
       \Tr(\tau_p^{(\gamma)}W_\gamma)
  \le\Forb(\gamma,p).
\end{aligned}
\]
Combining \eqref{eq:unified-maximizing-sequence} with
\eqref{eq:unified-fidelity-to-moment} yields
\(F_*\le\Forb(\gamma,p)\).

Assume henceforth that \(k=d-1\).  For \(u\in L^1(\mathsf E)\), define
\[
  \mathcal R_j(u)(A)
  =\int_A\Tr\bigl([\Lambda_j(u\tau_p)](y)\bigr)\,\dd y.
\]
The maps \(\mathcal R_j:L^1(\mathsf E)\to C_0(\mathsf E)^*\) are positive
contractions.  Using countable dense subsets of \(L^1(\mathsf E)\) and
\(C_0(\mathsf E)\), pass to a further subsequence converging weak-* on every \(u\)
to a positive contraction \(\mathcal R\).  Taking \(z=0\) in
\eqref{eq:unified-approx-covariance}, pairing with a function in
\(C_0(\mathsf E)\), and passing to the limit gives
\[
  \mathcal R(u(\,\cdot-h))
  =T_{\sqrt\gamma h}\mathcal R(u).
\]
Let \(\mu\) be the measure from \cref{lem:scalar-convolution}; then
\(\mu(\mathsf E)\le1\).

Finite-rank projections, \eqref{eq:weighted-map-limit}, the weak-*
convergence of \(\mathcal R_j(f_1)\), and
\(\chi\in C_0(\mathsf E)\) give
\[
\begin{aligned}
  \Tr\Gamma^\chi(\tau_p)
  &\le\liminf_j\Tr\Gamma_j^\chi(\tau_p)\\
  &=\lim_j\int_{\mathsf E}\chi(y)\,\mathcal R_j(f_1)(\dd y)
   =\int_{\mathsf E}\chi(y)\,\mathcal R(f_1)(\dd y).
\end{aligned}
\]
By \eqref{eq:scalar-convolution}, the last integral equals
\[
  \int_{\mathsf E}\mu(\dd r)
  \int_{\mathsf E}f_1(x)\chi(r+\sqrt\gamma x)\,\dd x.
\]
Completing the square gives
\[
  \int_{\mathsf E}f_1(x)\chi(r+\sqrt\gamma x)\,\dd x
  =c_k\exp\!\left[-\frac{\gamma-1}{2\gamma(\gamma+1)}
    \langle r,\Sigma_p^{-1}r\rangle\right].
\]
Since \(\mu(\mathsf E)\le1\), it follows that
\(\Tr\Gamma^\chi(\tau_p)\le c_k\).

Finally write \(\Tr\Gamma^\chi(X)=\Tr(AX)\).  Covariance implies that \(A\)
commutes with every Weyl displacement, so \(A=cI\), with
\[
  c=\Tr\Gamma^\chi(\tau_p)\le c_k\le1.
\]
Thus \(\Gamma^\chi\) is trace nonincreasing, and
\cref{prop:least-noise-moment} gives
\[
\begin{aligned}
  \lim_j M_j
  &=\Tr[\Gamma^\chi(\tau_p)W_\gamma]\\
  &\le\Tr\Gamma^\chi(\tau_p)\,
       \Tr(\tau_p^{(\gamma)}W_\gamma)
  \le c_k\Forb(\gamma,p).
\end{aligned}
\]
Combining this with \eqref{eq:unified-maximizing-sequence} and
\eqref{eq:unified-fidelity-to-moment} yields
\[
  F_*\le c_k\Forb(\gamma,p),
\]
which proves the hybrid classical--quantum case.
\end{proof}


\begin{thebibliography}{BCFJS96}
\bibitem[WZ82]{WoottersZurek1982}
W. K. Wootters and W. H. Zurek,
A single quantum cannot be cloned,
\doi{10.1038/299802a0}{\emph{Nature} \textbf{299}, 802--803 (1982)}.

\bibitem[Die82]{Dieks1982}
D. Dieks,
Communication by EPR devices,
\doi{10.1016/0375-9601(82)90084-6}{\emph{Physics Letters A} \textbf{92}, 271--272 (1982)}.

\bibitem[BH96]{BuzekHillery1996}
V. Buzek and M. Hillery,
Quantum copying: Beyond the no-cloning theorem,
\doi{10.1103/PhysRevA.54.1844}{\emph{Phys. Rev. A} \textbf{54}, 1844--1852 (1996)}.

\bibitem[GM97]{GisinMassar1997}
N. Gisin and S. Massar,
Optimal quantum cloning machines,
\doi{10.1103/PhysRevLett.79.2153}{\emph{Phys. Rev. Lett.} \textbf{79}, 2153--2156 (1997)}.

\bibitem[BDE+98]{BrussDiVincenzoEkertFuchsMacchiavelloSmolin1998}
D. Bruss, D. P. DiVincenzo, A. Ekert, C. A. Fuchs, C. Macchiavello, and J. A. Smolin,
Optimal universal and state-dependent quantum cloning,
\doi{10.1103/PhysRevA.57.2368}{\emph{Phys. Rev. A} \textbf{57}, 2368--2378 (1998)},
\arxiv{quant-ph/9705038}.

\bibitem[CB99]{CheflesBarnett1999}
A. Chefles and S. M. Barnett,
Strategies and networks for state-dependent quantum cloning,
\doi{10.1103/PhysRevA.60.136}{\emph{Phys. Rev. A} \textbf{60}, 136--144 (1999)},
\arxiv{quant-ph/9812035}.

\bibitem[Wer98]{Werner1998}
R. F. Werner,
Optimal cloning of pure states,
\doi{10.1103/PhysRevA.58.1827}{\emph{Phys. Rev. A} \textbf{58}, 1827--1832 (1998)},
\arxiv{quant-ph/9804001}.

\bibitem[KW99]{KeylWerner1999}
M. Keyl and R. F. Werner,
Optimal cloning of pure states, testing single clones,
\doi{10.1063/1.532887}{\emph{J. Math. Phys.} \textbf{40}, 3283--3299 (1999)},
\arxiv{quant-ph/9807010}.

\bibitem[MP95]{MassarPopescu1995}
S. Massar and S. Popescu,
Optimal extraction of information from finite quantum ensembles,
\doi{10.1103/PhysRevLett.74.1259}{\emph{Phys. Rev. Lett.} \textbf{74}, 1259--1263 (1995)}.

\bibitem[BEM98]{BrussEkertMacchiavello1998}
D. Bruss, A. Ekert, and C. Macchiavello,
Optimal universal quantum cloning and state estimation,
\doi{10.1103/PhysRevLett.81.2598}{\emph{Phys. Rev. Lett.} \textbf{81}, 2598--2601 (1998)}.

\bibitem[BA06]{BaeAcin2006}
J. Bae and A. Acin,
Asymptotic quantum cloning is state estimation,
\doi{10.1103/PhysRevLett.97.030402}{\emph{Phys. Rev. Lett.} \textbf{97}, 030402 (2006)},
\arxiv{quant-ph/0603078}.

\bibitem[CD06]{ChiribellaDAriano2006}
G. Chiribella and G. M. D'Ariano,
Quantum information becomes classical when distributed to many users,
\doi{10.1103/PhysRevLett.97.250503}{\emph{Phys. Rev. Lett.} \textbf{97}, 250503 (2006)}.

\bibitem[CY14]{ChiribellaYang2014Optimal}
G. Chiribella and Y. Yang,
Optimal asymptotic cloning machines,
\doi{10.1088/1367-2630/16/6/063005}{\emph{New J. Phys.} \textbf{16}, 063005 (2014)},
\arxiv{1404.0990}.

\bibitem[KC22]{KimChitambar2022}
C. Kim and E. Chitambar,
Process-optimized phase-covariant quantum cloning,
\doi{10.1103/PhysRevA.106.022405}{\emph{Phys. Rev. A} \textbf{106}, 022405 (2022)},
\arxiv{2107.03042}.

\bibitem[CKNWW05]{CerfKrugerNavezWernerWolf2005}
N. J. Cerf, O. Kr\"uger, P. Navez, R. F. Werner, and M. M. Wolf,
Non-Gaussian cloning of quantum coherent states is optimal,
\doi{10.1103/PhysRevLett.95.070501}{\emph{Phys. Rev. Lett.} \textbf{95}, 070501 (2005)},
\arxiv{quant-ph/0410058}.

\bibitem[DG98]{DuanGuo1998}
L.-M. Duan and G.-C. Guo,
Probabilistic cloning and identification of linearly independent quantum states,
\doi{10.1103/PhysRevLett.80.4999}{\emph{Phys. Rev. Lett.} \textbf{80}, 4999--5002 (1998)}.

\bibitem[NG98]{NiuGriffiths1998}
C.-S. Niu and R. B. Griffiths,
Optimal copying of one quantum bit,
\doi{10.1103/PhysRevA.58.4377}{\emph{Phys. Rev. A} \textbf{58}, 4377--4393 (1998)},
\arxiv{quant-ph/9805063}.

\bibitem[BCDM00]{BrussCinchettiDArianoMacchiavello2000}
D. Bruss, M. Cinchetti, G. M. D'Ariano, and C. Macchiavello,
Phase-covariant quantum cloning,
\doi{10.1103/PhysRevA.62.012302}{\emph{Phys. Rev. A} \textbf{62}, 012302 (2000)},
\arxiv{quant-ph/9909046}.

\bibitem[Cer00]{Cerf2000Asymmetric}
N. J. Cerf,
Asymmetric quantum cloning in any dimension,
\doi{10.1080/09500340008244036}{\emph{J. Mod. Opt.} \textbf{47}, 187--209 (2000)},
\arxiv{quant-ph/9805024}.

\bibitem[MJPV99]{MuraoJonathanPlenioVedral1999}
M. Murao, D. Jonathan, M. B. Plenio, and V. Vedral,
Quantum telecloning and multiparticle entanglement,
\doi{10.1103/PhysRevA.59.156}{\emph{Phys. Rev. A} \textbf{59}, 156--161 (1999)},
\arxiv{quant-ph/9806082}.

\bibitem[CDP08]{ChiribellaDArianoPerinotti2008}
G. Chiribella, G. M. D'Ariano, and P. Perinotti,
Optimal cloning of unitary transformation,
\doi{10.1103/PhysRevLett.101.180504}{\emph{Phys. Rev. Lett.} \textbf{101}, 180504 (2008)},
\arxiv{0804.0129}.

\bibitem[CYY13]{ChiribellaYangYao2013}
G. Chiribella, Y. Yang, and A. C.-C. Yao,
Quantum replication at the Heisenberg limit,
\doi{10.1038/ncomms3915}{\emph{Nat. Commun.} \textbf{4}, 2915 (2013)},
\arxiv{1304.2910}.

\bibitem[GCMBC14]{GendraCalsamigliaMunozTapiaBaganChiribella2014}
B. Gendra, J. Calsamiglia, R. Mu\~noz-Tapia, E. Bagan, and G. Chiribella,
Probabilistic metrology attains macroscopic cloning of quantum clocks,
\doi{10.1103/PhysRevLett.113.260402}{\emph{Phys. Rev. Lett.} \textbf{113}, 260402 (2014)},
\arxiv{1406.2218}.

\bibitem[DSS15]{DurSekatskiSkotiniotis2015}
W. D\"ur, P. Sekatski, and M. Skotiniotis,
Deterministic superreplication of one-parameter unitary transformations,
\doi{10.1103/PhysRevLett.114.120503}{\emph{Phys. Rev. Lett.} \textbf{114}, 120503 (2015)},
\arxiv{1410.6008}.

\bibitem[CYH15]{ChiribellaYangHuang2015}
G. Chiribella, Y. Yang, and C. Huang,
Universal superreplication of unitary gates,
\doi{10.1103/PhysRevLett.114.120504}{\emph{Phys. Rev. Lett.} \textbf{114}, 120504 (2015)}.

\bibitem[CY16]{ChiribellaYang2016}
G. Chiribella and Y. Yang,
Quantum superreplication of states and gates,
\doi{10.1007/s11467-016-0556-7}{\emph{Front. Phys.} \textbf{11}, 110304 (2016)},
\arxiv{1512.05839}.

\bibitem[SSD15]{SekatskiSkotiniotisDur2015NoSignaling}
P. Sekatski, M. Skotiniotis, and W. D\"ur,
No-signaling bounds for quantum cloning and metrology,
\doi{10.1103/PhysRevA.92.022355}{\emph{Phys. Rev. A} \textbf{92}, 022355 (2015)},
\arxiv{1412.3361}.

\bibitem[HT24]{HuTomamichel2024}
Y. Hu and M. Tomamichel,
Fundamental limits on quantum cloning from the no-signaling principle,
\doi{10.1103/PhysRevA.109.022221}{\emph{Phys. Rev. A} \textbf{109}, 022221 (2024)},
\arxiv{2305.02002}.

\bibitem[SGKS25]{SekatskiGuryanovaKothakondaSkotiniotis2025}
P. Sekatski, Y. Guryanova, N. B. T. Kothakonda, and M. Skotiniotis,
Cloning quantum channels,
\arxiv{2509.08059} (2025).

\bibitem[BGSQ26]{BrzicGrinkoStudzinskiQuintino2026}
V. Brzi\'c, D. Grinko, M. Studzi\'nski, and M. T. Quintino,
Optimal pure state cloning and transposition are complementary channels,
\arxiv{2603.23628} (2026).

\bibitem[SIGA05]{ScaraniIblisdirGisinAcin2005}
V. Scarani, S. Iblisdir, N. Gisin, and A. Acin,
Quantum cloning,
\doi{10.1103/RevModPhys.77.1225}{\emph{Rev. Mod. Phys.} \textbf{77}, 1225--1256 (2005)},
\arxiv{quant-ph/0511088}.

\bibitem[FWJ+14]{FanWangJingYueShiZhangMu2014}
H. Fan, Y.-N. Wang, L. Jing, J.-D. Yue, H.-D. Shi, Y.-L. Zhang, and L.-Z. Mu,
Quantum cloning machines and the applications,
\doi{10.1016/j.physrep.2014.06.004}{\emph{Phys. Rep.} \textbf{544}, 241--322 (2014)},
\arxiv{1301.2956}.

\bibitem[BCFJS96]{BarnumCavesFuchsJozsaSchumacher1996}
H. Barnum, C. M. Caves, C. A. Fuchs, R. Jozsa, and B. Schumacher,
Noncommuting mixed states cannot be broadcast,
\doi{10.1103/PhysRevLett.76.2818}{\emph{Phys. Rev. Lett.} \textbf{76}, 2818--2821 (1996)},
\arxiv{quant-ph/9511010}.

\bibitem[KH08]{KalevHen2008}
A. Kalev and I. Hen,
No-broadcasting theorem and its classical counterpart,
\doi{10.1103/PhysRevLett.100.210502}{\emph{Phys. Rev. Lett.} \textbf{100},
210502 (2008)}.

\bibitem[CEM99]{CiracEkertMacchiavello1999}
J. I. Cirac, A. K. Ekert, and C. Macchiavello,
Optimal purification of single qubits,
\doi{10.1103/PhysRevLett.82.4344}{\emph{Phys. Rev. Lett.} \textbf{82}, 4344--4347 (1999)},
\arxiv{quant-ph/9812075}.

\bibitem[Ras03a]{Rastegin2003Upper}
A. E. Rastegin,
Upper bound on the global fidelity for mixed-state cloning,
\doi{10.1103/PhysRevA.67.012305}{\emph{Phys. Rev. A} \textbf{67}, 012305 (2003)}.

\bibitem[Ras03b]{Rastegin2002Relative}
A. E. Rastegin,
A lower bound on the relative error of mixed-state cloning and related operations,
\emph{J. Opt. B: Quantum Semiclass. Opt.} \textbf{5}, S647--S650 (2003),
\arxiv{quant-ph/0208159}.

\bibitem[Ras03c]{Rastegin2003StateDependent}
A. E. Rastegin,
Global-fidelity limits of state-dependent cloning of mixed states,
\doi{10.1103/PhysRevA.68.032303}{\emph{Phys. Rev. A} \textbf{68}, 032303 (2003)},
\arxiv{quant-ph/0301132}.

\bibitem[Fan03]{Fan2003}
H. Fan,
Quantum cloning of mixed states in symmetric subspace,
\doi{10.1103/PhysRevA.68.054301}{\emph{Phys. Rev. A} \textbf{68}, 054301 (2003)},
\arxiv{quant-ph/0308058}.

\bibitem[FLS07]{FanLiuShi2007}
H. Fan, B. Y. Liu, and K. J. Shi,
Quantum cloning of identical mixed qubits,
\doi{10.26421/QIC7.5-6-8}{\emph{Quantum Inf. Comput.} \textbf{7}, 551--558 (2007)},
\arxiv{quant-ph/0601017}.

\bibitem[Kaz10]{Kazakov2007}
A. Ya. Kazakov,
``Partial'' quantum cloning and quantum cloning of mixed states,
\doi{10.1142/S0219749910006186}{\emph{Int. J. Quantum Inf.}
\textbf{8}, 435--442 (2010)},
\arxiv{0711.2860}.

\bibitem[CC07]{ChenChen2007MixedBroadcast}
L. Chen and Y.-X. Chen,
Mixed qubit cannot be universally broadcast,
\doi{10.1103/PhysRevA.75.062322}{\emph{Phys. Rev. A} \textbf{75}, 062322 (2007)},
\arxiv{quant-ph/0612205}.

\bibitem[LW17]{LemmWilde2017}
M. Lemm and M. M. Wilde,
Information-theoretic limitations on approximate quantum cloning and broadcasting,
\doi{10.1103/PhysRevA.96.012304}{\emph{Phys. Rev. A} \textbf{96}, 012304 (2017)},
\arxiv{1608.07569}.

\bibitem[DMP05]{DArianoMacchiavelloPerinotti2005}
G. M. D'Ariano, C. Macchiavello, and P. Perinotti,
Superbroadcasting of mixed states,
\doi{10.1103/PhysRevLett.95.060503}{\emph{Phys. Rev. Lett.} \textbf{95}, 060503 (2005)},
\arxiv{quant-ph/0506251}.

\bibitem[BDMP05]{BuscemiDArianoMacchiavelloPerinotti2005Optimal}
F. Buscemi, G. M. D'Ariano, C. Macchiavello, and P. Perinotti,
Optimal superbroadcasting of mixed qubit states,
\arxiv{quant-ph/0510155} (2005).

\bibitem[BDMP06]{BuscemiDArianoMacchiavelloPerinotti2006}
F. Buscemi, G. M. D'Ariano, C. Macchiavello, and P. Perinotti,
Universal and phase-covariant superbroadcasting for mixed qubit states,
\doi{10.1103/PhysRevA.74.042309}{\emph{Phys. Rev. A} \textbf{74}, 042309 (2006)},
\arxiv{quant-ph/0602125}.

\bibitem[DF07]{DangFan2007}
G.-F. Dang and H. Fan,
Optimal broadcasting of mixed states,
\doi{10.1103/PhysRevA.76.022323}{\emph{Phys. Rev. A} \textbf{76}, 022323 (2007)},
\arxiv{quant-ph/0609182}.

\bibitem[DPS06]{DArianoPerinottiSacchi2006Continuous}
G. M. D'Ariano, P. Perinotti, and M. F. Sacchi,
Superbroadcasting of continuous variable mixed states,
\doi{10.1088/1367-2630/8/6/099}{\emph{New J. Phys.} \textbf{8}, 99 (2006)},
\arxiv{quant-ph/0601114}.

\bibitem[PFBC24]{ParzygnatFullwoodBuscemiChiribella2024}
A. J. Parzygnat, J. Fullwood, F. Buscemi, and G. Chiribella,
Virtual quantum broadcasting,
\doi{10.1103/PhysRevLett.132.110203}{\emph{Phys. Rev. Lett.} \textbf{132}, 110203 (2024)},
\arxiv{2310.13049}.

\bibitem[BGA11]{BowlesGutaAdesso2011Purification}
P. Bowles, M. Gu{\c t}{\u a}, and G. Adesso,
Asymptotically optimal purification and dilution of mixed qubit and Gaussian states,
\doi{10.1103/PhysRevA.84.022320}{\emph{Phys. Rev. A} \textbf{84}, 022320 (2011)},
\arxiv{1102.2341}.

\bibitem[BGA12]{BowlesGutaAdesso2012Reversal}
P. Bowles, M. Gu{\c t}{\u a}, and G. Adesso,
Asymptotically optimal quantum channel reversal for qudit ensembles and multimode Gaussian states,
\doi{10.1088/1367-2630/14/11/113041}{\emph{New J. Phys.} \textbf{14}, 113041 (2012)},
\arxiv{1208.3569}.

\bibitem[Cav82]{Caves1982LinearAmplifiers}
C. M. Caves,
Quantum limits on noise in linear amplifiers,
\doi{10.1103/PhysRevD.26.1817}{\emph{Phys. Rev. D} \textbf{26}, 1817--1839 (1982)}.

\bibitem[CIR00]{CerfIpeRottenberg2000}
N. J. Cerf, A. Ipe, and X. Rottenberg,
Cloning of continuous quantum variables,
\doi{10.1103/PhysRevLett.85.1754}{\emph{Phys. Rev. Lett.} \textbf{85}, 1754--1757 (2000)},
\arxiv{quant-ph/9909037}.

\bibitem[BCI+01]{BraunsteinCerfIblisdirVanLoockMassar2001}
S. L. Braunstein, N. J. Cerf, S. Iblisdir, P. van Loock, and S. Massar,
Optimal cloning of coherent states with a linear amplifier and beam splitters,
\doi{10.1103/PhysRevLett.86.4938}{\emph{Phys. Rev. Lett.} \textbf{86}, 4938--4941 (2001)},
\arxiv{quant-ph/0012046}.

\bibitem[GJ07]{GutaJencova2007LAN}
M. Gu{\c t}{\u a} and A. Jen{\v c}ov\'a,
Local asymptotic normality in quantum statistics,
\doi{10.1007/s00220-007-0340-1}{\emph{Commun. Math. Phys.} \textbf{276}, 341--379 (2007)},
\arxiv{quant-ph/0606213}.

\bibitem[GK06]{GutaKahn2006LANQubits}
M. Gu{\c t}{\u a} and J. Kahn,
Local asymptotic normality for qubit states,
\doi{10.1103/PhysRevA.73.052108}{\emph{Phys. Rev. A} \textbf{73}, 052108 (2006)},
\arxiv{quant-ph/0512075}.

\bibitem[KG09]{KahnGuta2009LANFiniteDimensional}
J. Kahn and M. Gu{\c t}{\u a},
Local asymptotic normality for finite dimensional quantum systems,
\doi{10.1007/s00220-009-0787-3}{\emph{Commun. Math. Phys.} \textbf{289}, 597--652 (2009)},
\arxiv{0804.3876}.

\bibitem[YFG13]{YamagataFujiwaraGill2013LAN}
K. Yamagata, A. Fujiwara, and R. D. Gill,
Quantum local asymptotic normality based on a new quantum likelihood ratio,
\doi{10.1214/13-AOS1147}{\emph{Ann. Statist.} \textbf{41}, 2197--2217 (2013)},
\arxiv{1210.3749}.

\bibitem[GM06]{GutaMatsumoto2006MixedGaussianCloning}
M. Gu{\c t}{\u a} and K. Matsumoto,
Optimal cloning of mixed Gaussian states,
\doi{10.1103/PhysRevA.74.032305}{\emph{Phys. Rev. A} \textbf{74}, 032305 (2006)},
\arxiv{quant-ph/0605161}.

\bibitem[ZC17]{ZhaoChiribella2017}
X. Zhao and G. Chiribella,
Quantum amplification and purification of noisy coherent states,
\doi{10.1103/PhysRevA.95.042303}{\emph{Phys. Rev. A} \textbf{95}, 042303 (2017)},
\arxiv{1701.04602}.

\bibitem[LN24]{LahiryNussbaum2024}
S. Lahiry and M. Nussbaum,
Minimax estimation of low-rank quantum states and their linear functionals,
\doi{10.3150/23-BEJ1610}{\emph{Bernoulli} \textbf{30}, 610--635 (2024)},
\arxiv{2111.03279}.

\bibitem[Kup02]{Kuperberg2002}
G. Kuperberg,
Random words, quantum statistics, central limits, random matrices,
\emph{Methods Appl. Anal.} \textbf{9}, 101--119 (2002),
\arxiv{math/9909104}.

\bibitem[Joh01]{Johansson2001}
K. Johansson,
Discrete orthogonal polynomial ensembles and the Plancherel measure,
\doi{10.2307/2661375}{\emph{Ann. of Math.} \textbf{153}, 259--296 (2001)},
\arxiv{math/9906120}.

\bibitem[KW01]{KeylWerner2001Spectrum}
M. Keyl and R. F. Werner,
Estimating the spectrum of a density operator,
\doi{10.1103/PhysRevA.64.052311}{\emph{Phys. Rev. A} \textbf{64}, 052311 (2001)},
\arxiv{quant-ph/0102027}.

\bibitem[PRV67]{ParthasarathyRangaRaoVaradarajan1967}
K. R. Parthasarathy, R. Ranga Rao, and V. S. Varadarajan,
Representations of complex semi-simple Lie groups and Lie algebras,
\doi{10.2307/1970351}{\emph{Annals of Mathematics} \textbf{85}, 383--429 (1967)}.

\bibitem[FH91]{FultonHarris1991}
W. Fulton and J. Harris,
\emph{Representation Theory: A First Course},
Graduate Texts in Mathematics 129, Springer (1991).

\bibitem[TWZ25]{TangWrightZhandry2025}
E. Tang, J. Wright, and M. Zhandry,
Conjugate queries can help,
\arxiv{2510.07622} (2025).

\bibitem[FGSS26]{FanizzaGrinkoScharnhorstSpilecki2026}
M. Fanizza, D. Grinko, T. Scharnhorst, and J. Spilecki,
Optimal cloning of mixed states,
\arxiv{2608.27298} (2026).

\bibitem[JSO26]{JeonSohnOh2026}
S. Jeon, V. Sohn, and C. Oh,
Optimal copy complexity of quantum state cloning,
\arxiv{2608.24484} (2026).

\bibitem[LTHC26]{LiTheilHarrowChuang2026}
Z. Li, E. Theil, A. W. Harrow, and I. Chuang,
An exponential sample-complexity advantage for coherent quantum inference,
\arxiv{2605.21457} (2026).

\bibitem[LTHC26b]{LiTheilHarrowChuangQPA2026}
Z. Li, E. Theil, A. W. Harrow, and I. Chuang,
Quantum purity amplification for arbitrary eigenstates and multiple outputs,
\arxiv{2605.21570} (2026).

\end{thebibliography}

\end{document}
